% Le Cameroun à la veille de la colonisation — Histoire 1re Tech — Cours signé OnBuch+
% Programme officiel MINESEC. Compiler avec : tectonic N04-cameroun-veille-colonisation.tex
\newcommand{\DOCMATIERE}{Histoire}
\newcommand{\DOCNIVEAU}{1re Tech}
\newcommand{\DOCMODULE}{Histoire – classe de Première technique}
\newcommand{\DOCLECON}{N04}
\newcommand{\DOCTITRE}{Le Cameroun à la veille de la colonisation}
% OnBuch+ — préambule « cours signés » ENRICHI (compilé avec tectonic / XeLaTeX)
% Système visuel « Pop » d'OnBuch+ : fond crème (jamais blanc), bordures encre
% épaisses, ombres « dures » décalées, titres Archivo Black en capitales.
% Version MATHÉMATIQUES : graphes (pgfplots/tikz), boîtes pédagogiques
% (définition, propriété, méthode, attention, le savais-tu, exercice résolu,
% exercice, TP, bilan), page de garde et signature OnBuch+ ancrée en bas.
%
% Macros attendues dans main.tex AVANT \input{preamble} :
%   \DOCMATIERE  \DOCNIVEAU  \DOCMODULE  \DOCLECON (numéro)  \DOCTITRE
\documentclass[11pt]{article}

\usepackage{fontspec}
\usepackage[french]{babel}
\usepackage[a4paper,margin=2cm,top=3.4cm,bottom=2.8cm,headheight=1.4cm,footskip=1.3cm]{geometry}
\usepackage{amsmath,amssymb,amsfonts}
\usepackage{mathtools} % \xmapsto, \coloneqq... souvent utilisés par les modèles
\usepackage{cancel} % simplifications barrées (bilans de forces)
% \abs{z} (module, valeur absolue) et \norm{u} (norme), utilisés par les modèles
\providecommand{\abs}{}\let\abs\relax\DeclarePairedDelimiter{\abs}{\lvert}{\rvert}
\providecommand{\norm}{}\let\norm\relax\DeclarePairedDelimiter{\norm}{\lVert}{\rVert}
\usepackage[table]{xcolor} % [table] : fournit aussi \rowcolors (alternance de lignes)
\usepackage{pagecolor}
\usepackage{tikz}
\usetikzlibrary{arrows.meta,positioning,calc,shapes.geometric,shapes.misc,shapes.arrows,decorations.pathmorphing,decorations.pathreplacing,decorations.markings,patterns,shadows,mindmap,trees,fit,backgrounds,matrix,intersections,angles,quotes}
\usepackage{pgfplots}
\usepackage[version=4]{mhchem} % équations nucléaires \ce{^{235}_{92}U}
\usepackage[european,siunitx]{circuitikz} % schémas électriques
\pgfplotsset{compat=1.18}
\usepgfplotslibrary{fillbetween,groupplots} % aires entre courbes (calcul intégral)
\usepackage{enumitem}
\usepackage{fancyhdr}
% [most] charge toutes les bibliothèques tcolorbox (listings, minted, raster,
% documentation...) et gonfle inutilement la mémoire TeX sur un document long
% (~300 boîtes) : on ne charge que ce qui est réellement utilisé.
\usepackage[skins,breakable]{tcolorbox}
\usepackage{graphicx}
\usepackage{colortbl}
\usepackage{booktabs}
\usepackage{tabularx}
\usepackage{multirow}
\usepackage{diagbox} % case d'en-tête diagonale des tableaux à double entrée
% Colonnes extensibles centrée / gauche / droite (C, L, R), souvent utilisées par les modèles
\newcolumntype{C}{>{\centering\arraybackslash}X}
\newcolumntype{L}{>{\raggedright\arraybackslash}X}
\newcolumntype{R}{>{\raggedleft\arraybackslash}X}
\usepackage{array}
\usepackage{multicol}
\usepackage{siunitx}
% parse-numbers=false : accepte aussi \SI{2,5 \times 10^{-3}}{...}
\sisetup{locale=FR,output-decimal-marker={,},group-minimum-digits=4,parse-numbers=false}
\DeclareSIUnit\ohm{\ensuremath{\symup{\Omega}}} % U+2126 absent de la police
\DeclareSIUnit\division{div} % réglages d'oscilloscope (ms/div, V/div)
\DeclareSIUnit\tr{tr} % tours (vitesse de rotation, ex. \SI{1500}{\tr\per\minute})
\DeclareSIUnit\molar{M} % concentration molaire (ex. \SI{3}{\molar}, \SI{1}{\micro\molar})
\DeclareSIUnit\an{an} % année (le modèle confond parfois avec \year, un compteur TeX) — ex. \SI{5}{\kilo\meter\per\an}
\DeclareSIUnit\FCFA{FCFA} % franc CFA (ex. \SI{500}{\FCFA\per\kilo\gram})
\DeclareSIUnit\degreeBrix{\textdegree Brix} % teneur en sucre d'un sirop/jus (réfractométrie)
\DeclareSIUnit\month{mois} % le modèle utilise parfois \month, un compteur TeX, comme unité de durée
\DeclareSIUnit\microliter{\micro\liter} % le modèle écrit parfois \microliter en un seul token
\DeclareSIUnit\million{millions} % dénombrement (ex. \SI{15}{\million\per\mL} spermatozoïdes)
\usepackage{qrcode}
% \degree : commande fréquemment utilisée par les modèles pour un angle ou une
% température en dehors de \SI{}{} (non fournie par les paquets chargés).
\providecommand{\degree}{^{\circ}}
% \qed (fin de démonstration), utilisé par les modèles sans amsthm
\providecommand{\qed}{\hfill$\square$}
% \barre : double barre des valeurs interdites dans un tableau de variations (tkz-tab)
\providecommand{\barre}{\ensuremath{\|}}
% environnement proof (amsthm non chargé) utilisé par les modèles
\ifdefined\proof\else\newenvironment{proof}[1][Démonstration]{\par\noindent\textit{#1.}\ }{\qed\par}\fi
\usepackage{lastpage}
\usepackage{hyperref}
\hypersetup{hidelinks}

% ============================================================
% Palette « Pop » OnBuch+ — mêmes hex que la classe P (plus_ui.dart)
% ============================================================
\definecolor{poporange}{HTML}{F59321}
\definecolor{popdark}{HTML}{E07A0C}
\definecolor{popcream}{HTML}{FFF6E8}
\definecolor{popink}{HTML}{1C1714}
\definecolor{poppurple}{HTML}{7A5AE0}
\definecolor{popgreen}{HTML}{2E9E6A}
\definecolor{poppink}{HTML}{E0457B}
\definecolor{popblue}{HTML}{2F7DE1}
\definecolor{popgold}{HTML}{C98A12}
\definecolor{popmuted}{HTML}{8A7F76}
\definecolor{popline}{HTML}{EADFD2}
\definecolor{popgray}{HTML}{8A7F76}
\colorlet{popgrey}{popgray}
% Alias tolérants (le modèle nomme parfois les couleurs différemment)
\colorlet{popwhite}{white}
\colorlet{popblack}{popink}
\colorlet{popred}{poppink}
\colorlet{popyellow}{popgold}
% Teintes claires (fonds de tableaux/encadrés) — le modèle nomme parfois
% les couleurs avec un suffixe L (ex. popblueL) sans les avoir définies.
\colorlet{poporangeL}{poporange!20}
\colorlet{popdarkL}{popdark!20}
\colorlet{poppurpleL}{poppurple!20}
\colorlet{popgreenL}{popgreen!20}
\colorlet{poppinkL}{poppink!20}
\colorlet{popblueL}{popblue!20}
\colorlet{popgoldL}{popgold!20}
\colorlet{popredL}{popred!20}
\colorlet{popgrayL}{popgray!20}
% Teintes claires pour fonds de tableaux / zones
\colorlet{popblueL}{popblue!10!white}
\colorlet{poporangeL}{poporange!14!white}
\colorlet{popgreenL}{popgreen!12!white}
\colorlet{poppurpleL}{poppurple!12!white}
\colorlet{poppinkL}{poppink!12!white}
\colorlet{popgoldL}{popgold!14!white}

\pagecolor{popcream}

% ============================================================
% Typographie
% ============================================================
\defaultfontfeatures{Ligatures=TeX}
\setmainfont{PlusJakartaSans-Medium.ttf}[Path=../../../fonts/,BoldFont=PlusJakartaSans-Bold.ttf]
% Maths en sans-serif (Fira Math) avec chiffres et lettres droites de Plus
% Jakarta, pour que les formules se fondent dans le texte.
\usepackage[math-style=ISO,bold-style=ISO]{unicode-math}
\setmathfont{FiraMath-Regular.otf}[Path=../../../fonts/]
\setmathfont{PlusJakartaSans-Medium.ttf}[Path=../../../fonts/,range={up/num,up/{Latin,latin},\mathup}]
\usepackage{icomma} % virgule décimale sans espace en mode math
% Caractères mathématiques Unicode absents des polices de texte (Plus Jakarta,
% Archivo Black) : rendus par leur équivalent mathématique, en texte comme en
% formule — sinon ils s'affichent en carré vide (ex. « Arithmétique dans ℤ »).
\usepackage{newunicodechar}
\newunicodechar{ℝ}{\ensuremath{\mathbb{R}}}
\newunicodechar{ℤ}{\ensuremath{\mathbb{Z}}}
\newunicodechar{ℂ}{\ensuremath{\mathbb{C}}}
\newunicodechar{ℕ}{\ensuremath{\mathbb{N}}}
\newunicodechar{ℚ}{\ensuremath{\mathbb{Q}}}
\newunicodechar{𝔻}{\ensuremath{\mathbb{D}}}
\newunicodechar{α}{\ensuremath{\alpha}}
\newunicodechar{β}{\ensuremath{\beta}}
\newunicodechar{γ}{\ensuremath{\gamma}}
\newunicodechar{δ}{\ensuremath{\delta}}
\newunicodechar{ε}{\ensuremath{\varepsilon}}
\newunicodechar{θ}{\ensuremath{\theta}}
\newunicodechar{λ}{\ensuremath{\lambda}}
\newunicodechar{μ}{\ensuremath{\mu}}
\newunicodechar{σ}{\ensuremath{\sigma}}
\newunicodechar{τ}{\ensuremath{\tau}}
\newunicodechar{φ}{\ensuremath{\varphi}}
\newunicodechar{ψ}{\ensuremath{\psi}}
\newunicodechar{ω}{\ensuremath{\omega}}
\newunicodechar{Σ}{\ensuremath{\Sigma}}
% majuscules produites par \MakeUppercase dans les titres (α-aminés -> Α)
\newunicodechar{Α}{\ensuremath{\alpha}}
\newunicodechar{Β}{\ensuremath{\beta}}
\newunicodechar{Γ}{\ensuremath{\Gamma}}
\newunicodechar{Δ}{\ensuremath{\Delta}}
\newunicodechar{Ε}{\ensuremath{\varepsilon}}
\newunicodechar{Θ}{\ensuremath{\Theta}}
\newunicodechar{Λ}{\ensuremath{\Lambda}}
\newunicodechar{Μ}{\ensuremath{\mu}}
\newunicodechar{Τ}{\ensuremath{\tau}}
\newunicodechar{Φ}{\ensuremath{\Phi}}
\newunicodechar{Ψ}{\ensuremath{\Psi}}
\newunicodechar{Ω}{\ensuremath{\Omega}}
\newunicodechar{↦}{\ensuremath{\mapsto}}
\newunicodechar{∈}{\ensuremath{\in}}
\newunicodechar{∉}{\ensuremath{\notin}}
\newunicodechar{∧}{\ensuremath{\wedge}}
\newunicodechar{⇔}{\ensuremath{\Leftrightarrow}}
\newunicodechar{⟺}{\ensuremath{\Longleftrightarrow}}
\newunicodechar{⇒}{\ensuremath{\Rightarrow}}
\newunicodechar{⟹}{\ensuremath{\Longrightarrow}}
\newunicodechar{∘}{\ensuremath{\circ}}
\newunicodechar{∪}{\ensuremath{\cup}}
\newunicodechar{∩}{\ensuremath{\cap}}
\newunicodechar{≡}{\ensuremath{\equiv}}
\newunicodechar{⊂}{\ensuremath{\subset}}
\newunicodechar{⊥}{\ensuremath{\perp}}
\newunicodechar{∃}{\ensuremath{\exists}}
\newunicodechar{∀}{\ensuremath{\forall}}
\newunicodechar{∛}{\ensuremath{\sqrt[3]{\,}}}
\newunicodechar{∜}{\ensuremath{\sqrt[4]{\,}}}
\newunicodechar{⃗}{\ensuremath{{}^{\to}}}
\newunicodechar{ⁿ}{\ifmmode^{n}\else\textsuperscript{\ensuremath{n}}\fi}
\newunicodechar{ˣ}{\ifmmode^{x}\else\textsuperscript{\ensuremath{x}}\fi}
\newunicodechar{ᵇ}{\ifmmode^{b}\else\textsuperscript{\ensuremath{b}}\fi}
\newunicodechar{ʳ}{\ifmmode^{r}\else\textsuperscript{\ensuremath{r}}\fi}
\newunicodechar{ᵘ}{\ifmmode^{u}\else\textsuperscript{\ensuremath{u}}\fi}
\newunicodechar{ʸ}{\ifmmode^{y}\else\textsuperscript{\ensuremath{y}}\fi}
\newunicodechar{ᵃ}{\ifmmode^{a}\else\textsuperscript{\ensuremath{a}}\fi}
\newunicodechar{ᵉ}{\ifmmode^{e}\else\textsuperscript{\ensuremath{e}}\fi}
\newunicodechar{ᵏ}{\ifmmode^{k}\else\textsuperscript{\ensuremath{k}}\fi}
\newunicodechar{ᵖ}{\ifmmode^{p}\else\textsuperscript{\ensuremath{p}}\fi}
\newunicodechar{ᴺ}{\ifmmode^{N}\else\textsuperscript{\ensuremath{N}}\fi}
\newunicodechar{ᶜ}{\ifmmode^{c}\else\textsuperscript{\ensuremath{c}}\fi}
\newunicodechar{ʰ}{\ifmmode^{h}\else\textsuperscript{\ensuremath{h}}\fi}
\newunicodechar{ᵠ}{\ifmmode^{q}\else\textsuperscript{\ensuremath{q}}\fi}
\newunicodechar{ₙ}{\ifmmode_{n}\else\textsubscript{\ensuremath{n}}\fi}
\newunicodechar{ₐ}{\ifmmode_{a}\else\textsubscript{\ensuremath{a}}\fi}
\newunicodechar{ᵢ}{\ifmmode_{i}\else\textsubscript{\ensuremath{i}}\fi}
\newunicodechar{ᵣ}{\ifmmode_{r}\else\textsubscript{\ensuremath{r}}\fi}
\newunicodechar{ₖ}{\ifmmode_{k}\else\textsubscript{\ensuremath{k}}\fi}
\newunicodechar{ᵦ}{\ifmmode_{\beta}\else\textsubscript{\ensuremath{\beta}}\fi}
\newunicodechar{ₓ}{\ifmmode_{x}\else\textsubscript{\ensuremath{x}}\fi}
\newfontfamily\popdisplay{ArchivoBlack-Regular.ttf}[Path=../../../fonts/]
\newfontfamily\poplogo{SpaceGrotesk-Bold.ttf}[Path=../../../fonts/]
\newfontfamily\popsemi{PlusJakartaSans-SemiBold.ttf}[Path=../../../fonts/]
\newfontfamily\popxbold{PlusJakartaSans-ExtraBold.ttf}[Path=../../../fonts/]
\newfontfamily\popxboldls{PlusJakartaSans-ExtraBold.ttf}[Path=../../../fonts/,LetterSpace=3.0]

\setlist[enumerate]{leftmargin=1.7em,itemsep=5pt,topsep=4pt}
\setlist[itemize]{leftmargin=1.7em,itemsep=4pt,topsep=4pt,label={\color{poporange}\textbullet}}
\setlength{\parindent}{0pt}
\setlength{\parskip}{5pt}
\renewcommand{\arraystretch}{1.25}

% pgfplots : style Pop par défaut
\pgfplotsset{
  every axis/.append style={axis line style={popink,line width=1pt},tick style={popink},
    grid style={popline,line width=0.5pt},label style={font=\small},tick label style={font=\footnotesize}},
  popcurve/.style={poporange,line width=1.8pt,no markers,smooth},
}
% Même style disponible dans un \draw TikZ simple (hors pgfplots)
\tikzset{popcurve/.style={poporange,line width=1.8pt}}
\tikzset{popgrid/.style={popline,very thin}}
\colorlet{popcurve}{poporange} % le nom du style est parfois utilisé comme couleur

% ============================================================
% Pill / squiggle
% ============================================================
\newcommand{\poppill}[3]{%
  \tikz[baseline=-0.55ex]\node[draw=popink,line width=1.2pt,rounded corners=2.6pt,%
    fill=#2,inner xsep=6.5pt,inner ysep=3pt]{%
    \popxboldls\fontsize{7}{7}\selectfont\color{#3}\MakeUppercase{#1}};%
}
\newcommand{\popsquiggle}[1]{%
  \tikz[baseline=-0.4ex]{
    \draw[#1,line width=2.6pt,line cap=round]
      plot [smooth] coordinates {(0,0.09) (0.34,0.01) (0.68,0.21) (1.02,0.01) (1.36,0.10)};
  }%
}
% Mot-clé surligné (marqueur orange sous le texte)
\newcommand{\cle}[1]{{\popxbold\color{popdark}#1}}

% ============================================================
% En-tête / pied
% ============================================================
\pagestyle{fancy}
\fancyhf{}
\renewcommand{\headrulewidth}{0pt}
\renewcommand{\footrulewidth}{0pt}
\lhead{\raisebox{-0.15ex}{\poplogo\fontsize{13}{13}\selectfont\color{popink}OnBuch\color{poporange}+}}
\rhead{\poppill{\DOCMATIERE\ \textbullet\ \DOCNIVEAU\ \textbullet\ Leçon \DOCLECON}{white}{popink}}
\lfoot{\popxbold\fontsize{7.6}{7.6}\selectfont\color{popmuted}\MakeUppercase{Cours signé OnBuch+}\ \textbullet\ {\popsemi\DOCTITRE}}
\rfoot{\tikz[baseline=-0.6ex]\node[rounded corners=8.5pt,fill=popink,minimum height=17pt,minimum width=17pt,inner xsep=5pt,inner sep=0]{\popxbold\fontsize{8}{8}\selectfont\color{popcream}\thepage\,/\,\pageref*{LastPage}};}

% ============================================================
% Titres
% ============================================================
\newcommand{\chaptitre}[2]{%
  \begin{tcolorbox}[enhanced,colback=poporange,colframe=popink,boxrule=2.6pt,arc=11pt,
      drop shadow=popink,left=15pt,right=15pt,top=12pt,bottom=11pt,before skip=3pt,after skip=18pt]
    \poppill{#2}{popink}{popcream}\par\vspace{9pt}
    {\raggedright\hyphenpenalty=10000\exhyphenpenalty=10000\popdisplay\fontsize{22}{24}\selectfont\color{popcream}\MakeUppercase{#1}\par}
  \end{tcolorbox}
}
% Section numérotée : pastille orange avec le numéro + titre Archivo
\newcounter{popsec}
\newcommand{\coursec}[1]{%
  \stepcounter{popsec}\setcounter{popsub}{0}%
  \par\vspace{16pt}\noindent
  \begin{minipage}{\linewidth}
  \tikz[baseline=-0.7ex]\node[draw=popink,line width=1.6pt,fill=poporange,rounded corners=4pt,minimum size=22pt,inner sep=0]{\popdisplay\fontsize{12}{12}\selectfont\color{popcream}\thepopsec};%
  \hspace{8pt}{\popdisplay\fontsize{15}{17}\selectfont\color{popink}\MakeUppercase{#1}}\par
  \vspace{4pt}\noindent{\color{poporange}\rule{36pt}{3.6pt}}
  \end{minipage}\par\nopagebreak\vspace{8pt}
}
\newcounter{popsub}
\newcommand{\courssub}[1]{%
  \stepcounter{popsub}%
  \par\vspace{9pt}\noindent{\popxbold\fontsize{11.5}{13}\selectfont\color{popink}{\color{poporange}\thepopsec.\thepopsub}\hspace{6pt}#1}\par\nopagebreak\vspace{4pt}
}

% ============================================================
% Boîtes pédagogiques — même recette : fond blanc, bordure épaisse colorée,
% ombre dure encre, badge plein. Titre optionnel [#1].
% ============================================================
\tcbset{popbase/.style={breakable,enhanced,colback=white,boxrule=2.3pt,arc=9pt,
  shadow={3pt}{-3pt}{0pt}{popink},left=14pt,right=14pt,top=11pt,bottom=11pt,before skip=11pt,after skip=15pt,
  fontupper=\popsemi\fontsize{10.6}{14.5}\selectfont\color{popink},
  fontlower=\popsemi\fontsize{10.6}{14.5}\selectfont\color{popink}}}
% Coche dessinée (le glyphe ✓ n'existe pas dans les polices du document)
\newcommand{\popcheck}[1]{\tikz[baseline=-0.5ex]\draw[#1,line width=1.6pt,line cap=round,line join=round](-0.13,0.0)--(-0.04,-0.09)--(0.13,0.1);}
\providecommand{\coche}{\popcheck{popgreen}}
\providecommand{\resultat}[1]{\par\smallskip\noindent\fcolorbox{popgreen}{popgreenL}{\parbox{\dimexpr\linewidth-2\fboxsep-2\fboxrule\relax}{\textbf{Résultat.} #1}}\par\smallskip}
\newcommand{\popboxhead}[3]{\poppill{#1}{#2}{white}\ifx\relax#3\relax\else\hspace{6pt}{\popxbold\fontsize{10.6}{12}\selectfont\color{popink}#3}\fi\par\vspace{7pt}}

\newtcolorbox{aretenir}[1][]{popbase,colframe=popgreen,before upper={\popboxhead{À retenir}{popgreen}{#1}}}
\newtcolorbox{exemplebox}[1][]{popbase,colframe=poporange,before upper={\popboxhead{Exemple}{poporange}{#1}}}
\newtcolorbox{definition}[1][]{popbase,colframe=popblue,before upper={\popboxhead{Définition}{popblue}{#1}}}
\newtcolorbox{propriete}[1][]{popbase,colframe=poppurple,before upper={\popboxhead{Propriété}{poppurple}{#1}}}
\newtcolorbox{methode}[1][]{popbase,colframe=popgold,before upper={\popboxhead{Méthode}{popgold}{#1}}}
\newtcolorbox{attention}[1][]{popbase,colframe=poppink,before upper={\popboxhead{Attention}{poppink}{#1}}}
\newtcolorbox{experience}[1][]{popbase,colframe=popblue,colback=popblueL,before upper={\popboxhead{Expérience}{popblue}{#1}}}
% « Le savais-tu ? » : carte violette pleine (contexte, culture, Cameroun)
\newtcolorbox{savaistu}[1][]{popbase,colback=poppurple,colframe=popink,
  fontupper=\popsemi\fontsize{10.4}{14.2}\selectfont\color{white},
  before upper={\poppill{Le savais-tu\,?}{white}{poppurple}\ifx\relax#1\relax\else\hspace{6pt}{\popxbold\color{white}#1}\fi\par\vspace{7pt}}}
% Exercice résolu : énoncé en haut, \tcblower puis la solution sur fond crème
\newcounter{popexo}\newcounter{popexr}
\newtcolorbox{exoresolu}[1][]{popbase,colframe=poporange,colbacklower=popcream,
  segmentation style={popink,line width=1.2pt,dashed},
  before upper={\stepcounter{popexr}\popboxhead{Exercice résolu \thepopexr}{poporange}{#1}},
  before lower={\poppill{Solution}{popink}{popcream}\par\vspace{6pt}}}
% Exercice (non résolu en place) — la correction est dans la partie Corrigés
\newtcolorbox{exercice}[1][]{popbase,colframe=popink,
  before upper={\stepcounter{popexo}\popboxhead{Exercice \thepopexo}{popink}{#1}}}
% Corrigé
\newtcolorbox{corrige}[1][]{popbase,colframe=popgreen,colback=popgreenL,
  before upper={\popboxhead{Corrigé}{popgreen}{#1}}}
% Objectifs / prérequis / bilan
\newtcolorbox{objectifs}{popbase,colframe=popink,colback=poporangeL,
  before upper={\popboxhead{Objectifs de la leçon}{popink}{}}}
\newtcolorbox{prerequis}{popbase,colframe=popink,colback=popblueL,
  before upper={\popboxhead{Prérequis}{popblue}{}}}
\newtcolorbox{bilan}{popbase,colframe=popink,colback=white,boxrule=2.8pt,
  before upper={\popboxhead{Fiche bilan}{poporange}{L'essentiel à connaître pour l'examen}}}
% Figure encadrée avec légende
\newtcolorbox{popfigure}[1][]{enhanced,breakable=false,colback=white,colframe=popline,boxrule=1.4pt,arc=8pt,
  left=10pt,right=10pt,top=10pt,bottom=8pt,before skip=10pt,after skip=12pt,halign=center,
  lower separated=false,halign lower=center,
  fontlower=\popsemi\fontsize{9}{11.5}\selectfont\color{popmuted},
  #1}
\newcommand{\legende}[1]{\par\vspace{6pt}{\popsemi\fontsize{9}{11.5}\selectfont\color{popmuted}#1}}

% ============================================================
% Signature OnBuch+
% ============================================================
\newcommand{\poplogoinline}[1]{{\poplogo\fontsize{#1}{#1}\selectfont\color{popink}OnBuch\color{poporange}+}}
\newcommand{\signatureonbuch}{%
  \begin{tcolorbox}[enhanced,colback=popink,colframe=popink,boxrule=2.2pt,arc=10pt,
      drop shadow=poporange,left=14pt,right=14pt,top=10pt,bottom=10pt,before skip=0pt,after skip=0pt]
    \tikz[baseline=-0.7ex]\node[circle,fill=popgreen,draw=popcream,line width=1.3pt,minimum size=22pt,inner sep=0]{\popcheck{white}};%
    \hspace{9pt}\begin{minipage}[c]{0.62\linewidth}
      {\popxbold\fontsize{12}{13}\selectfont\color{popcream}Signé\ \textbullet\ {\poplogo OnBuch\color{poporange}+}}\par\vspace{2pt}
      {\raggedright\popsemi\fontsize{8}{10.5}\selectfont\color{popline}\MakeUppercase{Équipe pédagogique OnBuch+}\\\MakeUppercase{Conforme au programme officiel MINESEC}\par}
    \end{minipage}\hfill
    {\poplogo\fontsize{22}{22}\selectfont\color{popcream}OnBuch\color{poporange}+}
  \end{tcolorbox}%
}

% ============================================================
% Page de garde
% #1 titre · #2 matière · #3 niveau · #4 module · #5 n° leçon · #6 durée ·
% #7 description / objectifs (texte ou itemize)
% ============================================================
\newcommand{\pagegarde}[7]{%
  \thispagestyle{empty}
  \newgeometry{a4paper,margin=1.7cm,top=1.6cm,bottom=1.4cm}%
  \begin{tikzpicture}[remember picture,overlay]
    \fill[poporange!18] ([xshift=-3.2cm,yshift=-3.6cm]current page.north east) circle (4.6cm);
    \fill[poppurple!14] ([xshift=2.4cm,yshift=7.6cm]current page.south west) circle (3.8cm);
  \end{tikzpicture}%
  {\poplogo\fontsize{30}{30}\selectfont\color{popink}OnBuch\color{poporange}+}\hfill
  \raisebox{0.6ex}{\poppill{Cours signé \textbullet\ Premium}{popink}{popcream}}\par
  \vspace{1pt}\popsquiggle{poppurple}\par
  \vspace{8pt}
  \begin{tcolorbox}[enhanced,colback=poporange,colframe=popink,boxrule=2.8pt,arc=13pt,
      drop shadow=popink,left=18pt,right=18pt,top=12pt,bottom=13pt,after skip=10pt]
    \poppill{#2 \textbullet\ #3}{popink}{popcream}\hspace{5pt}\poppill{#4}{white}{popink}\par\vspace{8pt}
    {\popdisplay\fontsize{14}{14}\selectfont\color{popink}LEÇON #5}\par\vspace{4pt}
    {\raggedright\hyphenpenalty=10000\exhyphenpenalty=10000\popdisplay\fontsize{25}{27}\selectfont\color{popcream}\MakeUppercase{#1}\par}
    \vspace{8pt}
    {\popxbold\fontsize{9.5}{9.5}\selectfont\color{popink}\MakeUppercase{Durée conseillée : #6}}
  \end{tcolorbox}
  \begin{tcolorbox}[enhanced,colback=white,colframe=popink,boxrule=2.2pt,arc=10pt,
      drop shadow=popink,left=16pt,right=16pt,top=10pt,bottom=10pt,after skip=8pt]
    \poppill{Dans cette leçon}{poppurple}{white}\par\vspace{5pt}
    \popsemi\fontsize{9.3}{12.6}\selectfont\color{popink}#7
  \end{tcolorbox}
  \noindent\begin{minipage}[t]{0.487\linewidth}
    \begin{tcolorbox}[enhanced,colback=popgreen,colframe=popink,boxrule=2.2pt,arc=10pt,
        drop shadow=popink,left=11pt,right=11pt,top=10pt,bottom=10pt,height=3.1cm,valign=center]
      \tcbox[colback=white,colframe=popink,boxrule=1.3pt,arc=5pt,boxsep=0pt,nobeforeafter,
        left=3pt,right=3pt,top=3pt,bottom=3pt,valign=center]{\qrcode[height=1.9cm]{https://play.google.com/store/apps/details?id=com.luvvix.onbuch&hl=fr}}%
      \hspace{8pt}\begin{minipage}[c]{0.5\linewidth}
        {\popdisplay\fontsize{11}{12.5}\selectfont\color{white}\MakeUppercase{Télécharge OnBuch}}\par\vspace{4pt}
        {\raggedright\popsemi\fontsize{8.4}{11}\selectfont\color{white}Gratuit \textbullet\ Google Play\\Tuteur IA, annales, concours, résultats\par}
      \end{minipage}
    \end{tcolorbox}
  \end{minipage}\hfill
  \begin{minipage}[t]{0.487\linewidth}
    \begin{tcolorbox}[enhanced,colback=poppurple,colframe=popink,boxrule=2.2pt,arc=10pt,
        drop shadow=popink,left=11pt,right=11pt,top=10pt,bottom=10pt,height=3.1cm,valign=center]
      \tikz[baseline=-0.7ex]\node[circle,fill=white,draw=popink,line width=1.3pt,minimum size=1.6cm,inner sep=0]{\popdisplay\fontsize{24}{24}\selectfont\color{poppurple}+};%
      \hspace{8pt}\begin{minipage}[c]{0.6\linewidth}
        {\popdisplay\fontsize{11}{12.5}\selectfont\color{white}\MakeUppercase{Passe à OnBuch+}}\par\vspace{4pt}
        {\raggedright\popsemi\fontsize{8.4}{11}\selectfont\color{white}Tous les cours signés, Tuteur IA illimité, fiches PDF — onglet OnBuch+ de l'app\par}
      \end{minipage}
    \end{tcolorbox}
  \end{minipage}
  \vspace{8pt}
  \popcontenu
  \vfill
  \signatureonbuch
  \restoregeometry
  \newpage
}

% Rangée « Ce cours contient » (page de garde)
\newcommand{\poptile}[3]{% #1 couleur #2 gros texte #3 légende
  \begin{tcolorbox}[enhanced,colback=#1,colframe=popink,boxrule=2pt,arc=9pt,shadow={2.5pt}{-2.5pt}{0pt}{popink},
      left=6pt,right=6pt,top=9pt,bottom=9pt,halign=center,valign=center,height=1.75cm,width=\linewidth,nobeforeafter]
    {\popdisplay\fontsize{12.5}{13}\selectfont\color{white}\MakeUppercase{#2}}\par\vspace{3pt}
    {\centering\popsemi\fontsize{7.8}{9.5}\selectfont\color{white}#3\par}
  \end{tcolorbox}}
\newcommand{\popcontenu}{%
  \par\noindent{\popxboldls\fontsize{7.5}{7.5}\selectfont\color{popmuted}CE COURS SIGNÉ CONTIENT}\par\vspace{7pt}
  \noindent\begin{minipage}[t]{0.235\linewidth}\poptile{popblue}{Cours}{complet, illustré, pas à pas}\end{minipage}\hfill
  \begin{minipage}[t]{0.235\linewidth}\poptile{poporange}{Méthodes}{et exercices résolus}\end{minipage}\hfill
  \begin{minipage}[t]{0.235\linewidth}\poptile{poppink}{Exercices}{type examen + corrigés}\end{minipage}\hfill
  \begin{minipage}[t]{0.235\linewidth}\poptile{popgreen}{Fiche bilan}{l'essentiel à retenir}\end{minipage}\par}

% Sommaire
\newtcolorbox{sommaire}{enhanced,colback=white,colframe=popink,boxrule=2.2pt,arc=10pt,
  drop shadow=popink,left=18pt,right=18pt,top=13pt,bottom=13pt,before skip=6pt,after skip=20pt,
  before upper={\poppill{Sommaire}{poppurple}{white}\par\vspace{9pt}\popsemi\fontsize{10.6}{14.5}\selectfont\color{popink}}}

% Signature de fin de cours — poussée tout en bas de la dernière page
\newcommand{\findecours}{%
  \par\vfill\nopagebreak
  \begin{center}\popsquiggle{poporange}\end{center}\vspace{4pt}
  \signatureonbuch
}
\AtBeginDocument{\def\llbracket{\mathopen{[\mkern-3.5mu[}}\def\rrbracket{\mathclose{]\mkern-3.5mu]}}} % intervalles d'entiers (glyphe ⟦ absent de la police)
% Glyphes absents de Fira Math : redéfinis à partir de glyphes disponibles
\AtBeginDocument{%
  \def\setminus{\mathbin{\backslash}}\let\smallsetminus\setminus
  \def\vdots{\mathord{\vbox{\baselineskip3.2pt\lineskiplimit0pt\kern4pt\hbox{.}\hbox{.}\hbox{.}}}}%
  \def\ddots{\mathinner{\mkern1mu\raise6.5pt\hbox{.}\mkern2mu\raise3.6pt\hbox{.}\mkern2mu\raise0.7pt\hbox{.}\mkern1mu}}%
  \def\triangle{\mathord{\scalebox{0.9}{$\symup{\Delta}$}}}%
  \def\nabla{\mathord{\raisebox{\depth}{\scalebox{1}[-1]{$\symup{\Delta}$}}}}%
  \def\checkmark{\popcheck{popdark}}%
  \def\star{\mathbin{\ast}}\def\bigstar{\mathord{\ast}}%
  \def\diamond{\mathbin{\scalebox{0.6}{$\Diamond$}}}%
  \def\bigcup{\mathop{\mathchoice{\raisebox{-0.3em}{\scalebox{1.7}{$\cup$}}}{\scalebox{1.25}{$\cup$}}{\cup}{\cup}}\displaylimits}%
  \def\bigcap{\mathop{\mathchoice{\raisebox{-0.3em}{\scalebox{1.7}{$\cap$}}}{\scalebox{1.25}{$\cap$}}{\cap}{\cap}}\displaylimits}%
  \def\lesseqqgtr{\mathrel{\lessgtr}}\def\bigoplus{\mathop{\oplus}\displaylimits}%
}
\providecommand{\ssi}{si et seulement si} % abréviation fréquente
\AtBeginDocument{\providecommand{\wideparen}{\overparen}} % arc de cercle

%% FIGFIX-BEGIN v5 — correctifs globaux des figures (docs/onbuch-generation/tools/figfix)
\usepackage{adjustbox}\usepackage{etoolbox}\tcbuselibrary{raster}
% toute figure TikZ/pgfplots plus large que la ligne est réduite à la largeur disponible
\BeforeBeginEnvironment{tikzpicture}{\begin{adjustbox}{max width=\linewidth}}
\AfterEndEnvironment{tikzpicture}{\end{adjustbox}}
% légendes pgfplots lisibles par-dessus les courbes
\pgfplotsset{every axis/.append style={legend style={fill=white,fill opacity=0.92,text opacity=1,draw=black!15,inner sep=3pt}}}
% étiquettes d'arêtes (syntaxe edge["..."]) avec fond blanc
\tikzset{every edge quotes/.append style={fill=white,inner sep=1.2pt}}
%% FIGFIX-END
\begin{document}

\pagegarde{Le Cameroun à la veille de la colonisation}{Histoire}{1re Tech}{Histoire – classe de Première technique}{N04}{2 séances (fiche de progression harmonisée)}{Cette leçon étudie les premières entités politiques qui structuraient le territoire camerounais avant la colonisation : le lamidat du Nord, la chefferie des Grassfields et le sultanat bamoun. À travers l'analyse de cartes, de photographies et de données, elle montre que le Cameroun précolonial n'était pas un espace désorganisé mais un territoire doté d'institutions politiques élaborées.
\vspace{4pt}
\begin{multicols}{2}\small
\begin{itemize}
\item À la fin de cette leçon, je sais situer sur une carte les principales entités politiques du Cameroun précolonial (lamidats, chefferies des Grassfields, sultanat bamoun)
\item À la fin de cette leçon, je sais définir les notions de lamidat, de chefferie et de sultanat
\item À la fin de cette leçon, je sais décrire l'organisation politique, sociale et économique d'un lamidat peul du Nord-Cameroun
\item À la fin de cette leçon, je sais expliquer le fonctionnement d'une chefferie des Grassfields (pouvoir du fo, sociétés secrètes, hiérarchie)
\item À la fin de cette leçon, je sais présenter l'organisation du sultanat bamoun et le rôle du roi Njoya
\item À la fin de cette leçon, je sais analyser des documents (cartes, photographies, textes) pour établir des faits historiques
\end{itemize}\end{multicols}}

\begin{sommaire}
\begin{multicols}{2}
\begin{enumerate}
\item Introduction : le Cameroun précolonial, un espace politiquement organisé
\item Le lamidat : l'exemple des lamidats peuls du Nord-Cameroun
\item La chefferie des Grassfields
\item Le sultanat bamoun
\item Étude comparée : trois modèles d'organisation politique
\item TD : l'organisation d'une chefferie de la localité de l'établissement
\item Activité d'intégration
\item Méthodes et exercices résolus
\item Exercices
\item Corrigés des exercices
\item Fiche bilan
\end{enumerate}
\end{multicols}
\end{sommaire}

\begin{objectifs}
\begin{itemize}
\item À la fin de cette leçon, je sais situer sur une carte les principales entités politiques du Cameroun précolonial (lamidats, chefferies des Grassfields, sultanat bamoun)
\item À la fin de cette leçon, je sais définir les notions de lamidat, de chefferie et de sultanat
\item À la fin de cette leçon, je sais décrire l'organisation politique, sociale et économique d'un lamidat peul du Nord-Cameroun
\item À la fin de cette leçon, je sais expliquer le fonctionnement d'une chefferie des Grassfields (pouvoir du fo, sociétés secrètes, hiérarchie)
\item À la fin de cette leçon, je sais présenter l'organisation du sultanat bamoun et le rôle du roi Njoya
\item À la fin de cette leçon, je sais analyser des documents (cartes, photographies, textes) pour établir des faits historiques
\item À la fin de cette leçon, je sais construire un croquis ou un diagramme représentant l'organisation d'une entité politique précoloniale
\item À la fin de cette leçon, je sais rédiger et justifier une démarche, une réponse ou une conclusion à partir de documents
\item À la fin de cette leçon, je sais appliquer ces notions à l'étude d'une chefferie de ma localité (TD)
\end{itemize}
\end{objectifs}

\begin{prerequis}
\begin{itemize}
\item La localisation des grandes régions naturelles et administratives du Cameroun (géographie de 4e/3e)
\item La notion d'État et d'organisation politique (éducation à la citoyenneté)
\item Les grandes étapes de la pénétration européenne en Afrique au XIXe siècle (Histoire de 3e)
\item La diversité des peuples et des langues du Cameroun
\item La lecture et l'interprétation d'une carte historique
\end{itemize}
\end{prerequis}

\newpage

\coursec{Introduction : le Cameroun précolonial, un espace politiquement organisé}

\courssub{Situation d'accroche et problématique}

Chaque année, à Foumban, la fête traditionnelle du \cle{Nguon} rassemble des milliers de visiteurs venus de tout le Cameroun et de l'étranger. Devant le palais du sultan des Bamouns, bâti en 1917 sur les fondations de résidences royales bien plus anciennes, les spectateurs admirent les danses, les tambours sacrés et les insignes royaux. Une question revient alors sur toutes les lèvres : comment ce royaume a-t-il pu exister, s'organiser et prospérer bien avant l'arrivée des Allemands en 1884 ?

Le même étonnement saisit le voyageur qui traverse le Nord-Cameroun : à Garoua, à Maroua, à Ngaoundéré ou à Rey-Bouba, les \cle{lamibé} (pluriel de lamido, chefs peuls) trônent encore dans leurs palais, entourés de leurs conseillers et de leurs gardes. Dans l'Ouest et le Nord-Ouest, les \cle{fo} (chefs) des Grassfields président les danses rituelles, rendent la justice et perçoivent les tributs, comme le faisaient leurs ancêtres il y a deux siècles.

Que sait-on réellement de ces entités politiques qui structuraient le territoire camerounais à la veille de la colonisation ? Le Cameroun d'avant 1884 était-il un espace désorganisé, sans chefs ni lois, comme l'ont longtemps prétendu les colonisateurs pour justifier leur conquête ?

\textbf{Problématique :} comment le territoire camerounais était-il politiquement organisé avant la colonisation, et quelles formes prenaient ces organisations ?

\textbf{Annonce du plan :} après avoir défini le cadre chronologique et réfuté l'idée reçue d'un Cameroun « sans organisation », nous présenterons la diversité des formes politiques précoloniales, puis les facteurs qui expliquent leur constitution, avant de localiser sur une carte les trois entités que nous étudierons en détail : le lamidat, la chefferie des Grassfields et le sultanat bamoun.

\courssub{Le cadre chronologique : qu'entend-on par « à la veille de la colonisation » ?}

L'expression « à la veille de la colonisation » désigne la période qui précède immédiatement l'installation des puissances européennes sur le territoire camerounais. Pour le Cameroun, la date charnière est le \textbf{12 juillet 1884} : ce jour-là, les chefs douala \cle{Bell} et \cle{Akwa} signent avec les représentants des commerçants allemands (notamment la maison Woermann) le \cle{traité germano-douala}. Par ce traité, les chefs cèdent aux Allemands des droits commerciaux et acceptent le protectorat de l'Allemagne sur la côte. Cet acte marque le début officiel de la colonisation allemande du Cameroun.

Étudier le Cameroun « à la veille de la colonisation », c'est donc étudier le territoire \textbf{avant 1884}, essentiellement aux XVIII\textsuperscript{e} et XIX\textsuperscript{e} siècles, période durant laquelle se constituent et se consolident les grandes entités politiques que nous allons découvrir.

\begin{attention}[Une date à ne pas confondre]
Ne confonds pas le traité germano-douala du 12 juillet 1884 avec la conférence de Berlin (novembre 1884 -- février 1885), qui fixe les règles générales du partage de l'Afrique entre puissances européennes. Pour le Cameroun, c'est bien le traité du 12 juillet 1884 qui ouvre la colonisation. Retenir : \textbf{avant 1884 = Cameroun précolonial ; après 1884 = Cameroun colonial}.
\end{attention}

\begin{popfigure}
\begin{center}
\begin{tikzpicture}[line cap=round]
% Frise chronologique
\draw[very thick, popblue, -{Stealth[length=3mm]}] (0,0) -- (15.5,0);
\foreach \x/\annee in {1/{XVIII\textsuperscript{e} s.}, 4.2/{1804}, 7.4/{vers 1830}, 10.6/{1850-1880}, 13.8/{12 juillet 1884}}{
  \draw[thick, popdark] (\x,0.15) -- (\x,-0.15);
  \node[below, font=\small\bfseries, text=popdark] at (\x,-0.2) {\annee};
}
% Événements au-dessus
\node[above, align=center, font=\small, text=popink, text width=2.6cm] at (1,0.25) {Formation des chefferies des Grassfields (migrations Tikar)};
\node[above, align=center, font=\small, text=popink, text width=2.4cm] at (4.2,0.25) {Début du djihad peul dans la région};
\node[above, align=center, font=\small, text=popink, text width=2.4cm] at (7.4,0.25) {Fondation des lamidats (Garoua, Maroua, Ngaoundéré...)};
\node[above, align=center, font=\small, text=popink, text width=2.6cm] at (10.6,0.25) {Apogée du sultanat bamoun et des chefferies};
\fill[poporange] (13.8,0) circle (0.14);
\node[above, align=center, font=\small\bfseries, text=poporange, text width=2.6cm] at (13.8,0.3) {Traité germano-douala : début de la colonisation};
% Bandeau période
\draw[decorate, decoration={brace, amplitude=5pt}, popgreen, thick] (0.2,1.9) -- (13.4,1.9);
\node[above, font=\small\bfseries, text=popgreen] at (6.8,2.15) {Cameroun précolonial : espace politiquement organisé};
\end{tikzpicture}
\end{center}
\legende{Frise chronologique : de la formation des entités politiques (XVIII\textsuperscript{e}--XIX\textsuperscript{e} siècles) au traité germano-douala du 12 juillet 1884, qui marque la fin de la période précoloniale.}
\end{popfigure}

\courssub{Réfuter une idée reçue : le Cameroun précolonial n'était pas un espace « sans organisation »}

Pendant longtemps, les colonisateurs et certains récits européens ont présenté l'Afrique précoloniale comme un continent « sans histoire, sans États, sans institutions », peuplé de groupes dispersés et désorganisés. Cette vision servait à justifier la conquête : coloniser aurait consisté à « apporter l'ordre » à des populations qui en manquaient. C'est une \textbf{idée reçue} que les faits historiques réfutent totalement.

\begin{attention}[Erreur fréquente à éviter]
L'erreur la plus fréquente consiste à croire (ou à écrire dans une copie) que l'Afrique précoloniale, et le Cameroun en particulier, ne connaissaient ni États ni institutions avant l'arrivée des Européens. C'est faux : les documents le prouvent. Les palais des sultans bamouns à Foumban, les palais des lamibé du Nord, les cases royales et les sociétés régulatrices des Grassfields, les récits des voyageurs et les cartes anciennes attestent l'existence de pouvoirs structurés, d'armées, de systèmes de justice et de perception d'impôts. Dans un commentaire de document, il faut toujours s'appuyer sur ces preuves matérielles et écrites.
\end{attention}

En réalité, le territoire camerounais d'avant 1884 était couvert d'\cle{entités politiques} de tailles et de formes variées, dotées de chefs reconnus, de territoires délimités, de lois et d'institutions.

\begin{definition}[Entité politique]
Une \cle{entité politique} est un groupe humain organisé sous une autorité reconnue (chef, roi, sultan, lamido) qui exerce un pouvoir sur un territoire délimité et sur une population. Elle possède généralement des institutions (conseil, armée, justice, impôt) et des règles de succession au pouvoir.
\end{definition}

\begin{exemplebox}[Des preuves concrètes d'organisation politique]
\begin{itemize}
\item Au \textbf{Nord}, les lamidats peuls possèdent un chef (le lamido), une capitale fortifiée, une cavalerie, des impôts prélevés sur les populations et des tribunaux appliquant la loi islamique.
\item Dans les \textbf{Grassfields}, chaque chefferie s'organise autour du \emph{fo}, entouré de sociétés régulatrices (comme le \emph{Takumbeng} ou les sociétés de masques) qui contrôlent son pouvoir, punissent les fautes et protègent la communauté.
\item Chez les \textbf{Bamouns}, le sultan règne depuis un palais, commande une armée et entretient une administration de quartiers. Sous le règne du sultan \cle{Njoya} (à partir de 1886), le royaume connaît même la création d'une écriture originale, le \emph{shu-mom}, preuve de la vitalité d'un État déjà ancien et structuré.
\end{itemize}
\end{exemplebox}

\courssub{Une diversité de formes politiques}

L'organisation politique du Cameroun précolonial n'est pas uniforme : on distingue trois grands types d'entités, dont deux seront étudiés en détail dans cette leçon.

\textbf{1. Les États centralisés.} Ce sont des entités de grande taille, dirigées par un souverain puissant dont l'autorité s'exerce sur tout un territoire grâce à une hiérarchie de dignitaires et de chefs subalternes. Au Cameroun, on trouve :
\begin{itemize}
\item les \cle{lamidats} peuls du Nord (Garoua, Maroua, Ngaoundéré, Rey-Bouba, Banyo...), nés du djihad du XIX\textsuperscript{e} siècle et dirigés par des lamibé ;
\item le \cle{sultanat bamoun} de Foumban, royaume centralisé de l'Ouest, dirigé par un sultan.
\end{itemize}

\textbf{2. Les chefferies hiérarchisées.} De taille plus modeste, elles sont nombreuses dans les \cle{Grassfields} (Ouest et Nord-Ouest actuels) : Bafoussam, Bamenda, Bafut, Kom, Bali... Chacune est dirigée par un \emph{fo} dont le pouvoir, sacralisé, est encadré et contrôlé par des sociétés secrètes et des conseils de notables.

\textbf{3. Les sociétés segmentaires.} Pour le contraste, mentionnons brièvement les sociétés dites « segmentaires » ou « sans État », fréquentes dans la forêt du Sud (chez certains groupes bantu comme les Beti ou les Bassa) : elles ne connaissent pas de chef central ; l'autorité est exercée par les chefs de lignage et les conseils d'anciens, à l'échelle du village ou du clan. Cela ne signifie pas absence d'organisation : ces sociétés possèdent leurs propres règles, leurs alliances et leurs mécanismes de règlement des conflits.

\begin{aretenir}[Trois modèles d'organisation politique]
À la veille de la colonisation, le Cameroun connaît trois grands modèles : les \textbf{États centralisés} (lamidats du Nord, sultanat bamoun), les \textbf{chefferies hiérarchisées} (Grassfields) et les \textbf{sociétés segmentaires} (forêt du Sud). Cette diversité prouve que le territoire était pleinement organisé avant 1884.
\end{aretenir}

\courssub{Les facteurs de constitution des entités politiques}

Comment ces États et chefferies se sont-ils formés ? Quatre grands facteurs expliquent leur apparition et leur consolidation entre le XVIII\textsuperscript{e} et le XIX\textsuperscript{e} siècle.

\textbf{a) L'islamisation au Nord.} Au début du XIX\textsuperscript{e} siècle, le djihad (guerre sainte) lancé par le peul \cle{Ousman dan Fodio} au Nigeria voisin s'étend à l'Adamaoua et au Nord-Cameroun. Les chefs peuls (\emph{ardobé}, singulier \emph{ardo}) conquièrent les populations locales, fondent des lamidats et y installent l'islam comme religion d'État. L'islam fournit un droit écrit, une organisation fiscale et une légitimité religieuse au pouvoir du lamido.

\textbf{b) Les migrations Tikar.} Dans les Grassfields, les vagues de migrations des populations \cle{Tikar}, parties de la plaine de l'Adamaoua et de la région de Bankim entre les XVII\textsuperscript{e} et XIX\textsuperscript{e} siècles, donnent naissance à de nombreuses chefferies. Chaque groupe migrant, conduit par un chef, s'installe sur un territoire, absorbe ou repousse les populations antérieures et fonde une dynastie : c'est l'origine de chefferies comme Bafoussam, Bamenda ou Kom.

\textbf{c) Le commerce.} Les échanges à longue distance enrichissent les chefs et financent leurs États. Les caravanes transportent la \cle{kola} (noix stimulante très demandée au Sahel), les \cle{esclaves}, les \cle{tissus}, le sel, le fer et les cauris (coquillages servant de monnaie). Le contrôle des routes commerciales — vers la côte atlantique au sud, vers le Sahel au nord — est un enjeu de pouvoir : le sultanat bamoun, placé au carrefour des routes entre la forêt et les savanes, doit une partie de sa prospérité à cette position.

\textbf{d) La guerre.} Les conflits entre groupes jouent un double rôle. D'un côté, la guerre permet la conquête de territoires et la capture d'esclaves (lamidats peuls, guerres d'expansion bamoun) ; de l'autre, la nécessité de se défendre pousse les populations à se regrouper autour d'un chef guerrier capable de les protéger, ce qui renforce les chefferies des Grassfields, souvent installées sur des sites fortifiés ou entourées de fossés défensifs.

\begin{savaistu}[D'où vient le mot « Cameroun » ?]
Le nom « Cameroun » vient du portugais \emph{Rio dos Camarões}, c'est-à-dire « rivière aux crevettes ». À la fin du XV\textsuperscript{e} siècle (vers 1472), le navigateur portugais Fernando Póo remonte l'estuaire du Wouri et, frappé par l'abondance des crevettes dans ses eaux, lui donne ce nom. Les Allemands en feront plus tard « Kamerun », appliqué à tout le territoire de leur colonie à partir de 1884.
\end{savaistu}

\courssub{Localisation des trois entités étudiées}

Les trois entités que nous étudierons en détail dans cette leçon occupent des positions géographiques distinctes, qui correspondent à trois milieux naturels différents :

\begin{itemize}
\item les \textbf{lamidats} (Garoua, Maroua, Ngaoundéré, Rey-Bouba) dans les savanes du \textbf{Nord} ;
\item les \textbf{chefferies des Grassfields} (Bafoussam, Bamenda, Bafut, Kom) sur les hauts plateaux de l'\textbf{Ouest} et du \textbf{Nord-Ouest} ;
\item le \textbf{sultanat bamoun} (capitale : Foumban) à la lisière est des Grassfields, dans l'\textbf{Ouest}.
\end{itemize}

\begin{popfigure}
\begin{center}
\begin{tikzpicture}[scale=1.05, font=\small]
% Contour très simplifié du Cameroun
\draw[very thick, popdark, fill=popcream] 
  (2.2,0) -- (4.6,0.2) -- (5.4,1.2) -- (5.2,2.4) -- (4.4,3.2) -- (4.6,4.2)
  -- (3.8,5.2) -- (3.4,6.4) -- (2.6,7.4) -- (2.9,8.4) -- (2.2,9.2)
  -- (1.2,8.6) -- (1.6,7.4) -- (1.0,6.4) -- (1.4,5.2) -- (0.8,4.2)
  -- (1.2,3.0) -- (0.6,2.0) -- (1.2,1.0) -- cycle;
% Zones
\fill[popgreenL, opacity=0.75] (1.6,5.4) -- (3.6,5.4) -- (3.3,6.5) -- (2.5,7.5) -- (2.8,8.3) -- (2.1,9.0) -- (1.3,8.5) -- (1.7,7.3) -- (1.1,6.4) -- cycle;
\fill[poporangeL, opacity=0.8] (1.0,3.2) -- (3.4,3.3) -- (3.7,4.4) -- (3.0,5.3) -- (1.5,5.2) -- (0.9,4.3) -- cycle;
\fill[poppurpleL, opacity=0.8] (3.2,3.5) -- (4.3,3.4) -- (4.5,4.3) -- (3.8,4.9) -- (3.3,4.4) -- cycle;
% Étiquettes de zones
\node[font=\small\bfseries, text=popgreen!60!black, align=center] at (2.4,7.6) {LAMIDATS\\ (Nord)};
\node[font=\small\bfseries, text=poporange!70!black, align=center] at (1.6,4.6) {CHEFFERIES\\ Grassfields};
\node[font=\small\bfseries, text=poppurple!80!black, align=center] at (4.6,5.0) {SULTANAT\\ bamoun};
% Villes lamidats
\fill[popgreen!60!black] (2.2,8.6) circle (0.09); \node[right, font=\footnotesize] at (2.3,8.6) {Maroua};
\fill[popgreen!60!black] (2.5,6.9) circle (0.09); \node[right, font=\footnotesize] at (2.6,6.9) {Garoua};
\fill[popgreen!60!black] (3.2,5.7) circle (0.09); \node[right, font=\footnotesize] at (3.3,5.7) {Ngaoundéré};
\fill[popgreen!60!black] (1.7,7.9) circle (0.09); \node[left, font=\footnotesize] at (1.6,7.9) {Rey-Bouba};
% Villes Grassfields
\fill[poporange!70!black] (1.4,4.0) circle (0.09); \node[left, font=\footnotesize] at (1.3,4.0) {Bafoussam};
\fill[poporange!70!black] (1.2,4.7) circle (0.09); \node[left, font=\footnotesize] at (1.1,4.7) {Bamenda};
\fill[poporange!70!black] (1.9,4.9) circle (0.09); \node[above, font=\footnotesize] at (1.9,5.0) {Bafut};
\fill[poporange!70!black] (2.2,4.5) circle (0.09); \node[below right, font=\footnotesize] at (2.25,4.35) {Kom};
% Foumban
\fill[poppurple!80!black] (3.8,4.1) circle (0.11); \node[right, font=\footnotesize\bfseries] at (3.9,4.1) {Foumban};
% Côte
\node[font=\footnotesize\itshape, text=popblue] at (0.9,0.6) {Océan Atlantique};
% Nord
\draw[-{Stealth[length=2.5mm]}, very thick, popdark] (5.6,8.2) -- (5.6,9.2);
\node[above, font=\small\bfseries, text=popdark] at (5.6,9.2) {N};
% Échelle indicative
\draw[thick, popdark] (0.6,-0.4) -- (1.6,-0.4);
\draw[thick, popdark] (0.6,-0.5) -- (0.6,-0.3);
\draw[thick, popdark] (1.6,-0.5) -- (1.6,-0.3);
\node[below, font=\footnotesize] at (1.1,-0.5) {$\approx$ 150 km};
\end{tikzpicture}
\end{center}
\legende{Carte schématique du Cameroun précolonial (croquis simplifié, non exact). \textbf{1} : zone des lamidats peuls au Nord (Maroua, Garoua, Ngaoundéré, Rey-Bouba). \textbf{2} : chefferies des Grassfields à l'Ouest et au Nord-Ouest (Bafoussam, Bamenda, Bafut, Kom). \textbf{3} : sultanat bamoun autour de Foumban.}
\end{popfigure}

\begin{methode}[Analyser une carte historique]
Pour tirer l'information d'une carte historique comme celle-ci, procède toujours en cinq étapes :
\begin{enumerate}
\item \textbf{Lire le titre} : il indique le sujet, le lieu et la période (ici : le Cameroun avant 1884).
\item \textbf{Lire la légende} : elle donne la signification des couleurs et des symboles (zones politiques, villes).
\item \textbf{Repérer l'échelle et l'orientation} : la flèche du nord et l'échelle permettent d'évaluer distances et positions.
\item \textbf{Situer les lieux} : localise les entités et les villes les unes par rapport aux autres (Nord, Ouest, côte).
\item \textbf{Dégager l'information principale} : formule en une ou deux phrases ce que la carte démontre (ici : le territoire camerounais était couvert d'entités politiques organisées et réparties selon les milieux).
\end{enumerate}
\end{methode}

\begin{exoresolu}[Lecture de la carte du Cameroun précolonial]
\textbf{Énoncé.} À l'aide de la carte ci-dessus, réponds aux questions suivantes :
\begin{enumerate}
\item Dans quelle partie du Cameroun se trouvent les lamidats ? Cite deux de leurs capitales.
\item Quelle entité politique a pour capitale Foumban ?
\item Compare la localisation des chefferies des Grassfields et celle du sultanat bamoun.
\item Quelle conclusion générale cette carte permet-elle de tirer sur le Cameroun à la veille de la colonisation ?
\end{enumerate}
\tcblower
\textbf{Solution détaillée.}
\begin{enumerate}
\item Les lamidats se trouvent dans le \textbf{Nord} du Cameroun, dans la zone des savanes. Deux de leurs capitales sont \textbf{Garoua} et \textbf{Maroua} (on peut aussi citer Ngaoundéré et Rey-Bouba).
\item Foumban est la capitale du \textbf{sultanat bamoun}.
\item Les chefferies des Grassfields occupent les \textbf{hauts plateaux de l'Ouest et du Nord-Ouest} (Bafoussam, Bamenda, Bafut, Kom), tandis que le sultanat bamoun se situe \textbf{plus à l'est}, en bordure des Grassfields, autour de Foumban. Les deux entités sont donc voisines, ce qui explique leurs échanges commerciaux et leurs conflits.
\item Conclusion : la carte montre que, bien avant 1884, le territoire camerounais était \textbf{entièrement couvert d'entités politiques organisées} — États centralisés au Nord et à l'Ouest, chefferies hiérarchisées sur les plateaux — ce qui réfute l'idée d'un Cameroun précolonial « sans organisation ».
\end{enumerate}
\end{exoresolu}

\begin{aretenir}[L'essentiel de l'introduction]
\begin{itemize}
\item « À la veille de la colonisation » signifie \textbf{avant le 12 juillet 1884}, date du traité germano-douala.
\item Le Cameroun précolonial était un espace \textbf{politiquement organisé} : États centralisés (lamidats, sultanat bamoun), chefferies hiérarchisées (Grassfields) et sociétés segmentaires (Sud forestier).
\item La constitution de ces entités s'explique par quatre facteurs : \textbf{l'islamisation} au Nord, les \textbf{migrations Tikar}, le \textbf{commerce} (kola, esclaves, tissus) et la \textbf{guerre}.
\item Nous étudierons en détail trois modèles : le \textbf{lamidat}, la \textbf{chefferie des Grassfields} et le \textbf{sultanat bamoun}.
\end{itemize}
\end{aretenir}

\coursec{Le lamidat : l'exemple des lamidats peuls du Nord-Cameroun}

Nous venons de voir, dans l'introduction, que le Cameroun précolonial n'était pas un espace désorganisé : bien avant l'arrivée des Européens, des sociétés y avaient construit des États structurés, avec des chefs, des lois, des impôts et des armées. Nous allons maintenant étudier le premier grand modèle d'organisation politique : le \cle{lamidat}, qui a dominé toute la partie septentrionale du Cameroun au XIX\textsuperscript{e} siècle. Pour comprendre le lamidat, il faut d'abord remonter à son origine : un événement majeur de l'histoire de l'Afrique de l'Ouest, le djihad peul.

\courssub{Les origines : le djihad d'Ousman dan Fodio et la conquête du Nord-Cameroun}

Au début du XIX\textsuperscript{e} siècle, dans l'actuel Nigeria, vivent les \cle{Peuls} (ou Foulbés), un peuple d'éleveurs en grande partie islamisé. En 1804, un érudit religieux peul, \cle{Ousman dan Fodio}, déclare le \cle{djihad} contre les souverains haoussas qu'il accuse de mal pratiquer l'islam et d'opprimer les populations.

\begin{definition}[Le djihad]
Le \cle{djihad} est une guerre sainte menée au nom de l'islam. Celui d'Ousman dan Fodio, lancé en 1804, aboutit à la création du \cle{califat de Sokoto}, un vaste empire musulman dont dépendront les lamidats camerounais.
\end{definition}

La victoire des Peuls aboutit à la fondation du \textbf{califat de Sokoto}, un empire dirigé par un calife (successeur religieux et politique). Pour administrer ce vaste territoire, le califat est divisé en provinces appelées \textbf{émirats} ou \textbf{lamidats}, confiées à des chefs ayant participé à la conquête.

C'est dans ce cadre qu'intervient \cle{Modibo Adama}. Savant peul originaire de la région, il reçoit d'Ousman dan Fodio, vers 1809, un drapeau, symbole de sa mission : conquérir et islamiser la région de l'Adamaoua (actuel Nord-Cameroun). De 1809 à sa mort en 1847, Modibo Adama mène une série de campagnes militaires contre les populations locales non musulmanes. Il fonde la ville de \textbf{Yola} (dans l'actuel Nigeria), capitale de son émirat, et installe des chefs peuls dans les principales localités conquises. Ainsi naissent les lamidats du Nord-Cameroun.

\begin{aretenir}[Les dates clés]
\begin{itemize}
\item \textbf{1804} : début du djihad d'Ousman dan Fodio au Nigeria ; fondation du califat de Sokoto.
\item \textbf{vers 1804} : fondation du lamidat de Rey-Bouba.
\item \textbf{1809--1847} : conquête du Nord-Cameroun par Modibo Adama, au nom du califat de Sokoto.
\end{itemize}
\end{aretenir}

\begin{popfigure}
\begin{center}
\begin{tikzpicture}[scale=0.95, every node/.style={transform shape}]
% Frise chronologique simplifiée
\draw[thick, popdark] (0, -1.2) -- (12, -1.2);
\foreach \x/\date/\event in {0/1804/Djihad Ousman dan Fodio, 2.5/1809/Modibo Adama reçoit le drapeau, 5/1835/Fondation Ngaoundéré, 7.5/1847/Mort Modibo Adama, 10/1880/Apogée lamidats, 12/1901/Conquête allemande} {
  \draw[thick, popdark] (\x, -1.2) -- (\x, -1.5) node[below, font=\tiny, align=center, text width=1.8cm] {\date: \event};
}
\end{tikzpicture}
\end{center}
\legende{Frise chronologique simplifiée : du djihad de 1804 à la conquête allemande. Repères temporels pour situer la naissance et l'apogée des lamidats.}
\end{popfigure}

\courssub{Définition du lamidat et exemples}

\begin{definition}[Le lamidat et le lamido]
Un \cle{lamidat} est un État dirigé par un \cle{lamido} (pluriel : \emph{lamibé}), souverain musulman qui détient à la fois le pouvoir \textbf{politique}, \textbf{militaire} et \textbf{religieux}. Le lamidat est une émanation du califat de Sokoto : le lamido reconnaît l'autorité du calife de Sokoto, à qui il doit allégeance et auquel il envoie parfois des tributs.
\end{definition}

Autrement dit, le lamido n'est pas un souverain totalement indépendant : il gouverne son territoire au nom du calife de Sokoto, comme un gouverneur investi d'une très large autonomie. Cette subordination est symbolisée par le drapeau remis par le calife lors de l'investiture du lamido.

Les principaux lamidats du Cameroun précolonial sont :

\begin{itemize}
\item le lamidat de \textbf{Rey-Bouba}, fondé vers 1804, l'un des plus anciens et des plus puissants ;
\item le lamidat de \textbf{Garoua}, sur le fleuve Bénoué ;
\item le lamidat de \textbf{Maroua}, dans l'Extrême-Nord ;
\item le lamidat de \textbf{Ngaoundéré}, fondé vers 1835 par Ardo Njobdi, dans l'Adamaoua ;
\item les lamidats de \textbf{Banyo} et de \textbf{Tibati}, au sud de l'Adamaoua.
\end{itemize}

\begin{exemplebox}[Un exemple chiffré : le lamidat de Rey-Bouba]
Le lamidat de Rey-Bouba couvrait environ \textbf{35 000 km\textsuperscript{2}}, soit une superficie supérieure à celle de nombreuses régions actuelles. C'était l'une des plus vastes entités politiques du Cameroun précolonial. Son territoire s'étendait de part et d'autre de la frontière actuelle entre le Cameroun et le Nigeria.
\end{exemplebox}

\begin{popfigure}
\begin{center}
\begin{tikzpicture}[scale=0.95, every node/.style={transform shape}]
% Fond : contour simplifié du Nord-Cameroun
\draw[thick, popdark, fill=popcream] (0.5,0) -- (2.2,-0.4) -- (4.5,0.2) -- (5.6,1.6) -- (5.2,3.2) -- (5.8,4.6) -- (4.6,6.2) -- (3.2,7.2) -- (1.8,6.6) -- (1.2,5.2) -- (0.2,4.0) -- (0.8,2.2) -- cycle;
% Flèche du nord
\draw[-{Stealth}, thick, popdark] (6.6,6.2) -- (6.6,7.2) node[above]{\textbf{N}};
% Lamidats
\fill[poporange] (3.3,6.3) circle (0.09) node[right, popdark]{\textbf{Rey-Bouba}};
\fill[poporange] (4.4,5.2) circle (0.09) node[right, popdark]{\textbf{Garoua}};
\fill[poporange] (3.6,4.4) circle (0.09) node[right, popdark]{\textbf{Maroua}};
\fill[poporange] (2.6,2.6) circle (0.09) node[right, popdark]{\textbf{Ngaoundéré}};
\fill[poporange] (1.6,1.4) circle (0.09) node[right, popdark]{\textbf{Banyo}};
\fill[poporange] (3.4,0.9) circle (0.09) node[right, popdark]{\textbf{Tibati}};
% Fleuve Bénoué (simplifié)
\draw[popblue, thick] (0.6,4.6) .. controls (2.0,5.4) and (3.4,5.6) .. (5.4,4.9) node[pos=0.35, above, popblue]{\footnotesize Bénoué};
% Axes caravaniers
\draw[-{Stealth}, very thick, popgreen] (2.6,2.6) -- (4.4,5.2);
\draw[-{Stealth}, very thick, popgreen] (4.4,5.2) -- (3.3,6.3);
\draw[-{Stealth}, very thick, popgreen] (1.6,1.4) -- (2.6,2.6);
\draw[-{Stealth}, very thick, popgreen] (4.4,5.2) -- (5.7,5.9) node[right, popgreen]{\footnotesize vers Sokoto};
% Echelle indicative
\draw[thick, popdark] (0.4,-0.9) -- (1.9,-0.9) node[midway, below]{\footnotesize 0 \quad 150 km};
\end{tikzpicture}
\end{center}
\legende{Carte schématique du Nord-Cameroun précolonial : les principaux lamidats (points orange), le fleuve Bénoué (bleu) et les axes du commerce caravanier (flèches vertes). Croquis simplifié, sans précision cartographique.}
\end{popfigure}

\courssub{L'organisation politique du lamidat}

Comment un seul homme peut-il gouverner un territoire de plusieurs milliers de kilomètres carrés ? La réponse tient en un mot : la \textbf{délégation}. Le lamido s'appuie sur toute une hiérarchie de dignitaires et de chefs subalternes.

Au sommet se trouve le \textbf{lamido}, qui concentre tous les pouvoirs : il commande l'armée, rend la justice en dernier ressort, lève les impôts et dirige la prière. Il est assisté de \textbf{conseillers} :

\begin{itemize}
\item le \cle{galadima}, sorte de premier ministre ou de conseiller principal, qui gère l'administration quotidienne ;
\item l'\cle{alkali}, le juge, qui applique la loi islamique (la charia) pour trancher les litiges ;
\item d'autres dignitaires chargés de la guerre, des finances ou du protocole.
\end{itemize}

Sur le terrain, le territoire est découpé en quartiers (dans les villes) et en villages (à la campagne) :

\begin{itemize}
\item les \cle{lawanes} sont les chefs de quartiers et de localités ;
\item les \cle{djauros} sont les chefs de villages, souvent chargés de l'encadrement des troupeaux et de la collecte des impôts au niveau local.
\end{itemize}

\begin{popfigure}
\begin{center}
\begin{tikzpicture}[
  node distance=0.55cm and 1.1cm,
  box/.style={draw=popdark, thick, rounded corners=2pt, fill=popblueL, text width=3.6cm, align=center, minimum height=0.85cm, font=\small},
  topbox/.style={box, fill=poporangeL, text width=4.6cm, font=\small\bfseries},
  link/.style={-{Stealth}, thick, popdark}]
\node[topbox, align=center] (lamido){Le LAMIDO\\ \footnotesize pouvoir politique, militaire et religieux};
\node[box, below left=0.7cm and -0.4cm of lamido, align=center] (galadima){Le galadima\\ \footnotesize conseiller principal};
\node[box, below right=0.7cm and -0.4cm of lamido, align=center] (alkali){L'alkali\\ \footnotesize juge (loi islamique)};
\node[box, below=1.9cm of lamido, align=center] (lawanes){Les lawanes\\ \footnotesize chefs de quartiers et de localités};
\node[box, below=of lawanes, align=center] (djauros){Les djauros\\ \footnotesize chefs de villages};
\node[box, below=of djauros, fill=popgreenL, align=center] (pop){Les populations\\ \footnotesize Peuls, agriculteurs, éleveurs, esclaves};
\draw[link] (lamido) -- (galadima);
\draw[link] (lamido) -- (alkali);
\draw[link] (galadima) |- (lawanes);
\draw[link] (alkali) |- (lawanes);
\draw[link] (lawanes) -- (djauros);
\draw[link] (djauros) -- (pop);
\end{tikzpicture}
\end{center}
\legende{Organigramme de la hiérarchie politique d'un lamidat : du lamido (sommet) jusqu'aux populations (base), en passant par les conseillers, les chefs de quartiers (lawanes) et les chefs de villages (djauros).}
\end{popfigure}

\begin{methode}[Analyser l'organisation politique d'un État précolonial]
Pour décrire rigoureusement l'organisation politique d'un État comme un lamidat, procède en quatre étapes :
\begin{enumerate}
\item Identifier le \textbf{chef suprême} : qui détient le pouvoir et quels pouvoirs possède-t-il ?
\item Repérer les \textbf{collaborateurs directs} : conseillers, ministres, juges.
\item Décrire l'\textbf{échelonnement territorial} : qui commande dans les provinces, les villes, les villages ?
\item Situer l'État dans son \textbf{ensemble plus vaste} : de quelle entité supérieure dépend-il (ici, le califat de Sokoto) ?
\end{enumerate}
\end{methode}

\courssub{L'organisation sociale : une société fortement hiérarchisée}

La société des lamidats est une société de \textbf{castes et de statuts}, très inégalitaire. Au sommet se trouvent les \textbf{Peuls musulmans}, conquérants et détenteurs du pouvoir : nobles, dignitaires, érudits religieux (les \emph{modibé}, savants de l'islam). Viennent ensuite les populations libres islamisées, intégrées progressivement à l'État.

En bas de l'échelle se trouvent les populations conquises, que les Peuls appelaient \emph{Kirdi} (terme signifiant « païens », c'est-à-dire non musulmans) : Mafa, Moundang, Toupouri, Guiziga, et bien d'autres. Ces populations étaient :
\begin{itemize}
\item soit \textbf{soumises} : elles conservaient leurs terres mais devaient payer un tribut au lamido ;
\item soit \textbf{réduites en esclavage} : les razzias (expéditions de capture) fournissaient des esclaves utilisés dans l'agriculture, le service domestique, l'armée, ou vendus aux caravaniers.
\end{itemize}

\begin{attention}[Un vocabulaire à manier avec rigueur]
Le terme \emph{Kirdi} est un terme péjoratif employé par les conquérants peuls pour désigner les populations non islamisées du Nord-Cameroun. Dans une copie, il faut l'employer entre guillemets ou préciser qu'il s'agit du vocabulaire de l'époque, et non d'un nom de peuple.
\end{attention}

\courssub{L'organisation économique : élevage, agriculture et commerce caravanier}

L'économie du lamidat repose sur trois piliers :

\begin{enumerate}
\item \textbf{L'élevage} : les Peuls sont d'abord des éleveurs de \textbf{zébus} (bovins à bosse). Le troupeau est un signe de richesse et de prestige social ; il fournit lait, viande et cuirs.
\item \textbf{L'agriculture} : les populations sédentaires cultivent le mil, le sorgho, le maïs, le coton et l'arachide. Une partie de la récolte est prélevée par l'État.
\item \textbf{Le commerce caravanier} : les lamidats sont traversés par de grandes routes commerciales qui relient l'Adamaoua au bassin du Tchad et au califat de Sokoto. Les caravanes transportent le sel, le kola, les étoffes, les cuirs, l'ivoire... et malheureusement aussi des esclaves.
\end{enumerate}

Pour financer son fonctionnement (palais, armée, dignitaires), l'État prélève des ressources :

\begin{itemize}
\item la \cle{zakkat}, impôt religieux musulman prélevé sur les troupeaux et les récoltes (taux canonique : 2,5 \% de la valeur du cheptel ou de la récolte) ;
\item les \textbf{tributs} payés par les populations conquises (en nature : mil, bétail, ou en captifs) ;
\item les taxes sur les marchés et les caravanes.
\end{itemize}

\begin{exoresolu}[Calculer la part d'un prélèvement fiscal]
Dans un lamidat, un éleveur possède un troupeau de 50 zébus. La zakkat sur le bétail représente environ un animal pour 40 (soit 2,5 \%). Combien de zébus l'éleveur doit-il verser au lamido ?
\tcblower
\textbf{Solution détaillée.} On calcule 2,5 \% de 50 :
\[ 50 \times \frac{2.5}{100} = 50 \times 0.025 = 1.25 \]
Le résultat théorique est 1,25 zébu. Or on ne peut pas prélever une fraction d'animal : dans la pratique, l'éleveur verse \textbf{un zébu} au représentant du lamido (le djauro chargé de la collecte). Cet exercice simple montre que la zakkat, bien que modique pour chaque contribuable, représentait une ressource considérable à l'échelle d'un lamidat comptant des dizaines de milliers de têtes de bétail.
\end{exoresolu}

\courssub{Religion et pouvoir : l'islam, ciment de l'État}

Dans le lamidat, la religion n'est pas séparée de la politique : elle en est le \textbf{fondement}. Le lamido gouverne au nom de l'islam et du calife de Sokoto ; sa légitimité est d'abord religieuse. Cette fusion du religieux et du politique se traduit concrètement :

\begin{itemize}
\item les \textbf{mosquées} sont des lieux de prière mais aussi de rassemblement politique, où le lamido s'adresse à la communauté, notamment lors de la grande prière du vendredi ;
\item les \textbf{écoles coraniques} forment les lettrés de l'État : futurs juges (alkali), secrétaires, conseillers. L'écriture arabe devient l'écriture de l'administration ;
\item le \textbf{calendrier musulman} rythme la vie publique (Ramadan, fêtes de l'Aïd), renforçant l'unité de la communauté autour du souverain.
\end{itemize}

L'islam joue ainsi le rôle de \textbf{ciment de l'État} : il unit des populations d'origines diverses autour d'une même foi, d'une même loi et d'un même chef.

\courssub{Étude de document : le palais du lamido, architecture et symboles du pouvoir}

\textbf{Description du document.} Le document étudié est une photographie du palais du lamido de Ngaoundéré (on pourrait aussi observer celui de Rey-Bouba). On y voit une vaste enceinte de hauts murs en \textbf{banco} (terre crue mélangée à de la paille), de couleur ocre, percée d'une porte monumentale encadrée de piliers massifs. Les murs sont surmontés de créneaux et de motifs en relief. À l'intérieur de l'enceinte se succèdent des cours, des cases aux toits de chaume coniques et une grande salle d'audience.

\textbf{Analyse.} Cette architecture n'est pas seulement fonctionnelle : elle est un \textbf{langage du pouvoir}.

\begin{enumerate}
\item Les \textbf{hauts murs} expriment la puissance militaire : le palais est une forteresse capable de résister à un siège.
\item La \textbf{porte monumentale unique} symbolise le contrôle : nul n'entre chez le lamido sans autorisation ; elle matérialise la distance entre le souverain et le peuple.
\item La \textbf{salle d'audience} est le cœur politique du palais : c'est là que le lamido reçoit, trône, rend la justice et accorde les audiences, entouré de ses conseillers.
\item L'emplacement \textbf{central} du palais dans la ville, souvent à proximité de la grande mosquée, affirme l'union du pouvoir politique et du pouvoir religieux.
\end{enumerate}

\textbf{Conclusion de l'étude.} Le palais du lamido est à la fois une résidence, une forteresse, un tribunal et un centre administratif : il concentre physiquement tous les pouvoirs que détient le lamido. Son observation permet donc de « lire » l'organisation politique du lamidat dans la pierre... ou plutôt dans le banco.

\begin{savaistu}[Un palais toujours debout]
Le palais du lamido de \textbf{Rey-Bouba}, avec ses hauts murs en banco, est encore visible aujourd'hui dans le département du Mayo-Rey (région du Nord). Il témoigne, deux siècles après sa fondation, de la puissance de cet État précolonial. Le lamidat de Rey-Bouba existe d'ailleurs toujours comme chefferie traditionnelle de premier degré : le lamido actuel y exerce encore une autorité morale et coutumière reconnue par l'État camerounais.
\end{savaistu}

\begin{attention}[Ne pas confondre lamidat et sultanat]
Piège classique du Probatoire : le \textbf{lamidat} relève du \textbf{califat de Sokoto} (Nigeria) et naît de la conquête peule du début du XIX\textsuperscript{e} siècle ; son chef est un \emph{lamido}. Le \textbf{sultanat bamoun}, au contraire, est un \textbf{royaume autonome}, fondé bien plus tôt (vers la fin du XIV\textsuperscript{e} siècle selon la tradition), qui ne dépend d'aucun empire extérieur et ne s'islamise que tardivement, sous le règne du sultan Njoya. Retenir : lamidat = Peuls + Sokoto ; sultanat = Bamouns + indépendance.
\end{attention}

\begin{exercice}[Pour s'entraîner]
\begin{enumerate}
\item Définis les termes suivants : lamidat, djihad, zakkat, djauro.
\item Explique en cinq lignes le lien entre le djihad d'Ousman dan Fodio et la création des lamidats du Nord-Cameroun.
\item Rédige un paragraphe argumenté (10 à 15 lignes) répondant à la question : « En quoi l'islam est-il le ciment de l'État dans un lamidat ? »
\end{enumerate}
\end{exercice}

\begin{corrige}[Exercice]
\begin{enumerate}
\item \textbf{Lamidat} : État dirigé par un lamido, souverain musulman détenant les pouvoirs politique, militaire et religieux, et relevant du califat de Sokoto. \textbf{Djihad} : guerre sainte menée au nom de l'islam ; celui de 1804 a fondé le califat de Sokoto. \textbf{Zakkat} : impôt religieux musulman prélevé sur les troupeaux et les récoltes (taux de 2,5 \%). \textbf{Djauro} : chef de village dans un lamidat, chargé notamment de la collecte des impôts.
\item Le djihad d'Ousman dan Fodio (1804) crée le califat de Sokoto au Nigeria. Pour étendre ce califat, Modibo Adama reçoit vers 1809 la mission de conquérir l'Adamaoua. Ses campagnes (1809--1847) soumettent le Nord-Cameroun, où il installe des chefs peuls : ce sont les lamidats (Garoua, Maroua, Ngaoundéré, Banyo, Tibati, Rey-Bouba), provinces du califat.
\item Attendu : introduction (l'islam n'est pas séparé du pouvoir), développement (légitimité religieuse du lamido ; allégeance au calife de Sokoto ; justice rendue selon la loi islamique par l'alkali ; rôle des mosquées et des écoles coraniques qui forment les cadres de l'État ; unité des populations autour d'une même foi), conclusion (l'islam assure la cohésion et la légitimité de l'État).
\end{enumerate}
\end{corrige}

\begin{aretenir}[L'essentiel de la section]
\begin{itemize}
\item Les lamidats du Nord-Cameroun naissent du djihad d'Ousman dan Fodio (1804) et de la conquête de Modibo Adama (1809--1847) ; ils relèvent du califat de Sokoto.
\item Le lamido concentre les pouvoirs politique, militaire et religieux ; il gouverne avec des conseillers (galadima, alkali) et des chefs locaux (lawanes, djauros).
\item La société est hiérarchisée : Peuls musulmans au sommet, populations conquises soumises ou esclavagisées.
\item L'économie repose sur l'élevage, l'agriculture, le commerce caravanier et les prélèvements (zakkat, tributs).
\item L'islam est le ciment de l'État ; le palais du lamido en est le symbole architectural.
\end{itemize}
\end{aretenir}

\coursec{La chefferie des Grassfields}

Après avoir étudié le lamidat, modèle d'État centralisé et islamisé du Nord-Cameroun né de la conquête peule, nous descendons maintenant vers le sud, sur les hautes terres de l'Ouest. Nous y découvrons un modèle politique très différent : au lieu d'un vaste État dirigé par un souverain musulman, nous trouvons une mosaïque de petites entités autonomes, les \cle{chefferies}, chacune dirigée par un chef sacré appelé \cle{fo} (ou \emph{fon}). Comprendre la chefferie des Grassfields, c'est comprendre comment une société dense, agricole et non islamisée a su inventer son propre système de gouvernement, avec ses contre-pouvoirs et ses symboles.

\courssub{Un cadre géographique et humain particulier}

\begin{definition}[Les Grassfields]
Les \cle{Grassfields} (littéralement « champs d'herbe ») désignent les hautes terres de l'Ouest et du Nord-Ouest du Cameroun actuels : plateaux bamiléké, monts Bamboutos, plateau Bamenda, grasslands du Nord-Ouest. C'est une région de collines et de plateaux (altitude 1\,000--2\,500 m), au climat tempéré de type camerounais, couverte de savane herbeuse favorable à l'agriculture et à l'élevage.
\end{definition}

\begin{popfigure}
\begin{center}
\begin{tikzpicture}[scale=0.85]
% Contour très simplifié du Cameroun (polygone approximatif)
\draw[popdark, thick] (0,0) -- (8,0) -- (9,2) -- (9,6) -- (7,8) -- (5,9) -- (3,8) -- (1,6) -- (0,3) -- cycle;
% Région des Grassfields (Ouest/Nord-Ouest) : zone hachurée
\fill[popgreen!25, pattern=north east lines, pattern color=popgreen!60] (3.5,3.5) -- (6,3.5) -- (6.5,5) -- (5.5,6) -- (4,5.5) -- (3,4.5) -- cycle;
% Villes repères
\foreach \x/\y/\name in {4.5/4.2/Bafoussam, 5.2/4.8/Bamenda, 4.0/5.0/Bafut, 3.8/4.0/Bangangte, 5.5/4.0/Dschang} {
  \fill[popred] (\x,\y) circle (2.5pt);
  \node[popdark, font=\tiny, anchor=west] at (\x+0.15,\y) {\name};
}
% Flèche Nord
\draw[popdark, thick, ->] (7.5,7.5) -- (7.5,8.3);
\node[popdark, font=\small\bfseries] at (7.5,8.6) {N};
% Légende
\node[popdark, font=\small, align=left] at (1.5,1.5) {Cameroun (contour simplifié)};
\node[popdark, font=\small, align=left] at (1.5,1.0) {Zone hachurée = Grassfields};
\end{tikzpicture}
\end{center}
\legende{Croquis schématique du Cameroun : localisation de la région des Grassfields (hachurée verte) et principales chefferies. Non contractuel, sans échelle.}
\end{popfigure}

Cette région se caractérise par une \cle{forte densité de population}, parmi les plus élevées de toute l'Afrique centrale (plus de 100 hab./km² par endroits). Le sol fertile (sols volcaniques sur les plateaux) et le climat sain ont permis une occupation humaine ancienne et continue. Or, une population dense pose un problème politique : comment organiser la vie en commun, répartir les terres, trancher les conflits ? La réponse des Grassfields fut la \cle{chefferie}.

\begin{definition}[Chefferie]
Une \cle{chefferie} est une entité politique dirigée par un chef (le \cle{fo} ou \emph{fon}) dont l'autorité est à la fois \textbf{politique}, \textbf{religieuse} et \textbf{judiciaire}, et qui s'exerce sur un territoire délimité et sur une lignée de sujets qui lui doivent obéissance, travail et tribut.
\end{definition}

\courssub{Origine des chefferies : les migrations Tikar}

D'où viennent ces chefferies ? Les traditions orales des Grassfields convergent : la plupart des groupes de la région se reconnaissent une origine commune, celle des \cle{Tikar}, peuple installé sur les plateaux de l'Adamaoua camerounais et des bords de la plaine du Mbam (région de Bankim, Ngambe-Tikar).

À partir de ce foyer tikar, des \cle{vagues de migrations} successives, entre le XVII\textsuperscript{e} et le XIX\textsuperscript{e} siècle environ, ont conduit des groupes humains vers l'ouest et le nord-ouest, à la recherche de terres nouvelles, fuyant parfois les conflits ou la pression démographique. Chaque groupe migrant, conduit par un chef de lignage, s'installe sur un territoire, soumet ou absorbe les populations antérieures (dites « autochtones »), et fonde une chefferie nouvelle.

\begin{popfigure}
\begin{center}
\begin{tikzpicture}[scale=0.9]
% Frise chronologique simple : migrations Tikar
\draw[popdark, thick] (0,0) -- (14,0);
\foreach \x/\label in {0/XVII\textsuperscript{e} s., 3.5/1700, 7/1750, 10.5/1800, 14/XIX\textsuperscript{e} s.} {
  \draw[popdark, thick] (\x,-0.2) -- (\x,0.2);
  \node[popdark, font=\small, anchor=north] at (\x,-0.4) {\label};
}
% Flèches de migration
\draw[poporange, thick, ->] (1,0.5) -- (3,1.5);
\node[poporange, font=\small, anchor=south west] at (1,0.6) {Départ du foyer tikar (Adamaoua/Mbam)};
\draw[poporange, thick, ->] (4,0.5) -- (6,1.5);
\node[poporange, font=\small, anchor=south west] at (4,0.6) {Vagues successives vers l'Ouest};
\draw[poporange, thick, ->] (8,0.5) -- (10,1.5);
\node[poporange, font=\small, anchor=south west] at (8,0.6) {Fondation des chefferies (Bafoussam, Bamenda, etc.)};
\draw[poporange, thick, ->] (11,0.5) -- (13,1.5);
\node[poporange, font=\small, anchor=south west] at (11,0.6) {Consolidation avant la colonisation};
\end{tikzpicture}
\end{center}
\legende{Frise simplifiée : migrations Tikar et formation des chefferies des Grassfields (XVII\textsuperscript{e}--XIX\textsuperscript{e} siècle).}
\end{popfigure}

C'est ainsi que naquirent de très nombreuses chefferies, dont les plus célèbres sont :
\begin{itemize}
\item dans l'Ouest (pays bamiléké) : \textbf{Bafoussam}, Bandjoun, Bafang, Bangangté, Dschang, Baham, Foumban (frontière Bamoun) ;
\item dans le Nord-Ouest : \textbf{Bamenda}, \textbf{Bafut}, \textbf{Kom}, \textbf{Nso} (Banso), \textbf{Bali}, Mankon, Oku.
\end{itemize}

\begin{attention}[Ne pas confondre chefferie et lamidat]
Contrairement au lamidat peul, État de conquête islamisé et vaste, la chefferie des Grassfields est une entité de \textbf{petite taille}, issue de \textbf{migrations de groupes de parenté}, sans religion révélée ni écriture. Il n'existe pas d'« empire » des Grassfields : chaque chefferie est indépendante, même si des liens de parenté et d'alliance unissent les fo entre eux.
\end{attention}

\courssub{Le fo : un chef sacré aux trois pouvoirs}

Au sommet de la chefferie se trouve le \cle{fo} (ou \emph{fon}). Pour comprendre sa nature, il faut abandonner l'image du simple « maire » ou du simple « roi politique » : le fo est un \textbf{chef sacré}.

\begin{definition}[Le fo]
Le \cle{fo} est le chef sacré de la chefferie. Il cumule trois pouvoirs :
\begin{itemize}
\item le \textbf{pouvoir politique} : il gouverne, lève les impôts, décide de la guerre et de la paix, nomme les chefs de quartier ;
\item le \textbf{pouvoir religieux} : il est l'intermédiaire entre les vivants et les ancêtres, garant de la fertilité des terres et des femmes ; il accomplit les rites et les sacrifices (notamment lors des fêtes de la nouvelle igname) ;
\item le \textbf{pouvoir judiciaire} : il est le juge suprême, qui tranche en dernier ressort les litiges de la chefferie (vol, adultère, conflits fonciers, succession).
\end{itemize}
\end{definition}

Le pouvoir du fo se reconnaît à des \cle{symboles} précis, que l'on retrouve sur les photographies et dans les musées :
\begin{itemize}
\item le \textbf{trône} (siège royal, souvent sculpté, réservé au seul fo) ;
\item la \textbf{coupe royale}, utilisée pour les libations (vin de palme, huile) et les rites ancestraux ;
\item les \textbf{cauris}, coquillages servant de monnaie et signes de richesse et de souveraineté ;
\item la \textbf{coiffe royale} (souvent en perles ou ornée de cauris) et les \textbf{peaux de léopard}, attributs réservés au souverain, le léopard étant l'animal royal par excellence.
\end{itemize}

Le fo est entouré de règles strictes (tabous) : il ne se déplace qu'entouré de sa suite, on ne le voit pas manger en public, sa mort est longtemps tenue secrète (on dit qu'il « s'est perdu dans la brousse » ou « qu'il est parti en voyage »). Son successeur est choisi parmi ses fils par les notables (souvent le \emph{mafo} ou le \emph{ngang}), puis intronisé selon un rituel complexe qui dure plusieurs mois.

\courssub{Les contre-pouvoirs : notables et sociétés secrètes}

\begin{attention}[Erreur fréquente]
Il est faux de présenter le fo comme un \textbf{despote absolu} qui déciderait seul de tout. En réalité, son pouvoir est \textbf{encadré et contrebalancé} par les notables et par les sociétés secrètes. Un fo qui gouvernerait sans tenir compte de ces institutions perdrait toute légitimité et pourrait être abandonné par son peuple, voire destitué rituellement.
\end{attention}

Le premier contre-pouvoir est le \textbf{conseil des notables} : conseillers héréditaires (titres transmis dans certaines familles), chefs de quartier (\emph{nkoms}), princes royaux. Le fo ne prend aucune décision importante (guerre, paix, impôts, nomination) sans les consulter. Ce conseil assure la continuité de l'État en cas de vacance du pouvoir.

Le second contre-pouvoir, plus original, est constitué par les \cle{sociétés secrètes} et confréries régulatrices.

\begin{definition}[Société secrète]
Une \cle{société secrète} est une association initiatique (par exemple le \cle{Kwifon} ou \emph{Ngumba} chez les hommes, la \cle{Takumbeng} chez les femmes) qui encadre la vie sociale, \textbf{exécute les décisions} du chef et du conseil, fait respecter l'ordre et la morale, et \textbf{contrebalance le pouvoir} du fo.
\end{definition}

Le \cle{Kwifon} (appelé aussi Ngumba selon les chefferies) est la plus puissante de ces confréries : il fonctionne comme un véritable \textbf{organe exécutif et judiciaire}. Il possède sa propre case au palais, fait exécuter les sentences, perçoit les amendes, veille sur les marchés et peut même rappeler le fo à l'ordre (sanction rituelle). Ses membres agissent masqués (masques \emph{juju}), ce qui renforce leur autorité : nul ne sait qui se cache derrière le masque, ce qui les place au-dessus des personnes physiques.

\begin{savaistu}[La Takumbeng, un pouvoir féminin toujours vivant]
La \cle{Takumbeng} est une société secrète de \textbf{femmes}, composée surtout de femmes âgées (post-ménopausées). Elle intervient pour punir les atteintes graves à la morale et à la paix sociale, en particulier les violences faites aux femmes, l'inceste ou la violation des serments. Fait remarquable : la Takumbeng joue \textbf{encore aujourd'hui} un rôle de régulation sociale et politique dans le Nord-Ouest du Cameroun, et elle est intervenue publiquement lors des crises politiques récentes de la région (crise anglophone depuis 2016). Preuve que les institutions précoloniales ne sont pas des « fossiles » : elles continuent de vivre et de s'adapter.
\end{savaistu}

\courssub{La hiérarchie sociale de la chefferie}

La société des Grassfields est fortement hiérarchisée. Du sommet à la base, on distingue :
\begin{enumerate}
\item le \textbf{fo}, chef sacré, au sommet, détenteur de la \cle{royauté sacrée} ;
\item les \textbf{notables} : conseillers héréditaires, chefs de quartier, membres des grandes familles, qui forment l'encadrement politique et administratif ;
\item les \textbf{hommes libres} : la masse des cultivateurs, artisans et commerçants, organisés en lignages patrilinéaires ; ils paient l'impôt, fournissent des corvées et constituent l'armée en cas de guerre ;
\item les \textbf{captifs} : prisonniers de guerre ou personnes réduites en dépendance (par dette, crime), au bas de l'échelle, qui peuvent toutefois s'intégrer progressivement à la société par adoption ou affranchissement.
\end{enumerate}

\begin{popfigure}
\begin{center}
\begin{tikzpicture}[scale=1]
% Pyramide à 4 niveaux avec trapèzes
\fill[popgold!80] (0,3.4) -- (-1.1,2.4) -- (1.1,2.4) -- cycle;
\fill[poporange!70] (-1.1,2.4) -- (1.1,2.4) -- (2.2,1.4) -- (-2.2,1.4) -- cycle;
\fill[popblue!45] (-2.2,1.4) -- (2.2,1.4) -- (3.3,0.4) -- (-3.3,0.4) -- cycle;
\fill[popmuted!45] (-3.3,0.4) -- (3.3,0.4) -- (4.4,-0.6) -- (-4.4,-0.6) -- cycle;
\draw[popdark,thick] (0,3.4) -- (-4.4,-0.6) (0,3.4) -- (4.4,-0.6);
\node[popdark,font=\small\bfseries] at (0,2.75) {Le fo};
\node[popdark,font=\small] at (0,1.9) {Les notables (conseillers, chefs de quartier)};
\node[popdark,font=\small] at (0,0.9) {Les hommes libres (cultivateurs, artisans, commerçants)};
\node[popdark,font=\small] at (0,-0.1) {Les captifs};
\end{tikzpicture}
\end{center}
\legende{Pyramide sociale d'une chefferie des Grassfields : le fo au sommet, puis les notables, les hommes libres et les captifs à la base.}
\end{popfigure}

\courssub{Le palais (la'ah) : cœur politique, économique et culturel}

Le \cle{palais}, appelé \emph{la'ah} dans plusieurs langues de la région (bamiléké, bamoun), est bien plus que la résidence du fo : c'est le \textbf{centre politique, économique et culturel} de toute la chefferie. C'est là que se tiennent les audiences, les procès, les cérémonies rituelles (fêtes des moissons, intronisation), les danses ; c'est là que sont conservés les trésors royaux (masques, statues, trônes, cauris) et les crânes des ancêtres (skulls) ; c'est là que vivent le fo, ses nombreuses épouses (souvent plus d'une centaine pour les grands fo), ses enfants et ses serviteurs.

L'organisation spatiale du palais est elle-même un langage politique : l'entrée unique contrôle les visiteurs, la grande cour accueille les cérémonies publiques, la case du fo et celles de ses épouses occupent le cœur de l'ensemble, et une case spéciale est réservée aux sociétés secrètes, dont le Kwifon.

\begin{popfigure}
\begin{center}
\begin{tikzpicture}[scale=1]
% Enceinte du palais
\draw[popdark,very thick,rounded corners] (-4.6,-3) rectangle (4.6,3.2);
\node[popdark,font=\small\itshape] at (0,3.55) {Enclos du palais (la'ah)};
% Entrée unique
\draw[popdark,thick] (-0.9,-3) -- (-0.9,-3.8) (0.9,-3) -- (0.9,-3.8);
\draw[popdark,thick,->] (0,-4.4) -- (0,-3.6);
\node[popdark,font=\small] at (0,-4.7) {Entrée unique (portail gardé)};
% Cour des cérémonies
\draw[fill=popblue!25,draw=popblue,thick] (-2.6,-2.6) rectangle (2.6,0.2);
\node[popdark,font=\small\bfseries] at (0,-1.2) {Cour des cérémonies};
\node[popdark,font=\footnotesize\itshape] at (0,-1.8) {(audiences, danses, marché des jours de fête)};
% Case du fo
\draw[fill=popgold!60,draw=popdark,thick] (-1.4,1.4) rectangle (1.4,2.9);
\node[popdark,font=\small\bfseries] at (0,2.15) {Case du fo};
% Cases des épouses (gauche)
\draw[fill=poporange!40,draw=poporange,thick] (-4.2,1.4) rectangle (-1.9,2.9);
\node[popdark,font=\small] at (-3.05,2.35) {Cases des};
\node[popdark,font=\small] at (-3.05,1.95) {épouses};
% Case des sociétés secrètes (droite)
\draw[fill=poppurple!30,draw=poppurple,thick] (1.9,1.4) rectangle (4.2,2.9);
\node[popdark,font=\small] at (3.05,2.35) {Case des sociétés};
\node[popdark,font=\small] at (3.05,1.95) {secrètes (Kwifon)};
% Greniers (bas gauche)
\draw[fill=popgreen!25,draw=popgreen,thick] (-4.2,-0.1) rectangle (-3.0,1.0);
\node[popdark,font=\footnotesize] at (-3.6,0.45) {Greniers};
% Ateliers (bas droite)
\draw[fill=popgreen!25,draw=popgreen,thick] (3.0,-0.1) rectangle (4.2,1.0);
\node[popdark,font=\footnotesize] at (3.6,0.45) {Ateliers};
\end{tikzpicture}
\end{center}
\legende{Plan schématique d'un palais de chefferie des Grassfields : 1. entrée unique ; 2. cour des cérémonies ; 3. case du fo ; 4. cases des épouses ; 5. case des sociétés secrètes ; 6. greniers et ateliers. Croquis simplifié, sans échelle.}
\end{popfigure}

\courssub{La vie économique : agriculture, artisanat et marchés}

La prospérité de la chefferie repose d'abord sur l'\cle{agriculture}. Les femmes, principales cultivatrices, produisent le \textbf{maïs}, le \textbf{taro} (macabo), les \textbf{ignames}, ainsi que les haricots et les arachides. Les hommes défrichent, construisent les cases et les greniers, et pratiquent un peu d'élevage (chèvres, moutons, poules).

L'\cle{artisanat} est très développé et souvent réservé à certains lignages spécialisés (castes d'artisans) :
\begin{itemize}
\item la \textbf{sculpture} sur bois (masques, statues royales, tabourets, pipes, portes sculptées) ;
\item la \textbf{vannerie} (paniers, nattes, tamis) et la poterie (jarres, canaris) ;
\item la \textbf{forge} : le forgeron fabrique houes, couteaux, lances, cloches, bracelets ; c'est un personnage respecté et parfois craint (maîtrise du feu et du fer).
\end{itemize}

Les échanges s'organisent autour des \cle{marchés hebdomadaires}, qui suivent la semaine traditionnelle de huit jours propre à plusieurs groupes de la région (semaine bamiléké : \emph{Nkwo, Eke, Orie, Afo}...). On y troque et on y achète denrées, bétail, objets d'artisanat. La monnaie utilisée est le \cle{cauri} (\emph{Cypraea moneta}), petit coquillage importé des côtes de l'océan Indien (via les réseaux sahéliens), qui sert d'unité de compte et de moyen de paiement, complété par le troc et, dans certaines zones, par les barres de fer (\emph{ngwa}) et les étoffes (\emph{ndop}, \emph{toghu}).

\courssub{Étude de document : lire une photographie de chefferie}

\begin{methode}[Analyser une photographie historique]
\begin{enumerate}
\item \textbf{Décrire} ce que l'on voit, sans interpréter : lieux, bâtiments, personnages, objets, attitudes, vêtements.
\item \textbf{Identifier} les personnages (chef, notables, danseurs, foule) et les symboles du pouvoir (trône, coupe, coiffe, cauris, peau de léopard, masques).
\item \textbf{Situer} le document : nature (photographie coloniale, carte postale, cliché d'anthropologue), époque, auteur probable, but de la prise de vue.
\item \textbf{Dégager} ce que la photographie révèle du pouvoir : centralité du fo, rôle des sociétés secrètes, richesse et organisation de la chefferie.
\item \textbf{Conclure} en rappelant les limites du document (regard extérieur, mise en scène possible, cadrage partiel).
\end{enumerate}
\end{methode}

\begin{exoresolu}[Photographie d'un palais et d'une danse rituelle]
\textbf{Énoncé.} On donne à des élèves deux documents : (1) une photographie montrant un vaste ensemble de cases à toits coniques entouré d'une palissade, avec une grande cour où est rassemblée une foule ; (2) une photographie d'une danse rituelle où des danseurs masqués entourent un personnage assis sur un siège sculpté, coiffé d'un bonnet orné de cauris et tenant une coupe. À l'aide de la méthode d'analyse, dégagez ce que ces documents révèlent de l'organisation politique des chefferies des Grassfields.
\tcblower
\textbf{Solution détaillée.}
\begin{enumerate}
\item \textbf{Description.} Document 1 : un ensemble de nombreuses cases groupées dans un enclos (palissade), une cour centrale spacieuse, une foule nombreuse rassemblée. Document 2 : des danseurs masqués (masques en bois, costumes de fibres), un personnage central assis sur un siège sculpté à dossier haut, portant une coiffe ornée de cauris et tenant une coupe à la main.
\item \textbf{Identification.} L'ensemble de cases est le palais (la'ah), cœur de la chefferie ; la foule rassemblée indique une cérémonie publique (fête, audience). Le personnage assis est le fo : le siège sculpté est le trône royal, la coiffe de cauris et la coupe sont les insignes du pouvoir sacré ; les danseurs masqués sont les membres d'une société secrète (Kwifon ou société de danse régulatrice comme \emph{Njom}).
\item \textbf{Interprétation.} Le document 1 montre que le palais est le centre politique et culturel de la chefferie, capable de rassembler tout le peuple (légitimité populaire). Le document 2 montre que le fo est un chef sacré, dont l'autorité s'exprime par des symboles visibles (trône, cauris, coupe), et que les sociétés secrètes encadrent la vie publique aux côtés du chef (contre-pouvoir rituel).
\item \textbf{Conclusion.} Ces photographies confirment que les Grassfields, loin d'être des sociétés « sans État », possédaient des institutions politiques organisées : un chef sacré aux pouvoirs politique, religieux et judiciaire, des symboles de souveraineté, et des contre-pouvoirs incarnés par les confréries masquées. Limite : ces photographies, souvent prises par des colonisateurs, des missionnaires ou des administrateurs (début XX\textsuperscript{e} s.), peuvent refléter un regard extérieur, une mise en scène ou un cadrage orienté.
\end{enumerate}
\end{exoresolu}

\begin{aretenir}[L'essentiel de la section]
\begin{itemize}
\item Les Grassfields (Ouest et Nord-Ouest) sont une région de hautes terres à forte densité de population (sols volcaniques, climat tempéré).
\item Les chefferies (Bafoussam, Bamenda, Bafut, Kom, Nso, Bali, Bandjoun...) sont issues des migrations des groupes \cle{Tikar} (XVII\textsuperscript{e}--XIX\textsuperscript{e} s.).
\item Le \cle{fo} est un chef sacré cumulant les pouvoirs politique, religieux et judiciaire ; ses symboles sont le trône, la coupe, les cauris, la coiffe royale et la peau de léopard.
\item Son pouvoir est encadré par le conseil des notables et les sociétés secrètes (\cle{Kwifon}/Ngumba pour les hommes, \cle{Takumbeng} pour les femmes).
\item La société est hiérarchisée : fo, notables, hommes libres, captifs.
\item Le palais (la'ah) est le centre politique, économique et culturel ; l'économie repose sur l'agriculture (féminine), l'artisanat (lignages spécialisés), les marchés hebdomadaires (semaine de 8 jours) et la monnaie-cauri.
\end{itemize}
\end{aretenir}

\begin{exercice}[Pour s'entraîner]
\begin{enumerate}
\item Définis les termes suivants : chefferie, fo, société secrète, cauri, la'ah.
\item Explique en cinq lignes pourquoi on ne peut pas qualifier le fo de « dictateur ».
\item Compare en un tableau de deux colonnes le lamidat peul et la chefferie des Grassfields (origine, taille, religion, chef, contre-pouvoirs).
\item \textbf{Travail de terrain (TD) :} Si ton établissement se trouve dans une zone de chefferie, rends-toi au palais local (avec l'autorisation de l'administration) et interroge un notable ou un guide sur : l'origine de la chefferie, le rôle actuel du fo, le fonctionnement du Kwifon ou de la Takumbeng. Rédige un compte-rendu structuré (2 pages max). Si ton établissement n'est pas en zone de chefferie, fais une recherche documentaire sur la chefferie la plus proche.
\end{enumerate}
\end{exercice}

\begin{corrige}[Exercice]
\begin{enumerate}
\item \textbf{Chefferie} : entité politique dirigée par un chef (fo) dont l'autorité est politique, religieuse et judiciaire, exercée sur un territoire et une lignée de sujets. \textbf{Fo} : chef sacré de la chefferie, détenteur des trois pouvoirs. \textbf{Société secrète} : association initiatique (Kwifon, Takumbeng) qui exécute les décisions, encadre la vie sociale et contrebalance le pouvoir du chef. \textbf{Cauri} : coquillage (\emph{Cypraea moneta}) servant de monnaie d'échange et de signe de richesse. \textbf{La'ah} : palais royal, centre politique, économique et culturel de la chefferie.
\item Le fo n'est pas un dictateur car son pouvoir est limité par le conseil des notables, qu'il doit consulter pour toute décision majeure, et par les sociétés secrètes comme le Kwifon, qui exécutent les décisions mais peuvent aussi rappeler le chef à l'ordre, voire le sanctionner rituellement. Un fo qui ignorerait ces institutions perdrait sa légitimité sacrée et politique.
\item 
\begin{center}
\begin{tabularx}{\linewidth}{|l|X|X|}
\hline
\rowcolor{popblueL} \textbf{Critère} & \textbf{Lamidat peul} & \textbf{Chefferie des Grassfields} \\
\hline
Origine & Conquête (djihad du XIX\textsuperscript{e} s., Ousman dan Fodio) & Migrations des groupes Tikar (XVII\textsuperscript{e}--XIX\textsuperscript{e} s.) \\
\hline
Taille & Vaste État (plusieurs chefferies vassales) & Petite entité locale (quelques km² à quelques centaines de km²) \\
\hline
Religion & Islam (charia, écriture arabe) & Religions traditionnelles (culte des ancêtres, oralité) \\
\hline
Chef & Lamido, souverain musulman, calife local & Fo, chef sacré aux trois pouvoirs (politique, religieux, judiciaire) \\
\hline
Contre-pouvoirs & Conseil de vizirs (conseil de la cour), cadis & Conseil des notables, sociétés secrètes (Kwifon, Takumbeng) \\
\hline
\end{tabularx}
\end{center}
\end{enumerate}
\end{corrige}

\coursec{Le sultanat bamoun}

Après avoir étudié les chefferies des Grassfields, ces entités politiques de taille modeste mais solidement structurées, nous nous tournons maintenant vers l'ouest des hauts plateaux, là où les Grassfields rencontrent les savanes du Nord. C'est là que s'est développé l'un des États les plus remarquables du Cameroun précolonial : le \cle{royaume bamoun}, devenu \cle{sultanat} au début du XX\textsuperscript{e} siècle. Contrairement aux chefferies bamiléké, souvent rivales entre elles, le royaume bamoun a su unifier sous l'autorité d'un seul souverain un vaste territoire, une armée permanente et une administration centralisée. Mieux encore : à la veille de la colonisation, sous le règne du roi \cle{Njoya}, cet État africain a connu une modernisation endogène — écriture, école, architecture — qui en fait un cas unique en Afrique centrale.

\courssub{Origine et fondation du royaume bamoun}

\textbf{Une migration depuis les terres Tikar.}\par
Selon la tradition orale bamoun, recueillie et mise par écrit par le roi Njoya lui-même, le royaume trouve son origine dans une migration partie des \cle{terres Tikar}, situées dans l'actuelle région de l'Adamaoua et des hauts plateaux de l'Ouest. Un groupe de guerriers, conduit par le prince \cle{Nchare Yen}, aurait quitté ces terres en quête de nouveaux espaces à conquérir. La tradition fixe la fondation du royaume vers \cle{1394}, date symbolique retenue par les chroniques bamoun.

En chemin, les migrants de Nchare Yen soumettent ou absorbent les populations installées sur les bords de la rivière \cle{Noun}. Nchare Yen devient le premier \cle{mfon} (roi) et fonde la dynastie qui régnera sans interruption jusqu'à la colonisation. La capitale est fixée à \cle{Foumban}, ville qui demeure aujourd'hui encore le cœur politique et culturel du pays bamoun.

\begin{definition}[Mfon]
Le \cle{mfon} est le roi des Bamouns. Souverain politique, militaire et religieux, il concentre tous les pouvoirs et règne depuis la capitale Foumban. Après la conversion du roi Njoya à l'islam, le titre de \emph{mfon} est complété par celui de \cle{sultan}.
\end{definition}

\textbf{Un royaume né de la conquête.}\par
Contrairement aux chefferies des Grassfields, qui se structurent autour d'une parenté ou d'un lignage, le royaume bamoun naît d'une \emph{conquête militaire}. Cette origine guerrière marque durablement l'État : l'armée y joue un rôle central, et l'expansion territoriale reste pendant des siècles l'objectif des souverains. Les rois successeurs de Nchare Yen, notamment \cle{Mbuembue} (au début du XIX\textsuperscript{e} siècle), mènent de grandes campagnes qui portent le royaume à son apogée.

\courssub{La nature de l'État bamoun : un royaume centralisé et militaire}

\begin{definition}[Sultanat]
Un \cle{sultanat} est un État dirigé par un \cle{sultan}, c'est-à-dire un souverain musulman. Le royaume bamoun devient un sultanat après la conversion du roi Njoya à l'islam, vers 1916-1918 : le pouvoir politique reste le même, mais il s'inscrit désormais dans le cadre religieux et culturel de l'islam.
\end{definition}

\textbf{Une organisation politique hiérarchisée.}\par
L'État bamoun repose sur une pyramide de pouvoir clairement organisée :

\begin{itemize}
\item au sommet, le \cle{mfon}, souverain absolu, chef des armées et garant de la prospérité du royaume ;
\item le roi est assisté de conseillers influents, parmi lesquels la \cle{mère du roi}, qui joue un rôle politique de premier plan, et de hauts dignitaires chargés notamment de la justice et de l'encadrement des princes ;
\item le territoire est divisé en \cle{provinces}, confiées à des chefs nommés par le roi et responsables devant lui ;
\item enfin, le royaume dispose d'une \cle{armée permanente}, véritable colonne vertébrale de l'État, qui assure les conquêtes, la défense des frontières et la sécurité intérieure.
\end{itemize}

\begin{popfigure}
\begin{center}
\begin{tikzpicture}[scale=1]
% Contour simplifié du royaume bamoun
\fill[popgoldL] (0,0) .. controls (1.2,1.4) and (2.6,1.8) .. (4,1.5)
 .. controls (5.4,1.2) and (6,0.2) .. (5.6,-1)
 .. controls (5.2,-2.2) and (3.6,-2.6) .. (2.2,-2.3)
 .. controls (0.8,-2) and (-0.6,-1.2) .. (0,0);
\draw[popdark,thick] (0,0) .. controls (1.2,1.4) and (2.6,1.8) .. (4,1.5)
 .. controls (5.4,1.2) and (6,0.2) .. (5.6,-1)
 .. controls (5.2,-2.2) and (3.6,-2.6) .. (2.2,-2.3)
 .. controls (0.8,-2) and (-0.6,-1.2) .. (0,0);
% Capitale Foumban
\fill[popink] (2.8,0.1) circle (0.09);
\node[popdark,font=\bfseries\small] at (2.8,0.45) {FOUMBAN (capitale)};
\draw[popink] (2.8,0.1) -- (2.8,0.38);
% Rivières
\draw[popblue,thick] (-0.8,1.8) .. controls (0.5,1) and (1,0.4) .. (1.4,-0.8) .. controls (1.7,-1.8) and (1.2,-2.6) .. (0.8,-3.1);
\node[popblue,font=\small] at (-0.6,2.1) {le Noun};
\draw[popblue,thick] (6.8,2.2) .. controls (6.2,1.2) and (6.4,0) .. (6.6,-1.2) .. controls (6.8,-2.2) and (6.4,-2.9) .. (6,-3.4);
\node[popblue,font=\small] at (7.1,2.4) {le Mbam};
% Voisins
\node[popgreen,font=\small\itshape] at (3.2,-3.4) {Chefferies bamiléké (Grassfields)};
\node[popgreen,font=\small\itshape] at (3.4,2.6) {Lamidats peuls (Ngaoundéré, Tibati)};
% Flèches d'expansion
\draw[-{Stealth[length=3mm]},poporange,very thick] (2.8,0.1) -- (4.8,1);
\draw[-{Stealth[length=3mm]},poporange,very thick] (2.8,0.1) -- (1,-1.6);
\draw[-{Stealth[length=3mm]},poporange,very thick] (2.8,0.1) -- (4.6,-1.8);
% Nord
\draw[-{Stealth[length=2.5mm]},popdark,thick] (-1.6,-3.2) -- (-1.6,-2.4);
\node[popdark,font=\small] at (-1.6,-3.5) {N};
% Échelle
\draw[popdark,thick] (7.6,-3.4) -- (9.1,-3.4);
\draw[popdark,thick] (7.6,-3.3) -- (7.6,-3.5);
\draw[popdark,thick] (9.1,-3.3) -- (9.1,-3.5);
\node[popdark,font=\footnotesize] at (8.35,-3.7) {50 km (indicatif)};
\end{tikzpicture}
\end{center}
\legende{Carte schématique du royaume bamoun à son apogée (XIX\textsuperscript{e} siècle). 1 : territoire contrôlé entre le Noun et le Mbam (zone dorée) ; 2 : capitale Foumban ; 3 : flèches d'expansion militaire ; 4 : voisins du royaume. Croquis indicatif, sans précision cartographique.}
\end{popfigure}

\textbf{Conquêtes et extension territoriale.}\par
Grâce à son armée, le royaume bamoun s'étend progressivement. À son \cle{apogée}, au XIX\textsuperscript{e} siècle sous le règne du roi Mbuembue, il contrôle un vaste territoire situé \emph{entre la rivière Noun à l'ouest et la rivière Mbam à l'est}, soit plusieurs milliers de kilomètres carrés. Les peuples voisins sont soumis, deviennent tributaires, ou sont repoussés. Cette puissance explique que ni les chefferies bamiléké au sud, ni les lamidats peuls au nord, n'aient jamais réussi à conquérir le pays bamoun : le royaume a notamment résisté aux cavaliers peuls issus du djihad d'Ousman dan Fodio au début du XIX\textsuperscript{e} siècle.

\courssub{L'islamisation : du royaume au sultanat}

\textbf{Une conversion tardive et pacifique.}\par
L'islam pénètre au pays bamoun par les échanges commerciaux et les contacts avec les Peuls du Nord. Mais c'est le roi \cle{Njoya} qui franchit le pas : vers \cle{1916-1918}, il se convertit officiellement à l'islam, entraînant avec lui une partie de sa cour et de son peuple. Le royaume bamoun devient alors un \cle{sultanat}, et le mfon prend le titre de sultan.

\begin{attention}[Ne pas confondre les deux islamisations]
Ne situe pas l'islamisation bamoun au même moment que le \cle{djihad peul}. Au Nord-Cameroun, l'islam s'impose au début du XIX\textsuperscript{e} siècle par la \emph{conquête militaire} des Peuls d'Ousman dan Fodio. Chez les Bamouns, la conversion intervient un siècle plus tard, au \emph{début du XX\textsuperscript{e} siècle}, de manière \emph{pacifique} et par la volonté du roi Njoya, \textbf{sans conquête peule}. Au contraire, le royaume bamoun avait victorieusement résisté aux Peuls au XIX\textsuperscript{e} siècle.
\end{attention}

\courssub{Le roi Njoya (vers 1876-1933) : un souverain modernisateur}

\textbf{Une figure majeure de l'histoire du Cameroun.}\par
Né vers 1876, Njoya accède au trône vers \cle{1886}, encore adolescent, après la mort de son père tué à la guerre. La régence est assurée par sa mère, qui protège le jeune roi des prétendants rivaux. Son long règne chevauche deux époques : la fin de l'indépendance bamoun et l'arrivée des colonisateurs allemands (1902) puis français (1916). En \cle{1931}, les autorités coloniales françaises le destituent et l'exilent à Yaoundé, où il meurt en \cle{1933}.

\textbf{Un inventeur : l'écriture Shümom.}\par
La réalisation la plus célèbre de Njoya est la création d'une \cle{écriture} pour la langue bamoun, appelée \cle{Shümom}. Élaborée à partir de 1895 environ et perfectionnée dans les années qui suivent, elle permet de mettre par écrit l'histoire, les lois et la pharmacopée du royaume.

\begin{exemplebox}[L'écriture Shümom en chiffres]
L'écriture Shümom, créée par le roi Njoya à partir de \cle{1895} environ, comptait à l'origine \textbf{plus de 400 signes} (des idéogrammes représentant des mots entiers). Njoya et ses scribes l'ont ensuite simplifiée en plusieurs étapes, jusqu'à un système syllabique d'\textbf{environ 80 signes}, beaucoup plus facile à apprendre. C'est l'une des très rares écritures inventées en Afrique subsaharienne avant ou pendant la colonisation.
\end{exemplebox}

\textbf{Un bâtisseur et un promoteur de l'école.}\par
Njoya fait construire à Foumban, en \cle{1917}, un somptueux \cle{palais} de briques, inspiré à la fois de l'architecture locale et des modèles coloniaux : c'est le célèbre palais de Foumban, toujours debout aujourd'hui. Il fonde également des \cle{écoles} où l'on enseigne l'écriture Shümom, et fait rédiger par ses scribes une \emph{histoire du peuple bamoun}, précieuse source pour les historiens. Sa cour, raffinée, accueille artistes, scribes, forgerons et commerçants étrangers.

\begin{popfigure}
\begin{center}
\begin{tikzpicture}[scale=1]
% Axe chronologique
\draw[-{Stealth[length=3mm]},popdark,very thick] (0,0) -- (13.5,0);
% Graduations
\foreach \x in {0.5, 3.2, 5.9, 8.2, 9.5, 11.6, 12.9}
  \draw[popdark,thick] (\x,-0.12) -- (\x,0.12);
\node[popdark,font=\footnotesize,below] at (0.5,-0.15) {1886};
\node[popdark,font=\footnotesize,below] at (3.2,-0.15) {1895};
\node[popdark,font=\footnotesize,below] at (5.9,-0.15) {1903};
\node[popdark,font=\footnotesize,below] at (8.2,-0.15) {1916};
\node[popdark,font=\footnotesize,below] at (9.5,-0.15) {1917};
\node[popdark,font=\footnotesize,below] at (11.6,-0.15) {1931};
\node[popdark,font=\footnotesize,below] at (12.9,-0.15) {1933};
% Événements
\node[popblue,font=\footnotesize,align=center,above] at (0.5,0.25) {Début du\\règne de Njoya};
\draw[popblue] (0.5,0.12) -- (0.5,0.3);
\node[popgreen,font=\footnotesize,align=center,above] at (3.2,0.25) {Début de\\l'écriture Shümom};
\draw[popgreen] (3.2,0.12) -- (3.2,0.3);
\node[popgreen,font=\footnotesize,align=center,below] at (5.9,-0.55) {Shümom simplifiée\\(environ 80 signes)};
\draw[popgreen] (5.9,-0.12) -- (5.9,-0.5);
\node[poporange,font=\footnotesize,align=center,above] at (8.2,0.25) {Conversion à l'islam\\(1916-1918)};
\draw[poporange] (8.2,0.12) -- (8.2,0.3);
\node[poppurple,font=\footnotesize,align=center,below] at (9.5,-0.55) {Construction du\\palais de Foumban};
\draw[poppurple] (9.5,-0.12) -- (9.5,-0.5);
\node[popink,font=\footnotesize,align=center,above] at (11.6,0.25) {Destitution et exil\\par les Français};
\draw[popink] (11.6,0.12) -- (11.6,0.3);
\node[popink,font=\footnotesize,align=center,below] at (12.9,-0.55) {Mort de Njoya\\à Yaoundé};
\draw[popink] (12.9,-0.12) -- (12.9,-0.5);
\end{tikzpicture}
\end{center}
\legende{Frise chronologique du règne du roi Njoya (1886-1931) et de ses principales réalisations : invention de l'écriture Shümom, islamisation du royaume, construction du palais de Foumban, puis destitution (1931) et mort en exil (1933).}
\end{popfigure}

\begin{savaistu}[Le palais de Foumban, un musée vivant]
Le palais de Foumban, construit en \cle{1917} sur l'ordre du sultan Njoya, abrite aujourd'hui un \cle{musée} qui conserve les trésors de la dynastie bamoun : trônes, armes, masques, et surtout les \emph{manuscrits écrits en Shümom}, dont l'histoire du royaume rédigée par les scribes de Njoya. Le palais est inscrit sur la liste indicative du patrimoine mondial de l'UNESCO.
\end{savaistu}

\courssub{Vie économique et culturelle du royaume}

\textbf{Une économie prospère.}\par
La richesse du royaume bamoun repose sur plusieurs piliers :

\begin{itemize}
\item l'\cle{agriculture} : maïs, mil, sorgho, arachides et ignames nourrissent une population dense ; les paysans versent des tributs au roi ;
\item l'\cle{artisanat d'art}, réputé dans toute la région : fonte du \cle{bronze} (statuettes, pipes), sculpture sur \cle{bois} (masques, trônes, portes sculptées), travail des \cle{perles} et tissage ;
\item le \cle{commerce} : Foumban est un carrefour commercial où convergent les produits du Nord (sel, kola, tissus) et ceux des Grassfields ; les marchands bamoun entretiennent des relations avec les Peuls, les Bamiléké et jusqu'à la côte.
\end{itemize}

\textbf{Une cour raffinée.}\par
La cour de Foumban est un centre de création artistique et intellectuelle. Le roi entretient des artistes, des musiciens, des scribes et des devins. Cette politique de mécénat atteint son sommet sous Njoya, qui fait de sa capitale un véritable centre culturel, admiré par les premiers observateurs européens.

\courssub{Étude de document : un État africain moderne à la veille de la colonisation}

\textbf{Description des documents.}\par
Le dossier documentaire comprend deux pièces :
\begin{itemize}
\item \emph{Document 1} : une photographie du \cle{palais de Foumban}, construit en 1917. On y voit un imposant bâtiment de briques à deux niveaux, coiffé d'une tour centrale surmontée du double serpent, symbole du pouvoir bamoun. L'architecture combine des techniques locales et des influences extérieures.
\item \emph{Document 2} : un portrait du \cle{sultan Njoya}, assis en tenue royale, entouré des insignes du pouvoir (coiffe, sceptre, perles).
\end{itemize}

\begin{exoresolu}[Commenter les deux documents : le royaume bamoun, un État moderne]
\textbf{Consigne} : à partir des deux documents et de tes connaissances, montre que le royaume bamoun était, à la veille de la colonisation, un État africain organisé et ouvert à la modernité. (Environ 15 lignes.)
\tcblower
\textbf{Corrigé rédigé.}\par
Les deux documents présentent le royaume bamoun au début du XX\textsuperscript{e} siècle. Le document 1 est une photographie du palais de Foumban, construit en 1917 par le roi Njoya ; le document 2 est le portrait de ce même souverain, qui régna de 1886 à 1931.\par
Ces documents montrent d'abord un État \emph{politiquement organisé} : le palais, siège du pouvoir royal, témoigne de l'existence d'une autorité centrale forte, capable de mobiliser des ressources et des artisans pour un tel chantier. Le portrait du sultan, entouré des insignes du pouvoir, confirme l'existence d'une monarchie hiérarchisée, à la tête d'un royaume centralisé disposant d'une armée permanente et de provinces administrées.\par
Ils révèlent ensuite un État \emph{ouvert à la modernité} : l'architecture du palais combine savoir-faire locaux et influences extérieures, signe que le royaume n'est pas fermé mais en contact avec le monde. Cette ouverture est confirmée par les autres réalisations de Njoya : l'invention de l'écriture Shümom, la création d'écoles, la conversion à l'islam.\par
En conclusion, ces documents prouvent que le Cameroun précolonial n'était pas un espace désorganisé : le royaume bamoun était un État moderne, dirigé par un souverain bâtisseur et réformateur, ce qui dément l'idée coloniale d'une Afrique « sans histoire ».
\end{exoresolu}

\begin{attention}[Piège de rédaction au commentaire de document]
Ne te contente pas de \emph{décrire} la photographie (« on voit un grand bâtiment »). Le correcteur attend une \emph{analyse} : chaque observation doit servir un argument. Relie toujours le document à une notion du cours (État centralisé, modernisation, souverain réformateur) et date les faits avec précision (palais : 1917 ; règne de Njoya : 1886-1931).
\end{attention}

\begin{aretenir}[Complément utile au Probatoire et au Baccalauréat technique]
Le roi \cle{Njoya} est souvent cité comme l'exemple même de la \cle{modernisation africaine endogène} face à la colonisation : bien avant et pendant l'arrivée des Européens, il a inventé une écriture, fondé des écoles, fait rédiger l'histoire de son peuple et construit un palais moderne. À l'examen, tu peux mobiliser cet exemple pour montrer que les sociétés africaines précoloniales disposaient d'États structurés et capables d'innovation. Retiens les dates clés : fondation du royaume vers 1394 (Nchare Yen), règne de Njoya 1886-1931, écriture Shümom à partir de 1895 environ, conversion à l'islam 1916-1918, palais de Foumban 1917, mort de Njoya en exil en 1933.
\end{aretenir}

\begin{exercice}[Pour s'entraîner]
1. Définis les termes suivants : \emph{mfon}, \emph{sultanat}, \emph{Shümom}.\par
2. Explique en quelques lignes pourquoi l'islamisation du royaume bamoun diffère de celle des lamidats peuls du Nord-Cameroun.\par
3. Rédige un paragraphe argumenté (10 à 15 lignes) répondant à la question : « En quoi le roi Njoya incarne-t-il la modernisation endogène d'un État africain ? »
\end{exercice}

\begin{corrige}[Exercice — Le sultanat bamoun]
1. Le \emph{mfon} est le roi des Bamouns, souverain politique, militaire et religieux régnant depuis Foumban. Le \emph{sultanat} est un État dirigé par un sultan, souverain musulman : le royaume bamoun le devient après la conversion de Njoya vers 1916-1918. Le \emph{Shümom} est l'écriture inventée par Njoya à partir de 1895 environ, passée de plus de 400 signes à environ 80.\par
2. Au Nord-Cameroun, l'islam s'impose au début du XIX\textsuperscript{e} siècle par la conquête militaire des Peuls (djihad d'Ousman dan Fodio), qui fondent des lamidats. Chez les Bamouns, la conversion est tardive (début du XX\textsuperscript{e} siècle), pacifique et décidée par le roi Njoya lui-même, sans aucune conquête peule : le royaume avait d'ailleurs résisté victorieusement aux Peuls au siècle précédent.\par
3. \emph{Exemple de paragraphe} : Njoya incarne la modernisation endogène car ses réformes naissent de sa propre initiative et non d'une imposition coloniale. Avant même l'arrivée effective des Allemands, il invente l'écriture Shümom et la simplifie pour la diffuser, fonde des écoles et fait rédiger l'histoire de son peuple. Bâtisseur, il fait ériger en 1917 le palais de Foumban, synthèse d'architecture locale et de techniques nouvelles. Enfin, sa conversion à l'islam montre sa capacité à intégrer des apports extérieurs en les adaptant. Njoya prouve ainsi qu'un État africain pouvait se moderniser de l'intérieur, tout en restant maître de ses choix.
\end{corrige}

\coursec{Étude comparée : trois modèles d'organisation politique}

Nous venons d'étudier successivement le \cle{lamidat} peul du Nord-Cameroun, la \cle{chefferie} des Grassfields et le \cle{sultanat bamoun}. Chacune de ces entités possède son histoire propre, ses institutions et ses coutumes. Mais l'historien ne s'arrête pas à la description : il compare. Comparer, c'est chercher à la fois ce qui rapproche et ce qui distingue plusieurs réalités, afin de mieux comprendre chacune d'elles. C'est l'objet de cette section : confronter les trois modèles d'organisation politique étudiés pour dégager leurs points communs et leurs spécificités, puis mesurer la diversité politique du Cameroun à la veille de la colonisation.

\courssub{Pourquoi et comment comparer ?}

\begin{methode}[Construire un tableau comparatif]
Comparer trois entités politiques exige une démarche rigoureuse, en quatre étapes :
\begin{enumerate}
\item \textbf{Définir des critères communs} : on ne peut comparer que des éléments de même nature. Nous retiendrons six critères : l'\cle{origine} de l'entité, la \textbf{nature du pouvoir}, la \textbf{religion}, la \textbf{hiérarchie sociale}, l'\textbf{économie} et le \textbf{rapport à l'islam}.
\item \textbf{Remplir le tableau à partir des documents} : pour chaque critère, on reporte les informations tirées des leçons précédentes, sans invention ni extrapolation.
\item \textbf{Dégager les ressemblances} : on repère les lignes du tableau où les réponses sont voisines (points communs).
\item \textbf{Dégager les différences} : on repère les lignes où les réponses divergent (spécificités), puis on les explique par l'histoire propre de chaque entité.
\end{enumerate}
\end{methode}

\begin{attention}[Éviter les généralisations abusives]
Le piège classique consiste à plaquer un modèle unique sur tout le Cameroun précolonial : « tous les chefs étaient des sultans », « toutes les populations étaient musulmanes », « toutes les sociétés avaient des esclaves ». C'est faux. Chaque entité a sa propre histoire : le lamidat de Rey-Bouba n'est pas la chefferie de Bafou, qui n'est pas le royaume de Foumban. Au Probatoire et au Baccalauréat technique, toute généralisation non nuancée est sanctionnée. Il faut toujours préciser \emph{de quelle entité} l'on parle et \emph{à quelle époque}.
\end{attention}

\courssub{Le tableau comparatif}

Appliquons la méthode. À partir des acquis des sections précédentes, nous construisons le tableau suivant.

\begin{popfigure}
\begin{tikzpicture}[x=1cm,y=1cm]
% En-têtes
\fill[popblue] (0,0) rectangle (3.2,1);
\fill[poporange] (3.2,0) rectangle (7.2,1);
\fill[popgreen] (7.2,0) rectangle (11.2,1);
\fill[poppurple] (11.2,0) rectangle (15.2,1);
\node[white,font=\bfseries] at (1.6,0.5) {Critère};
\node[white,font=\bfseries] at (5.2,0.5) {Lamidat (Nord)};
\node[white,font=\bfseries] at (9.2,0.5) {Chefferie (Grassfields)};
\node[white,font=\bfseries] at (13.2,0.5) {Sultanat bamoun};
% Lignes
\foreach \y/\crit/\lam/\chef/\sult in {
 -1.1/{Origine}/{Conquête peule (jihad de Modibo Adama, début XIX\textsuperscript{e} s.)}/{Chefferies anciennes issues de migrations et de fusions de groupes}/{Fondation par Ncharé Yen (XVI\textsuperscript{e}--XVII\textsuperscript{e} s.), dynastie bamoun},
 -2.2/{Chef}/{Lamido, nommé et confirmé par le califat de Sokoto}/{Fon (chef), appuyé par les sociétés secrètes (Takembeng, Kwifon)}/{Sultan (Mfon), monarque absolu et conquérant},
 -3.3/{Religion}/{Islam religion d'État}/{Religions traditionnelles (culte des ancêtres)}/{Islam adopté sous Njoya, coexistant avec les traditions},
 -4.4/{Institutions}/{Cour du lamido, fiefs, cadis (juges musulmans)}/{Palais, sociétés régulatrices, conseil des notables}/{Cour hiérarchisée, armée, écriture inventée (A-ka-u-ku)},
 -5.5/{Économie}/{Agriculture, élevage, esclavage, caravanes}/{Agriculture vivrière, marchés périodiques, artisanat}/{Agriculture, artisanat de cour, commerce longue distance},
 -6.6/{Rapport à l'islam}/{Entité issue du jihad, islamisée d'emblée}/{Islam absent ou marginal}/{Islam adopté tardivement et volontairement (vers 1916)}
}{
 \node[anchor=west,font=\bfseries\small,text=popdark] at (0.1,\y) {\crit};
 \node[anchor=west,font=\scriptsize,text width=3.7cm] at (3.35,\y) {\lam};
 \node[anchor=west,font=\scriptsize,text width=3.7cm] at (7.35,\y) {\chef};
 \node[anchor=west,font=\scriptsize,text width=3.7cm] at (11.35,\y) {\sult};
 \draw[popline] (0,\y+0.55) -- (15.2,\y+0.55);
}
\draw[popline] (0,0) -- (15.2,0);
\draw[popline] (0,-7.15) -- (15.2,-7.15);
\foreach \x in {0,3.2,7.2,11.2,15.2} \draw[popline] (\x,0) -- (\x,-7.15);
\end{tikzpicture}
\legende{Tableau comparatif des trois entités politiques étudiées : lamidat peul, chefferie des Grassfields, sultanat bamoun. Chaque ligne correspond à un critère de comparaison.}
\end{popfigure}

\courssub{Les points communs}

La lecture horizontale du tableau fait apparaître quatre traits partagés par les trois entités.

\textbf{Un pouvoir centralisé et sacralisé.} Dans les trois cas, le pouvoir est concentré entre les mains d'un chef unique : le \cle{lamido}, le \cle{fon} ou le \cle{sultan}. Ce pouvoir n'est pas seulement politique : il est \emph{sacralisé}, c'est-à-dire entouré d'une dimension religieuse. Le lamido gouverne au nom de l'islam ; le fon est le garant du lien avec les ancêtres ; le sultan bamoun est à la fois chef de guerre, chef religieux et descendant d'une dynastie fondatrice. Le chef n'est donc pas un simple administrateur : il incarne l'unité du groupe.

\textbf{Une hiérarchie sociale marquée.} Aucune de ces sociétés n'est égalitaire. Autour du chef s'ordonnent des catégories bien définies : princes et notables, dignitaires de cour, guerriers, agriculteurs, artisans, et parfois captifs ou esclaves. Cette hiérarchie est visible dans l'organisation de la cour, dans les titres honorifiques et dans les cérémonies.

\textbf{Le palais comme symbole du pouvoir.} Chez les Peuls comme dans les Grassfields et au pays bamoun, le palais du chef est le centre de la vie politique : on y rend la justice, on y reçoit les visiteurs, on y conserve les insignes du pouvoir. Le palais de Foumban, reconstruit sous le sultan \cle{Njoya}, en est l'illustration la plus célèbre.

\textbf{Une économie agricole et commerciale.} Les trois entités reposent sur l'agriculture (mil et sorgho au Nord, maïs et tubercules dans les Grassfields et chez les Bamouns), complétée par l'élevage au Nord, l'artisanat (tissage, poterie, forge) et le commerce, qu'il s'agisse des caravanes transsahariennes, des marchés périodiques des Grassfields ou des réseaux commerciaux bamouns.

\courssub{Les différences}

La lecture verticale du tableau révèle trois profils distincts.

\textbf{Le lamidat : une entité dépendante du califat de Sokoto.} Les lamidats peuls du Nord-Cameroun (Rey-Bouba, Banyo, Tibati, Ngaoundéré, Maroua) sont nés du jihad de \cle{Modibo Adama} au début du XIX\textsuperscript{e} siècle. Leur originalité est d'être \emph{subordonnés} au califat de \cle{Sokoto} (actuel Nigeria) : le lamido reçoit son investiture du calife et lui doit allégeance. L'islam y est religion d'État dès l'origine, et le droit musulman structure la justice.

\textbf{La chefferie des Grassfields : petite taille et sociétés secrètes.} À l'ouest, le plateau des Grassfields est morcelé en une multitude de chefferies de \emph{petite taille} (quelques villages à quelques milliers d'habitants). Le fon y est puissant mais \emph{encadré} : les \cle{sociétés secrètes} ou sociétés régulatrices (comme le \cle{Kwifon} ou la \cle{Takembeng}) contrôlent ses décisions, veillent au respect des coutumes et peuvent le sanctionner. C'est un pouvoir tempéré par des contre-pouvoirs traditionnels, et l'islam y est absent ou marginal.

\textbf{Le sultanat bamoun : un royaume conquérant et innovant.} Le royaume bamoun, fondé par \cle{Ncharé Yen}, s'est construit par la conquête militaire et l'assimilation des peuples vaincus. Sous le sultan Njoya (fin XIX\textsuperscript{e} -- début XX\textsuperscript{e} siècle), il se distingue par une remarquable capacité d'\emph{innovation} : invention d'une écriture originale (l'\cle{A-ka-u-ku}), rédaction d'une histoire du royaume, réformes administratives, adoption volontaire de l'islam. C'est un État centralisé, expansionniste et ouvert aux emprunts.

\begin{popfigure}
\begin{tikzpicture}[scale=0.92]
\draw[thick,poporange,fill=poporangeL] (0,0) circle (2.3);
\draw[thick,popgreen,fill=popgreenL] (3.2,0) circle (2.3);
\draw[thick,poppurple,fill=poppurpleL] (1.6,2.6) circle (2.3);
\node[font=\bfseries\small,text=poporange] at (-1.5,1.9) {Lamidat};
\node[font=\bfseries\small,text=popgreen] at (4.7,1.9) {Chefferie};
\node[font=\bfseries\small,text=poppurple] at (1.6,5.1) {Sultanat bamoun};
\node[font=\scriptsize,text width=2.6cm,align=center] at (-1.3,-0.6) {Dépend de Sokoto\\ islam d'État\\ jihad};
\node[font=\scriptsize,text width=2.6cm,align=center] at (4.5,-0.6) {Petite taille\\ sociétés secrètes\\ culte des ancêtres};
\node[font=\scriptsize,text width=2.8cm,align=center] at (1.6,3.4) {Conquérant\\ écriture inventée\\ innovations};
\node[font=\scriptsize,text width=2.2cm,align=center] at (1.6,0.9) {Chef sacralisé\\ hiérarchie\\ palais\\ agriculture};
\node[font=\scriptsize,text width=1.5cm,align=center] at (0.35,1.75) {Islam présent};
\node[font=\scriptsize,text width=1.7cm,align=center] at (2.85,1.75) {Palais central};
\end{tikzpicture}
\legende{Diagramme de Venn des trois entités : au centre, les traits communs (pouvoir centralisé et sacralisé, hiérarchie sociale, palais, économie agricole et commerciale) ; dans les parties exclusives, les spécificités de chaque modèle.}
\end{popfigure}

\begin{exoresolu}[Comparer pour conclure]
\textbf{Énoncé.} À partir du tableau comparatif, dégagez deux points communs et deux différences entre le lamidat peul et la chefferie des Grassfields, puis formulez une phrase de conclusion.
\tcblower
\textbf{Solution détaillée.}
\begin{itemize}
\item \emph{Points communs} : dans les deux cas, le pouvoir est détenu par un chef sacralisé (lamido / fon) dont le palais est le centre politique ; dans les deux cas, la société est hiérarchisée et l'économie repose sur l'agriculture et le commerce.
\item \emph{Différences} : le lamidat est issu du jihad peul et dépend du califat de Sokoto, tandis que la chefferie des Grassfields est une entité ancienne, indépendante et de petite taille ; l'islam est religion d'État dans le lamidat, alors que la chefferie pratique les religions traditionnelles et que le fon est contrôlé par les sociétés secrètes.
\item \emph{Conclusion} : « Ainsi, lamidats et chefferies partagent une même logique de pouvoir centralisé et sacralisé, mais ils diffèrent profondément par leur origine, leur religion et leurs institutions, ce qui prouve la diversité des organisations politiques du Cameroun précolonial. »
\end{itemize}
\end{exoresolu}

\courssub{Conclusion intermédiaire : diversité et richesse politique avant 1884}

\begin{methode}[Rédiger une conclusion justifiée]
Une conclusion d'étude comparée se construit en trois temps :
\begin{enumerate}
\item \textbf{Rappeler la problématique} : de quoi parlait-on ? (Ici : le Cameroun précolonial était-il politiquement organisé ?)
\item \textbf{Présenter les faits établis} : résumer sans détail les résultats de la comparaison (points communs, différences).
\item \textbf{Formuler un jugement argumenté} : tirer une leçon générale et, si possible, ouvrir sur la suite (la colonisation).
\end{enumerate}
\end{methode}

Appliquons cette méthode. À la veille de la colonisation, le territoire camerounais n'était ni un vide politique ni un ensemble uniforme. Lamidats peuls, chefferies des Grassfields et sultanat bamoun partagent des traits communs — pouvoir centralisé et sacralisé, hiérarchie sociale, palais, économie agricole et commerciale — qui prouvent l'existence d'États et de sociétés structurés. Mais ils diffèrent par leur origine, leur taille, leur religion et leurs institutions : entité subordonnée à Sokoto, mosaïque de petites chefferies encadrées par des sociétés secrètes, royaume conquérant et innovant. Cette \cle{diversité} est une richesse : elle montre que les populations camerounaises avaient inventé, bien avant 1884, plusieurs réponses originales au problème de l'organisation du pouvoir.

\begin{savaistu}[Des chefs précoloniaux aux « auxiliaires » de la colonisation]
Après 1884, les colonisateurs allemands, puis français et britanniques, ne détruisent pas toutes les entités précoloniales : ils s'appuient sur elles. C'est la politique d'\cle{administration indirecte} : les lamidos du Nord et les fons des Grassfields sont maintenus en place mais transformés en \emph{chefs auxiliaires} de l'administration coloniale, chargés de collecter les impôts, de recruter les travailleurs forcés et de transmettre les ordres. Certaines entités sont ainsi utilisées, d'autres affaiblies, d'autres détruites : le sultan Njoya, jugé trop indépendant, sera destitué par les Français en 1931 et exilé à Yaoundé, où il meurt en 1933. La colonisation bouleverse donc durablement ces organisations politiques héritées du passé.
\end{savaistu}

\courssub{TD : Organisation d'une chefferie de la localité de l'établissement}

Conformément au programme, ce travail dirigé vise à appliquer les notions étudiées à une situation concrète : la chefferie la plus proche de votre établissement.

\begin{methode}[Démarche du travail de terrain]
\begin{enumerate}
\item \textbf{Identification} : nom de la chefferie, groupement ethnique, localisation (arrondissement, département, région), nom du chef actuel et titre (fon, lamido, sultan, chef de groupement, etc.).
\item \textbf{Origine et histoire} : date et circonstances de la fondation (migration, conquête, scission), dynastie régnante, événements marquants (guerres, alliances, colonisation).
\item \textbf{Institutions politiques} : rôle du chef, conseil de notables, sociétés secrètes ou régulatrices (Kwifon, Takembeng, Ngiri, etc.), administration des quartiers/villages, justice coutumière.
\item \textbf{Organisation sociale} : hiérarchie (princes, notables, lignages, catégories d'âge, éventuels captifs/descendants d'esclaves), place des femmes (reine-mère, associations féminines).
\item \textbf{Économie et vie culturelle} : activités principales (agriculture, élevage, artisanat, commerce), marchés, fêtes annuelles (Ngoun, Medumba, Nyem-Nyem, etc.), patrimoine matériel et immatériel (palais, objets sacrés, langue, écriture).
\item \textbf{Rapport aux pouvoirs extérieurs} : relations avec l'administration coloniale (chefferie d'institution, canton), avec l'État actuel (chefferie de 1er, 2e, 3e degré), avec les confessions religieuses.
\end{enumerate}
\end{methode}

\begin{experience}[Enquête de terrain : fiche de restitution]
\textbf{Matériel} : questionnaire papier ou numérique, appareil photo (si autorisé), carnet de notes.
\textbf{Protocole} :
\begin{enumerate}
\item Prendre rendez-vous avec le secrétariat de la chefferie (via l'administration de l'établissement).
\item Effectuer la visite en petit groupe (3--4 élèves) avec un enseignant accompagnateur.
\item Recueillir les informations auprès du chef, d'un notable ou du secrétaire de la chefferie.
\item Photographier (si autorisé) le palais, les insignes, le musée éventuel, les lieux de culte.
\item Rédiger collectivement un rapport structuré (2--3 pages) illustré, à remettre sous quinzaine.
\end{enumerate}
\textbf{Observations attendues} : identifier les continuités et ruptures avec le modèle théorique étudié (lamidat, chefferie Grassfields, sultanat). Noter l'adaptation aux contextes moderne et étatique.
\textbf{Interprétation} : en quoi cette chefferie locale illustre-t-elle la diversité des organisations politiques camerounaises ? Comment le pouvoir traditionnel s'articule-t-il aujourd'hui avec l'administration décentralisée ?
\end{experience}

\begin{exercice}[À vous de comparer]
En vous aidant du tableau comparatif, rédigez un paragraphe de dix à quinze lignes répondant à la question : « Le Cameroun précolonial était-il un espace politiquement organisé ? » Votre paragraphe devra contenir au moins deux points communs, deux différences et une conclusion justifiée.
\end{exercice}

\begin{corrige}[Exercice : À vous de comparer]
Éléments de correction attendus :
\begin{itemize}
\item \emph{Introduction de la réponse} : reformulation de la problématique (existence d'organisations politiques avant 1884).
\item \emph{Points communs} : pouvoir centralisé entre les mains d'un chef sacralisé (lamido, fon, sultan) ; hiérarchie sociale marquée ; rôle central du palais ; économie fondée sur l'agriculture et le commerce.
\item \emph{Différences} : le lamidat est issu du jihad et dépend du califat de Sokoto ; la chefferie des Grassfields est de petite taille et son fon est contrôlé par les sociétés secrètes ; le sultanat bamoun est un royaume conquérant qui a inventé une écriture.
\item \emph{Conclusion justifiée} : le Cameroun précolonial était bien politiquement organisé, mais selon des modèles divers ; cette diversité sera ensuite utilisée, affaiblie ou détruite par l'administration coloniale.
\end{itemize}
Le paragraphe doit être rédigé en phrases liées, sans simple énumération, et chaque affirmation doit être illustrée par un exemple précis (Rey-Bouba, Kwifon, Njoya...).
\end{corrige}

\coursec{TD : l'organisation d'une chefferie de la localité de l'établissement}

Dans la section précédente, nous avons comparé trois modèles d'organisation politique du Cameroun précolonial : le \cle{lamidat} peul du Nord, la \cle{chefferie} des Grassfields et le \cle{sultanat bamoun}. Ces modèles ne sont pas des souvenirs figés dans les manuels : autour de votre établissement scolaire existe très probablement une chefferie vivante, avec son chef, ses notables, son palais et ses cérémonies. Ce travail dirigé (TD) vous invite à devenir, le temps d'une enquête, de véritables historiens de terrain : observer, interroger, croiser les sources et restituer vos découvertes.

\courssub{Objectifs du TD}

Ce TD poursuit deux objectifs complémentaires, directement issus du programme officiel.

\begin{itemize}
\item \textbf{Appliquer les notions de l'unité à une situation concrète de la vie courante} : les concepts étudiés en classe (chef, succession héréditaire, conseil des notables, sociétés secrètes ou régulatrices, palais) doivent vous servir de grille de lecture pour comprendre une chefferie réelle et proche de vous.
\item \textbf{Rédiger et justifier une démarche, une réponse ou une conclusion} : vous devrez produire un croquis légendé accompagné d'un texte qui explique comment vous avez procédé, quelles sources vous avez utilisées et pourquoi vos conclusions sont fiables.
\end{itemize}

\begin{definition}[Enquête historique locale]
Une \cle{enquête historique locale} est une démarche d'investigation qui consiste à recueillir des informations sur le passé et le présent d'une communauté proche (village, quartier, chefferie) à partir de \cle{sources orales} (témoignages, récits des anciens), de \cle{sources écrites} (documents administratifs, archives) et de \cle{sources matérielles} (monuments, palais, objets, photographies). Comme toute démarche historique, elle exige de croiser les sources pour vérifier les informations.
\end{definition}

\courssub{La démarche d'enquête en cinq étapes}

L'enquête se déroule en cinq étapes ordonnées : \textbf{identifier} la chefferie, \textbf{interroger} les informateurs, \textbf{observer} les lieux, \textbf{croquer} (réaliser un croquis), puis \textbf{rédiger} la restitution. Chaque étape prépare la suivante.

\begin{popfigure}
\begin{center}
\begin{tikzpicture}[font=\small]
\tikzset{etape/.style={rectangle, rounded corners=3pt, draw=popblue, fill=popblueL, thick, minimum width=2.4cm, minimum height=1.1cm, align=center}}
\node[etape, align=center] (e1) at (0,0){\textbf{1. Identifier}\\la chefferie};
\node[etape, align=center] (e2) at (3.4,0){\textbf{2. Interroger}\\chef, notables, anciens};
\node[etape, align=center] (e3) at (6.8,0){\textbf{3. Observer}\\palais, cérémonies};
\node[etape, align=center] (e4) at (10.2,0){\textbf{4. Croquer}\\plan légendé};
\node[etape, align=center] (e5) at (13.6,0){\textbf{5. Rédiger}\\texte justificatif};
\draw[-{Stealth[length=3mm]}, thick, poporange] (e1) -- (e2);
\draw[-{Stealth[length=3mm]}, thick, poporange] (e2) -- (e3);
\draw[-{Stealth[length=3mm]}, thick, poporange] (e3) -- (e4);
\draw[-{Stealth[length=3mm]}, thick, poporange] (e4) -- (e5);
\node[below=0.15cm of e1, text=popmuted, align=center, font=\scriptsize] {nom, localisation,\\chef actuel};
\node[below=0.15cm of e2, text=popmuted, align=center, font=\scriptsize] {questionnaire\\préparé};
\node[below=0.15cm of e3, text=popmuted, align=center, font=\scriptsize] {photos autorisées,\\notes de terrain};
\node[below=0.15cm of e4, text=popmuted, align=center, font=\scriptsize] {fond simplifié,\\légende organisée};
\node[below=0.15cm of e5, text=popmuted, align=center, font=\scriptsize] {sources citées,\\conclusion};
\end{tikzpicture}
\end{center}
\legende{Schéma de la démarche d'enquête en cinq étapes : identifier, interroger, observer, croquer, rédiger.}
\end{popfigure}

\textbf{Étape 1 : identifier la chefferie.} Avec votre professeur, repérez la chefferie la plus proche de l'établissement : son nom, son village ou son quartier, l'identité du chef actuel, sa place dans la hiérarchie administrative (chef de 1\iere{}, 2\ieme{} ou 3\ieme{} catégorie).

\textbf{Étape 2 : interroger.} Préparez un questionnaire et menez l'entretien auprès du chef, d'un notable ou d'un ancien du village.

\textbf{Étape 3 : observer.} Visitez le palais ou la place de la chefferie : notez la disposition des cases, la présence d'une place publique, d'un lieu de culte, d'arbres symboliques, des éléments d'architecture.

\textbf{Étape 4 : croquer.} Réalisez un croquis légendé du palais ou de l'organisation de la chefferie.

\textbf{Étape 5 : rédiger.} Rédigez un texte d'une quinzaine de lignes présentant vos résultats et justifiant votre démarche.

\begin{methode}[Mener une enquête historique locale]
\begin{enumerate}
\item \textbf{Préparer un questionnaire} avant la visite : questions simples, ouvertes, classées par thème (fondation, succession, institutions, activités).
\item \textbf{Prendre des notes} pendant l'entretien : nom et fonction de l'informateur, date et lieu de l'entretien.
\item \textbf{Croiser au moins deux sources} : un récit oral doit être confronté à un autre témoignage, à un document administratif ou à une observation matérielle. Une information donnée par une seule personne reste une hypothèse.
\item \textbf{Citer ses informateurs} dans la production finale : « selon M. X, notable de la chefferie de Y, entretien du... ». Citer ses sources est la marque de l'honnêteté scientifique.
\item \textbf{Distinguer les faits des opinions} : rapporter ce qui est attesté, et signaler ce qui relève de la tradition ou de la mémoire familiale.
\end{enumerate}
\end{methode}

\courssub{Les sources mobilisables}

L'historien de terrain dispose de plusieurs catégories de sources, qu'il doit combiner.

\begin{center}
\begin{tabularx}{\linewidth}{|l|X|X|}
\hline
\rowcolor{popblueL} \textbf{Type de source} & \textbf{Exemples} & \textbf{Ce qu'elle apporte} \\
\hline
Sources orales & Interview du chef ou d'un notable, récits des anciens, chants et traditions & Histoire de la fondation, généalogie des chefs, sens des cérémonies \\
\hline
Sources écrites & Arrêtés de reconnaissance des chefs traditionnels, registres administratifs, monographies de circonscription & Dates officielles, catégorie du chef, limites de la chefferie \\
\hline
Sources matérielles & Palais, cases à tourelles, statues, masques, objets royaux, photographies & Organisation de l'espace, architecture, symboles du pouvoir \\
\hline
\end{tabularx}
\end{center}

\begin{attention}[Respecter les us et coutumes]
La visite d'une chefferie n'est pas une visite ordinaire : c'est un lieu de pouvoir et de sacré. Quelques règles à respecter impérativement :
\begin{itemize}
\item se faire annoncer et \textbf{saluer le chef selon le protocole local} (salutations, gestes, tenue correcte) ;
\item demander l'\textbf{autorisation avant toute photographie}, en particulier à l'intérieur du palais ou lors d'une cérémonie ;
\item ne pas toucher les objets rituels (masques, statues, trônes) sans permission ;
\item remercier les informateurs à la fin de l'entretien.
\end{itemize}
Un manquement à ces règles peut compromettre l'enquête et offense la communauté qui vous accueille.
\end{attention}

\courssub{Les questions guides de l'enquête}

Pour structurer votre questionnaire, organisez vos questions autour de quatre thèmes, qui reprennent les notions de la leçon.

\begin{itemize}
\item \textbf{Qui dirige ?} Quel est le nom du chef actuel ? Depuis quand règne-t-il ? Quel est son titre exact ?
\item \textbf{Comment le chef est-il désigné ?} La succession est-elle \cle{héréditaire} ? Dans quelle famille ou quel lignage le chef est-il choisi ? Qui participe au choix (notables, sociétés) ? Y a-t-il des rites d'intronisation ?
\item \textbf{Quelles sont les institutions ?} Existe-t-il un \cle{conseil des notables} ? Des sociétés régulatrices ou des groupes chargés de la sécurité, de la justice ou des cérémonies ? Quels sont leurs rôles ?
\item \textbf{Quelles activités économiques et culturelles ?} Quelles sont les activités de la communauté (agriculture, artisanat, commerce) ? Quelles cérémonies, danses ou fêtes rythment la vie de la chefferie ?
\end{itemize}

\courssub{Réaliser le croquis : gabarit et méthode}

La production attendue comporte un \cle{croquis} légendé du palais ou de l'organisation de la chefferie. Le croquis n'est pas un dessin artistique : c'est un document géographique et historique simplifié, qui ne retient que l'essentiel.

\begin{methode}[Réaliser un croquis géographique et historique]
\begin{enumerate}
\item Tracer un \textbf{fond simplifié} : contour général du site, quelques repères (route, place, cases principales), sans chercher la précision cartographique.
\item \textbf{Sélectionner les éléments} utiles au sujet : case du chef, case des notables, place des cérémonies, arbre à palabres, entrée du palais. On ne reporte pas tout : on choisit.
\item Construire une \textbf{légende organisée} : symboles numérotés ou fléchés, classés (pouvoir, habitat, cérémonies).
\item Ajouter un \textbf{titre}, l'orientation (flèche du nord) et, si possible, une échelle indicative.
\end{enumerate}
\end{methode}

\begin{popfigure}
\begin{center}
\begin{tikzpicture}[font=\scriptsize, scale=0.95]
% Contour du site
\draw[thick, popdark, rounded corners=8pt] (0,0) rectangle (9,6);
% Route
\draw[thick, popmuted, dashed] (-0.4,1) -- (9.4,1.6) node[pos=0.9, above, sloped, text=popmuted]{route};
% Case du chef
\draw[fill=poporangeL, draw=poporange, thick] (4,3.4) circle (0.55);
\node[text=popdark] at (4,3.4) {\textbf{1}};
% Cases des notables
\draw[fill=popblueL, draw=popblue, thick] (2.2,4.4) circle (0.35);
\node[text=popdark] at (2.2,4.4) {\textbf{2}};
\draw[fill=popblueL, draw=popblue, thick] (5.9,4.5) circle (0.35);
\node[text=popdark] at (5.9,4.5) {\textbf{2}};
% Place des cérémonies
\draw[fill=popgreenL, draw=popgreen, thick] (4,1.9) ellipse (1.6 and 0.7);
\node[text=popdark] at (4,1.9) {\textbf{3}};
% Arbre à palabres
\node[popgreen] at (7.3,3.2) {\textbf{4}};
\draw[popgreen, thick] (7.3,2.6) -- (7.3,3.4);
\draw[fill=popgreenL, draw=popgreen, thick] (7.3,3.7) circle (0.45);
% Entrée
\draw[thick, poppink] (4,-0.02) -- (4,0.6);
\node[text=poppink, below] at (4,0) {\textbf{5}};
% Nord
\draw[-{Stealth[length=3mm]}, thick, popdark] (8.4,4.6) -- (8.4,5.6);
\node[above, text=popdark] at (8.4,5.6) {N};
% Légende
\node[anchor=west, text=popdark] at (9.6,5.2) {\textbf{Légende}};
\node[anchor=west] at (9.6,4.6) {\textbf{1} Case du chef (à compléter)};
\node[anchor=west] at (9.6,4.0) {\textbf{2} Cases des notables};
\node[anchor=west] at (9.6,3.4) {\textbf{3} Place des cérémonies};
\node[anchor=west] at (9.6,2.8) {\textbf{4} Arbre à palabres};
\node[anchor=west] at (9.6,2.2) {\textbf{5} Entrée principale};
\node[anchor=west, text=popmuted] at (9.6,1.4) {Échelle indicative :};
\draw[thick, popmuted] (9.6,1.0) -- (11.1,1.0);
\node[anchor=west, text=popmuted] at (11.2,1.0) {environ 20 m};
\end{tikzpicture}
\end{center}
\legende{Gabarit de croquis d'une chefferie locale, à adapter et compléter par les élèves lors de l'enquête : fond simplifié, éléments numérotés, légende organisée, titre et nord.}
\end{popfigure}

\begin{exemplebox}[Deux terrains d'enquête possibles]
\begin{itemize}
\item \textbf{À Foumban} (région de l'Ouest) : l'enquête portera sur le \cle{palais royal} des sultans bamouns et le musée qui l'accompagne. On y observe l'architecture du palais, les objets royaux, et l'on interroge les gardiens des traditions sur l'histoire de la dynastie fondée par Nchare Yen.
\item \textbf{À Bafoussam} (Grassfields) : l'enquête portera sur la chefferie et son architecture remarquable de \cle{cases à tourelles} (grandes cases coniques au toit de chaume). On y étudie le rôle du \emph{fo} (chef), du conseil des notables et des sociétés régulatrices comme le \emph{Takumbeng}.
\end{itemize}
Dans les deux cas, la démarche reste identique : identifier, interroger, observer, croquer, rédiger.
\end{exemplebox}

\begin{savaistu}[Les chefs traditionnels reconnus par l'État]
Au Cameroun, les chefs traditionnels sont officiellement reconnus par l'État et classés, depuis le décret de 1977, en \cle{chefs de 1\iere{}, 2\ieme{} et 3\ieme{} catégorie}. Ils sont nommés par arrêté administratif après désignation selon la coutume, et ils collaborent avec l'administration (collecte de certaines informations, maintien de l'ordre coutumier, mobilisation des communautés). C'est pourquoi les \textbf{arrêtés de reconnaissance} conservés dans les sous-préfectures sont de précieuses sources écrites pour votre enquête.
\end{savaistu}

\courssub{Restitution en classe et évaluation}

La dernière phase du TD se déroule en classe. Chaque groupe présente oralement son croquis en cinq minutes, puis la classe \textbf{confronte les résultats aux modèles étudiés} : la chefferie enquêtée ressemble-t-elle au modèle des Grassfields (chef sacré, conseil des notables, sociétés régulatrices) ? Au modèle du lamidat (pouvoir centralisé autour du lamido) ? Au modèle bamoun (cour organisée, artisans attachés au palais) ? Cette confrontation permet de vérifier que les notions de la leçon éclairent réellement la réalité locale.

\begin{exoresolu}[Exemple de restitution rédigée (modèle attendu)]
Après avoir enquêté dans la chefferie de votre localité, rédigez un court texte justifiant votre démarche et présentant une conclusion.
\tcblower
\textbf{Solution détaillée (modèle de rédaction).} « Notre groupe a enquêté sur la chefferie de [nom], située à [lieu], le [date]. Nous avons d'abord identifié la chefferie grâce au document administratif de la sous-préfecture : le chef actuel, [nom], est un chef de [catégorie], reconnu par arrêté. Nous avons ensuite interrogé [nom, fonction], qui nous a expliqué que la succession est héréditaire dans le lignage de [nom], le choix du successeur revenant aux notables réunis en conseil. Ces informations orales ont été croisées avec le récit d'un ancien du village, qui les a confirmées, et avec notre observation du palais : la case du chef, au centre, est entourée des cases des notables, ce qui traduit dans l'espace l'organisation du pouvoir. Nous en concluons que cette chefferie correspond au modèle des Grassfields étudié en classe : un chef sacré, un conseil de notables et des sociétés qui encadrent la vie sociale. Sources : entretien avec [informateur], arrêté de reconnaissance, photographies autorisées du palais. »
\end{exoresolu}

\begin{exercice}[À vous de jouer]
Par groupes de quatre, menez l'enquête dans la chefferie de la localité de votre établissement.
\begin{enumerate}
\item Rédigez un questionnaire de huit questions portant sur la désignation du chef, les institutions et les activités de la chefferie.
\item Réalisez un croquis légendé du palais ou de l'organisation de la chefferie (titre, nord, légende organisée).
\item Rédigez un texte de quinze lignes justifiant votre démarche et comparant la chefferie étudiée à l'un des trois modèles de la leçon.
\end{enumerate}
\end{exercice}

\begin{corrige}[Exercice : critères d'évaluation]
L'évaluation porte sur trois critères, annoncés à l'avance :
\begin{enumerate}
\item \textbf{Exactitude des informations} (environ la moitié des points) : les faits rapportés sont-ils exacts, datés, croisés entre au moins deux sources et attribués à des informateurs cités ?
\item \textbf{Qualité du croquis} : fond simplifié, éléments sélectionnés et pertinents, légende organisée, titre et orientation présents.
\item \textbf{Justification de la démarche} : le texte explique-t-il clairement les étapes suivies, les sources mobilisées et la conclusion comparée aux modèles de la leçon (lamidat, Grassfields, Bamoun) ?
\end{enumerate}
Un travail sans sources citées, ou un croquis sans légende ni titre, ne peut obtenir la totalité des points, même si le contenu est riche.
\end{corrige}

\begin{aretenir}[L'essentiel du TD]
\begin{itemize}
\item L'enquête historique locale applique les notions de la leçon à une chefferie réelle : elle suit cinq étapes — identifier, interroger, observer, croquer, rédiger.
\item On croise toujours au moins deux sources (orales, écrites, matérielles) et l'on cite ses informateurs.
\item Le croquis est un document simplifié : fond, éléments sélectionnés, légende organisée, titre et nord.
\item Le respect des us et coutumes (salutations, autorisation de photographier) est une condition de réussite de l'enquête.
\item La restitution confronte la chefferie étudiée aux trois modèles d'organisation politique du Cameroun précolonial.
\end{itemize}
\end{aretenir}

\newpage

\coursec{Activité d'intégration}

\courssub{Objectifs de l'activité}

Cette activité d'intégration clôt la leçon sur le Cameroun à la veille de la colonisation. Elle vise à faire mobiliser par l'élève, dans une situation concrète, l'ensemble des savoirs étudiés : le lamidat peul du Nord-Cameroun, la chefferie des Grassfields et le sultanat bamoun.

À l'issue de cette activité, l'élève doit être capable de :
\begin{itemize}
\item \cle{localiser} sur une carte muette les trois entités politiques étudiées ;
\item \cle{comparer} les trois modèles d'organisation politique à l'aide d'un tableau ;
\item \cle{réaliser} un croquis légendé de la chefferie de sa localité (ou de la chefferie la plus proche de l'établissement) ;
\item \cle{rédiger} un paragraphe argumenté répondant à une question de synthèse, en justifiant sa réponse par trois arguments tirés des documents ;
\item \cle{présenter} et défendre sa production devant la classe.
\end{itemize}

Ces tâches correspondent aux savoir-faire du programme officiel : appliquer les notions de l'unité à des situations concrètes de la vie courante, et rédiger et justifier une démarche, une réponse ou une conclusion.

\courssub{Situation}

\textbf{Contexte.} Dans le cadre de la semaine culturelle de ton lycée technique, le club d'Histoire organise une exposition intitulée \emph{« Le Cameroun avant les Blancs : des États organisés »}. Le proviseur a confié à ta classe de Première technique la réalisation d'un panneau d'exposition. Un visiteur européen invité à la cérémonie d'ouverture a déclaré, lors d'une visite préparatoire : \emph{« Avant la colonisation, il n'y avait ici ni États, ni organisation politique, seulement des villages dispersés. »} Choqués par cette affirmation, les élèves décident de démontrer le contraire à l'aide de documents historiques et d'une enquête sur la chefferie de leur localité.

\textbf{Le dossier documentaire fourni à chaque groupe :}

\begin{itemize}
\item \textbf{Document 1 :} carte muette du Cameroun avec les grands ensembles géographiques (Nord soudanien, Ouest montagneux des Grassfields, plateau de l'Adamaoua, Sud forestier).
\item \textbf{Document 2 :} photographie du palais royal de Foumban, construit par le sultan \cle{Ibrahim Njoya} (vers 1917), avec la légende suivante : « Le palais compte plusieurs bâtiments, une cour d'honneur et abrite la cour, les serviteurs et les ateliers des artisans royaux. »
\item \textbf{Document 3 :} photographie d'une chefferie des Grassfields (case du \emph{fondom}, place publique, cases des notables), montrant l'organisation spatiale autour du palais du \cle{fon}.
\item \textbf{Document 4 :} extrait de récit sur le roi Njoya : « Njoya, sultan des Bamouns de 1886 à 1931, inventa vers 1895 une écriture, le \emph{shümom}, pour son peuple. Il fit construire des écoles, rédigea une histoire de son royaume et fit dessiner une carte de son territoire. Son État comptait une armée, des ministres et des collecteurs d'impôts. »
\item \textbf{Document 5 :} données chiffrées sur le lamidat de \cle{Rey-Bouba} (Nord-Cameroun) : fondé au début du XIX\textsuperscript{e} siècle après le jihad peul de Modibo Adama (1804-1809), le lamidat s'étendait sur environ \SI{30000}{km^{2}} et comptait, vers 1900, près de \num{100000} habitants. Le \cle{lamido} commandait une cavalerie, prélevait l'impôt (le \emph{zakat}) et rendait la justice dans son palais entouré d'une enceinte.
\item \textbf{Document 6 :} informations recueillies par enquête auprès du chef et des notables de la localité : nom de la chefferie, nom du chef, nombre de quartiers, existence d'une société secrète ou d'un conseil des notables, fêtes traditionnelles, marché hebdomadaire.
\end{itemize}

\textbf{Problème posé :} le visiteur européen a-t-il raison ? Le Cameroun précolonial était-il un espace politiquement organisé ? Ta production devra le démontrer.

\courssub{Consignes}

Par groupes de quatre à six élèves, vous réaliserez les tâches suivantes, dans l'ordre :

\begin{enumerate}
\item \textbf{Localisation.} Sur la carte muette (document 1), placez et nommez les trois entités politiques étudiées : le lamidat de Rey-Bouba, le sultanat bamoun (capitale Foumban) et une chefferie des Grassfields de votre choix (par exemple Bafou, Bafut ou Bali). Ajoutez une flèche du nord et une légende.
\item \textbf{Comparaison.} Complétez le tableau comparatif fourni (ou reproduisez-le) en remplissant, pour chaque entité : le chef (titre), la capitale ou le centre du pouvoir, les institutions (armée, justice, impôt, conseil), et une originalité propre à l'entité.
\item \textbf{Enquête et croquis.} À partir des informations du document 6 (enquête locale), réalisez le \cle{croquis} légendé de la chefferie de votre localité (ou de la chefferie la plus proche de l'établissement). Le croquis doit comporter : un titre, une orientation (flèche du nord), une échelle indicative, et au moins cinq éléments légendés (palais ou case du chef, place publique, cases des notables, marché, quartiers, lieux de culte, etc.).
\item \textbf{Rédaction.} Rédigez un paragraphe de conclusion de 10 à 15 lignes répondant à la question : \emph{« Le Cameroun précolonial était-il un espace organisé ? Justifiez votre réponse par trois arguments tirés des documents. »} Chaque argument doit être appuyé sur un document précis (citez le numéro du document).
\item \textbf{Présentation.} Chaque groupe présente son panneau en 5 minutes. La production est évaluée selon la grille suivante : exactitude historique (8 points), qualité du croquis et de la légende (6 points), argumentation et justification de la conclusion (6 points).
\end{enumerate}

\begin{attention}[Pièges à éviter]
Ne confonds pas les titres des chefs : \cle{lamido} pour le lamidat peul, \cle{fon} (ou \emph{fo}) pour la chefferie des Grassfields, \cle{sultan} (ou \emph{mfon}) pour le royaume bamoun. Dans le croquis, n'oublie jamais le titre, l'orientation et la légende : un croquis sans légende est une carte muette, donc un travail incomplet. Dans la rédaction, une affirmation sans référence à un document n'est pas un argument.
\end{attention}

\courssub{Ressources à mobiliser}

Pour réussir cette activité, tu dois mobiliser les savoirs de la leçon :

\begin{aretenir}[Les trois modèles d'organisation politique]
\begin{itemize}
\item \textbf{Le lamidat} (Nord-Cameroun) : État centralisé né du jihad peul du début du XIX\textsuperscript{e} siècle. Le \cle{lamido} concentre les pouvoirs politique, militaire, judiciaire et religieux. Il prélève l'impôt, commande une cavalerie et s'appuie sur des subalternes à la tête des villages.
\item \textbf{La chefferie des Grassfields} (Ouest) : entité politique de taille modeste dirigée par le \cle{fon}, assisté d'un conseil de notables et encadré par des sociétés secrètes (comme la société des \emph{ngumba} ou équivalents) qui régulent le pouvoir. Le palais (\emph{fondom}) est le centre politique, économique et rituel.
\item \textbf{Le sultanat bamoun} (Ouest, capitale Foumban) : royaume centralisé fondé par Nchare Yen, dirigé au XIX\textsuperscript{e}-XX\textsuperscript{e} siècle par le sultan \cle{Ibrahim Njoya}, inventeur de l'écriture \emph{shümom}. L'État dispose d'une armée, de ministres, d'un impôt et d'une administration.
\end{itemize}
\end{aretenir}

\begin{methode}[Réaliser un croquis géographique légendé]
\begin{enumerate}
\item Donner un \textbf{titre} précis au croquis (nom de la chefferie, localisation).
\item Tracer le \textbf{plan} simplifié : les éléments sont représentés par des formes simples (rectangles pour les cases, points pour les lieux remarquables, traits pour les voies).
\item Indiquer l'\textbf{orientation} par une flèche du nord et une \textbf{échelle} indicative.
\item Numéroter les éléments et reporter les numéros dans une \textbf{légende} placée sous le croquis.
\item Vérifier que chaque élément de la légende figure bien sur le croquis, et inversement.
\end{enumerate}
\end{methode}

\courssub{Démarche et corrigé commenté}

\begin{exoresolu}[Corrigé commenté de l'activité d'intégration]
\textbf{Énoncé rappelé :} localiser les trois entités, compléter le tableau comparatif, réaliser le croquis de la chefferie locale et rédiger un paragraphe de conclusion argumenté.
\tcblower
\textbf{Étape 1 : la localisation sur la carte muette.}

Le lamidat de Rey-Bouba se place dans le \textbf{Nord-Cameroun} (région du Nord actuelle, plaine de la Bénoué) : c'est un héritage du jihad peul de Modibo Adama (1804-1809), qui a islamisé et organisé politiquement le Nord. Le sultanat bamoun se place sur les \textbf{hautes terres de l'Ouest}, avec Foumban pour capitale. Les chefferies des Grassfields (Bafou, Bafut, Bali, etc.) se placent dans les régions actuelles du \textbf{Nord-Ouest et de l'Ouest}, zone de collines et de plateaux herbeux — d'où le nom anglais de \emph{Grassfields}, « champs d'herbe ». La carte complétée montre donc un territoire couvert d'entités politiques identifiées, ce qui réfute déjà l'idée de « villages dispersés ».

\textbf{Étape 2 : le tableau comparatif.}

\begin{center}
\begin{tabularx}{\linewidth}{|l|X|X|X|}\hline
\rowcolor{popblueL} \textbf{Critère} & \textbf{Lamidat (Rey-Bouba)} & \textbf{Chefferie des Grassfields} & \textbf{Sultanat bamoun} \\ \hline
Chef & Lamido & Fon (ou fo) & Sultan (mfon) \\ \hline
Centre du pouvoir & Palais du lamido, ville de Rey-Bouba & Palais (fondom), place publique & Palais royal de Foumban \\ \hline
Institutions & Cavalerie, impôt (zakat), justice du lamido & Conseil de notables, sociétés secrètes, marché & Armée, ministres, impôts, écoles \\ \hline
Originalité & État issu du jihad peul, islamisé & Pouvoir contrôlé par les notables et les sociétés régulatrices & Écriture shümom inventée par Njoya \\ \hline
\end{tabularx}
\end{center}

\textbf{Étape 3 : le croquis de la chefferie locale.}

Le croquis ci-dessous est un modèle adaptable à la chefferie de la localité de l'établissement (ici, une chefferie type des Grassfields ; chaque groupe le remplace par les données de son enquête).

\begin{popfigure}
\begin{tikzpicture}[scale=1]
\draw[->,thick,popdark] (8.6,3.4) -- (8.6,4.4) node[above]{\textbf{N}};
\draw[fill=popgoldL,draw=popdark] (0,0) rectangle (2.2,1.6);
\node at (1.1,0.8) {\small 1};
\draw[fill=popgreenL,draw=popdark] (3,0.2) rectangle (4.4,1.4);
\node at (3.7,0.8) {\small 2};
\draw[fill=popgreenL,draw=popdark] (5,0.2) rectangle (6.4,1.4);
\node at (5.7,0.8) {\small 3};
\draw[fill=poporangeL,draw=popdark] (0.2,2.6) circle (0.9);
\node at (0.2,2.6) {\small 4};
\draw[fill=poppinkL,draw=popdark] (3,2.4) rectangle (6.4,3.4);
\node at (4.7,2.9) {\small 5};
\draw[thick,popmuted,dashed] (2.2,0.8) -- (3,0.8);
\draw[thick,popmuted,dashed] (4.4,0.8) -- (5,0.8);
\draw[thick,popmuted,dashed] (1.1,1.6) -- (0.6,2.2);
\draw[thick,popmuted,dashed] (1.4,1.6) -- (3.2,2.6);
\end{tikzpicture}
\legende{Croquis type d'une chefferie des Grassfields. 1 : palais du fon (fondom) ; 2 et 3 : cases des notables ; 4 : place publique et marché ; 5 : quartier des familles royales. Les pointillés figurent les voies reliant le palais aux autres lieux. Échelle indicative : le palais mesure environ 40 m de long.}
\end{popfigure}

\textbf{Étape 4 : le paragraphe de conclusion (modèle de rédaction).}

\emph{« Le Cameroun précolonial était un espace politiquement organisé, et non un ensemble de villages dispersés. Trois arguments le démontrent. D'abord, le Nord-Cameroun était structuré en lamidats centralisés : le lamidat de Rey-Bouba s'étendait sur environ \SI{30000}{km^{2}} et regroupait près de \num{100000} habitants dirigés par un lamido qui levait l'impôt, rendait la justice et commandait une cavalerie (document 5). Ensuite, les Grassfields étaient couverts de chefferies hiérarchisées où le fon gouvernait avec un conseil de notables et des sociétés secrètes, autour d'un palais et d'une place publique (document 3). Enfin, le sultanat bamoun était un véritable État administré : le sultan Ibrahim Njoya disposait de ministres, d'une armée et d'écoles, et il a même inventé une écriture, le shümom, pour administrer son royaume (documents 2 et 4). Ainsi, loin d'être désorganisé, le Cameroun précolonial connaissait des États dotés de chefs, d'institutions et de territoires définis. »}

\textbf{Justification de la démarche :} chaque argument cite son document (exigence du savoir-faire « rédiger et justifier »), reprend un savoir précis de la leçon (chiffres, titres des chefs, institutions) et contribue à répondre directement au problème posé.
\end{exoresolu}

\courssub{Bilan de l'activité}

Cette activité d'intégration a permis de mettre en évidence plusieurs acquis essentiels de la leçon :

\begin{itemize}
\item \textbf{Sur le plan des savoirs :} les trois entités étudiées (lamidat, chefferie des Grassfields, sultanat bamoun) constituent trois modèles distincts mais comparables d'organisation politique, chacun avec un chef, un centre du pouvoir et des institutions. Le Cameroun précolonial apparaît ainsi comme un espace politiquement structuré bien avant l'arrivée des Européens en 1884.
\item \textbf{Sur le plan des savoir-faire :} les élèves ont appliqué les notions du cours à une situation concrète de la vie courante (l'exposition du lycée, l'enquête sur la chefferie de leur localité), construit un croquis légendé selon une démarche rigoureuse, et rédigé une conclusion argumentée et justifiée par des documents.
\item \textbf{Sur le plan citoyen et culturel :} l'enquête auprès du chef et des notables a montré que les chefferies actuelles sont les héritières vivantes de ces entités précoloniales : le patrimoine politique traditionnel camerounais n'est pas un passé disparu, mais une réalité que chaque élève peut observer autour de lui.
\end{itemize}

\begin{savaistu}[Le saviez-vous ?]
L'écriture \emph{shümom} inventée par le sultan Ibrahim Njoya vers 1895 comptait à l'origine plus de 400 signes, avant d'être simplifiée à environ 80. Grâce à elle, les Bamouns ont rédigé des chroniques, des cartes et même des ordonnances de pharmacopée traditionnelle. Le palais de Foumban abrite aujourd'hui un musée qui conserve ces manuscrits : c'est l'un des rares cas en Afrique subsaharienne d'une écriture créée localement avant la colonisation.
\end{savaistu}

\newpage

\coursec{Méthodes et exercices résolus}

\coursec{Leçon N04 : Le Cameroun à la veille de la colonisation}

\courssub{Introduction : le Cameroun précolonial, un espace politiquement organisé}

\emph{Contexte.} Avant la signature des traités germano-douala de 1884, le territoire de l'actuel Cameroun n'est pas un vide politique. Il est structuré par une mosaïque d'entités politiques viables, fruit d'une histoire longue de migrations, de conquêtes et de constructions étatiques. Trois grands modèles dominent : les \cle{lamidats} peuls au Nord, les \cle{chefferies des Grassfields} à l'Ouest et au Nord-Ouest, et le \cle{sultanat bamoun} dans la région du Noun.

\begin{definition}[Entité politique précoloniale]
Structure de pouvoir organisée sur un territoire délimité, dotée d'institutions centrales (chef, administration, armée, justice), d'une cohésion sociale et de ressources économiques propres, capable d'assurer la défense et l'ordre intérieur.
\end{definition}

\begin{popfigure}
\begin{tikzpicture}[scale=0.7]
% Contour très simplifié du Cameroun (polygone approximatif)
\draw[thick, popdark, fill=popcream] 
  (0,0) -- (1,0.5) -- (0.5,2) -- (1.5,3) -- (3,3.5) -- (4.5,3) -- (5.5,2.5) -- (6,1.5) -- (5.5,0.5) -- (4,-0.2) -- (2.5,0) -- cycle;
% Zones
% Nord : Lamidats (zone hachurée poporangeL)
\pattern[pattern=north east lines, pattern color=poporange] (0.5,0.2) -- (1.5,1) -- (1,2.2) -- (2.5,2.5) -- (4,2.2) -- (5,1.2) -- (4.5,0.3) -- (2,0) -- cycle;
% Grassfields : Ouest/Nord-Ouest (zone hachurée popblueL)
\pattern[pattern=north west lines, pattern color=popblue] (2.5,2.5) -- (3,3.5) -- (4.5,3) -- (5.5,2.5) -- (5,1.5) -- (4,2.2) -- cycle;
% Bamoun : Noun (zone poppurpleL)
\fill[poppurpleL, opacity=0.6] (3.5,1.8) circle (0.8);
% Villes / Capitales
\node[poporange, font=\small\bfseries] at (1.5,0.5) {Garoua};
\node[poporange, font=\small\bfseries] at (3.5,1.2) {Ngaoundéré};
\node[poporange, font=\small\bfseries] at (5,0.8) {Maroua};
\node[popblue, font=\small\bfseries] at (3.5,3) {Bamenda};
\node[popblue, font=\small\bfseries] at (4.5,2.2) {Bafut};
\node[popblue, font=\small\bfseries] at (5.2,1.8) {Bandjoun};
\node[poppurple, font=\small\bfseries] at (3.5,1.8) {Foumban};
% Flèche Nord
\node[popdark] at (6.5,3.5) {\Large $\uparrow$ N};
\draw[popdark, ->] (6.5,3.2) -- (6.5,2.8);
% Échelle indicative
\draw[popdark, thick] (0.5,-0.5) -- (2.5,-0.5) node[midway, below, font=\small] {$\approx$ 200 km};
\end{tikzpicture}
\legende{Croquis schématique : localisation des trois grands modèles politiques au Cameroun précolonial (Nord : lamidats ; Grassfields : chefferies ; Noun : sultanat bamoun).}
\end{popfigure}

\begin{aretenir}[Idées forces de l'introduction]
\begin{itemize}
\item Le Cameroun précolonial connaît une \textbf{dynamique étatique} marquée au XIX\textsuperscript{e} siècle.
\item Les entités sont \textbf{centralisées} (pouvoir unique, administration, armée) ou \textbf{structurées} (chefferies à sociétés de régulation).
\item L'islam est facteur d'unification au Nord ; la tradition et le sacré structurent l'Ouest.
\end{itemize}
\end{aretenir}

\courssub{Le lamidat : l'exemple des lamidats peuls du Nord-Cameroun}

\emph{Origines.} Au début du XIX\textsuperscript{e} siècle, le \cle{djihad} lancé par \textbf{Ousman dan Fodio} depuis Sokoto (1804) se propage vers le sud. Son disciple \textbf{Modibo Adama} mène la conquête de l'Adamaoua et des plaines de la Bénoué (vers 1809-1847). Il fonde l'\cle{Émirat de l'Adamaoua} (Yola), subdivisé en \cle{lamidats} autonomes mais vassaux.

\begin{definition}[Lamidat]
Entité politique peule dirigée par un \emph{lamido} (chef politique, militaire, judiciaire et religieux), issue de la conquête djihadiste, organisée sur une base théocratique musulmane.
\end{definition}

\courssub{Organisation politique et administrative}
\begin{itemize}
\item \textbf{Le Lamido} : chef suprême, \emph{Amir al-Mu'minin} (commandeur des croyants) local. Pouvoir héréditaire (choix parmi les fils du défunt par un conseil d'électeurs : \emph{galadima}, \emph{modibo}, chefs de guerre).
\item \textbf{L'administration centrale} (cour du lamido) :
  \begin{itemize}
  \item \emph{Galadima} : premier ministre, chef de l'administration, régent en cas de vacance.
  \item \emph{Modibo} (ou \emph{Limane}) : grand imam, conseiller religieux, justice islamique (\emph{sharia}).
  \item \emph{Sarkin Yaki} : chef de l'armée, organisation de la défense et des razzias.
  \item \emph{Waziri} : trésorier, gestion des tributs et du butin.
  \end{itemize}
\item \textbf{L'administration territoriale} : Le lamidat est divisé en \emph{chefferies de canton} (confiées à des \emph{lawanes} pour les sédentaires) et en \emph{chefferies de campements} (confiées à des \emph{djaouros} ou \emph{ardos} pour les pasteurs nomades). Ces chefs perçoivent l'impôt (\emph{jangali} sur le bétail, \emph{kharaj} sur les terres) et l'acheminent au lamido.
\end{itemize}

\courssub{Organisation sociale et économique}
\begin{itemize}
\item \textbf{Stratification sociale} :
  \begin{enumerate}
  \item \emph{Fulbe} (Peuls) conquérants, musulmans, aristocrates guerriers.
  \item Populations \emph{converties} (Mboum, Toupouri, etc.), intégrées, payent l'impôt \emph{zakât}.
  \item \emph{Kirdi} (terme péjoratif signifiant « païens ») : populations non islamisées (Montagnards, Koma, etc.), soumises au tribut, pourvoyeuses d'esclaves (\emph{maccube}) pour le service domestique, l'armée ou le commerce transsaharien.
  \end{enumerate}
\item \textbf{Économie} : Élevage bovin (zébu) dominant, agriculture (mil, sorgho, maïs, coton), artisanat (tissage, cuir), commerce caravanier (kola, sel, bétail, esclaves) vers le Nord (Sokoto, Bornou) et la côte.
\end{itemize}

\begin{exemplebox}[Le lamidat de Garoua]
Fondé vers 1830 par \textbf{Modibo Haman Njoundi} (fils de Modibo Adama). Situé sur la Bénoué, port fluvial majeur. Devient un centre commercial prospère (coton, ivoire, esclaves) et un relais du djihad vers le sud. Le lamido y réside dans un palais fortifié (\emph{birni}).
\end{exemplebox}

\courssub{La chefferie des Grassfields}

\emph{Cadre géographique et humain.} Région de haute altitude (1000-2500 m), volcanique, densément peuplée (Bamiléké, Bamoun, Tikar, populations du Nord-Ouest : Bafut, Bamenda, Kom, Nso). Climat favorable à l'agriculture intensive.

\begin{definition}[Chefferie des Grassfields (Fon/Fom)]
Entité politique centralisée dirigée par un \emph{Fon} (ou \emph{Fom}), roi sacré, dont le pouvoir est encadré et contrebalancé par des \cle{sociétés secrètes} (ou sociétés de régulation) et un collège de notables.
\end{definition}

\courssub{Organisation politique : le Fon et les contre-pouvoirs}
\begin{itemize}
\item \textbf{Le Fon (Fom)} : Chef suprême, \textbf{sacré} (intermédiaire avec les ancêtres), garant de la fécondité (terre, femmes, bétail). Désignation : héréditaire (souvent le fils cadet ou un neveu maternel selon la coutume), intronisation par les sociétés secrètes et les notables. Il ne paraît pas en public (port du masque, voix déguisée).
\item \textbf{Les sociétés de régulation (contre-pouvoirs traditionnels)} :
  \begin{itemize}
  \item \emph{Kwifon} (ou \emph{Ngumba}, \emph{Mabusi}, \emph{Takumbeng} chez les femmes) : société masculine (ou féminine) masquée, gardienne de la tradition, police religieuse et judiciaire, contrôle le Fon (peut le destituer), organise les funérailles royales.
  \item \emph{Nguon} (Bamoun) : grande fête quinquennale de régénération du pouvoir.
  \end{itemize}
\item \textbf{Les notables} (\emph{Nkam}, \emph{Takam}, \emph{Kamveu}) : Chefs de lignages, de quartiers, porte-parole (\emph{Nkwi}), chefs de guerre. Ils forment le conseil du Fon, gèrent la justice coutumière, la distribution des terres, les travaux communautaires.
\item \textbf{Organisation territoriale} : Le royaume (chefferie) est divisé en \emph{quartiers} (\emph{tu}) dirigés par des chefs de quartier, eux-mêmes subdivisés en \emph{compounds} (concessions familiales).
\end{itemize}

\courssub{Économie et société}
\begin{itemize}
\item \textbf{Agriculture intensive} : Ignames, taro, maïs, banane plantain, café (introduit plus tard), cultures maraîchères. Fumure, terrasses, jachère courte.
\item \textbf{Élevage} : Petits ruminants (chèvres, moutons), porcs, volailles. Bovins (zébu) chez les Bamiléké de l'Est (Banyo, Tibati) et au Nord-Ouest.
\item \textbf{Artisanat et commerce} : Tissage (ndop, tissus à raphia), sculpture sur bois/ivoire, ferronnerie, poterie. Marchés périodiques (8 jours) très actifs, réseaux commerciaux à longue distance (kola, sel, perles).
\item \textbf{Société} : Hiérarchisée (Fon, notables, hommes libres, captifs/esclaves \emph{mbu}). L'esclavage est domestique, non chattel ; intégration possible.
\end{itemize}

\begin{exemplebox}[La chefferie de Bafut (Nord-Ouest)]
L'une des plus puissantes chefferies Tikar/Bamileke du Nord-Ouest. Le Fon de Bafut règne sur une confédération de villages. Résistance célèbre face aux Allemands (guerre de Bafut, 1901-1910). Palais traditionnel (Achum) classé patrimoine UNESCO.
\end{exemplebox}

\courssub{Le sultanat bamoun}

\emph{Origines et dynastie.} Peuple bamoun (langue bantoïde), installé dans la plaine du Noun (Foumban) depuis le XIV\textsuperscript{e} s. (dynastie Ncharé Yen). Au XIX\textsuperscript{e} s., le sultanat s'affirme comme une puissance régionale.

\begin{definition}[Sultanat bamoun (Mfon)]
Royaume centralisé de la plaine du Noun, dirigé par un \emph{Mfon} (sultan) aux pouvoirs politiques, militaires, judiciaires et culturels étendus, doté d'une administration de cour sophistiquée et d'une culture écrite originale.
\end{definition}

\courssub{L'apogée : le règne du sultan Njoya (vers 1886-1931)}
\begin{itemize}
\item \textbf{Modernisation politique} : Création d'une administration bureaucratique (ministres : \emph{Nji} pour l'intérieur, \emph{Mfoumpou} pour la guerre, \emph{Nguon} pour les finances), codification des lois, recensement de la population.
\item \textbf{Innovation culturelle majeure} : Invention de l'\cle{écriture shümom} (a-ka-u-ku), syllabaire d'environ 70 signes (vers 1896-1910), création d'une imprimerie, d'écoles, d'un musée, traduction du Coran et de l'histoire bamoun.
\item \textbf{Architecture} : Construction du \textbf{palais royal de Foumban} (1917-1922), synthèse de l'architecture traditionnelle (haute toiture conique) et coloniale (briques, symétrie).
\item \textbf{Diplomatie} : Alliance précoce avec les Allemands (1902) pour contrer les voisins (Nso, Bafut) et préserver l'autonomie interne. Collaboration intellectuelle avec l'ethnologue allemand Ernst Tessemacher.
\end{itemize}

\courssub{Organisation sociale et économique}
\begin{itemize}
\item \textbf{Société} : Sultan (\emph{Mfon}) $\rightarrow$ Princes/Princesses ($\emph{Nji}$, \emph{Mfoumpou}) $\rightarrow$ Notables/Administrateurs $\rightarrow$ Artisans/Commerçants (corporations organisées) $\rightarrow$ Paysans $\rightarrow$ Captifs (\emph{Mbu}).
\item \textbf{Économie} : Agriculture diversifiée (maïs, igname, arachide, coton), élevage, artisanat d'art réputé (bronze, perles, tissus \emph{ndop}), commerce régional (kola, sel, fer), tributs des provinces.
\end{itemize}

\begin{savaistu}[Le Shümom : une écriture africaine vivante]
L'écriture shümom, inventée par Njoya, a connu plusieurs réformes (de 500+ pictogrammes à ~70 syllabes). Elle est encore enseignée aujourd'hui à Foumban. C'est l'un des rares exemples d'écriture africaine indigène créée au contact de la modernité coloniale mais hors missionnarisation. Le palais de Foumban abrite les archives royales écrites en shümom.
\end{savaistu}

\courssub{Étude comparée : trois modèles d'organisation politique}

\begin{methode}[Comparer les trois modèles d'organisation politique]
\begin{enumerate}
\item Choisir des \cle{critères de comparaison} identiques pour les trois entités : chef, nature du pouvoir, organisation sociale, économie, religion, exemples.
\item Remplir le tableau ligne par ligne, avec des informations \textbf{brèves et exactes}.
\item Dégager ensuite les \cle{points communs} (pouvoir centralisé, hiérarchie sociale, économie agraire) et les \cle{différences} (religion, degré de centralisation, mode de désignation du chef).
\item Rédiger une courte conclusion comparative.
\end{enumerate}
\end{methode}

\begin{center}
\begin{tabularx}{\linewidth}{|l|X|X|X|}
\hline
\rowcolor{popblueL} \textbf{Critère} & \textbf{Lamidat peul (Nord)} & \textbf{Chefferie des Grassfields (Ouest/NO)} & \textbf{Sultanat bamoun (Noun)} \\
\hline
\textbf{Chef (Titre)} & Lamido & Fon (ou Fom) & Sultan (Mfon) \\
\hline
\textbf{Nature du pouvoir} & Théocratique musulman (politique + religieux), héréditaire (élection par notables) & Sacralisé, encadré par sociétés secrètes (Kwifon), héréditaire (coutume matrilinéaire/patrilinéaire) & Politique, militaire, culturel, bureaucratique, héréditaire \\
\hline
\textbf{Administration} & Centrale (Galadima, Modibo, Sarkin Yaki) + Territoriale (Lawanes, Djaouros) & Cour (Notables, Nkam) + Sociétés secrètes (Kwifon, Ngumba) + Quartiers & Bureaucratie de cour (Nji, Mfoumpou, scribes), provinces \\
\hline
\textbf{Organisation sociale} & Stricte : Fulbe (conquérants) > Convertis > Kirdi (tributaires/esclaves) & Fon > Notables/Sociétés > Hommes libres > Captifs (intégrables) & Sultan > Princes/Notables > Corporations d'artisans > Paysans > Captifs \\
\hline
\textbf{Religion / Légitimité} & Islam (Charia), Djihad & Croyances traditionnelles, Culte des ancêtres, Sacralité du Fon & Islam (minoritaire, cour) + Tradition + Culture écrite (Shümom) \\
\hline
\textbf{Économie de base} & Élevage bovin, agriculture pluviale, tributs, commerce caravanier (esclaves, kola, sel) & Agriculture intensive (igname, maïs), petit élevage, artisanat, marchés périodiques & Agriculture, artisanat d'art (bronze, perles), commerce, tributs \\
\hline
\textbf{Exemples majeurs} & Garoua, Maroua, Ngaoundéré, Banyo, Tibati, Rey-Bouba & Bafut, Bamenda, Bandjoun, Bana, Bangangté, Kumbo (Nso) & Foumban (capitale unique) \\
\hline
\end{tabularx}
\end{center}

\begin{aretenir}[Synthèse comparative]
\begin{itemize}
\item \textbf{Point commun majeur :} Les trois sont des \textbf{États centralisés} ou des chefferies \textbf{fortement structurées}, avec une administration, une armée, une justice, une fiscalité et une idéologie légitimante (Islam, Ancêtres, Culture royale).
\item \textbf{Différence fondamentale :} La \textbf{source de légitimité} et le \textbf{mode de contrôle du pouvoir}. Le Lamido tire son pouvoir de l'Islam et de la conquête (théocratie). Le Fon tire son pouvoir du sacré ancestral mais est \textbf{contrôlé par des sociétés secrètes} (système de freins et contrepoids). Le Sultan Njoya construit une \textbf{légitimité moderne et culturelle} (écriture, bureaucratie, diplomatie) tout en gérant l'héritage traditionnel.
\item \textbf{Conclusion :} À la veille de 1884, le Cameroun présente une \textbf{diversité politique cohérente} : des États aboutis, capables de diplomatie, de guerre, d'innovation culturelle et de gestion territoriale. La colonisation allemande ne rencontre pas un « chaos » mais des partenaires/rivaux organisés.
\end{itemize}
\end{aretenir}

\courssub{Travaux Dirigés : L'organisation d'une chefferie de la localité de l'établissement}

\emph{Objectif pédagogique.} Appliquer les notions du cours (chef, notables, sociétés, économie, rites) à une réalité locale observable. Développer l'esprit d'enquête, la rigueur de la collecte orale et la rédaction d'un rapport structuré.

\begin{methode}[Étudier l'organisation d'une chefferie de la localité de l'établissement (TD)]
\begin{enumerate}
\item \cle{Préparer l'enquête} : Choisir la chefferie (chef supérieur, chef de groupement, chef de village selon la localité). Rédiger un \textbf{guide d'entretien} structuré : Identité du chef (nom, titre, lignée) ; Mode de désignation (héritage, élection, rites d'intronisation) ; Dignitaires et leurs fonctions (notables, chefs de quartiers, sociétés traditionnelles, porte-parole) ; Justice coutumière ; Gestion du foncier ; Fêtes et rites annuels ; Relations avec l'administration moderne (sous-préfet, maire).
\item \cle{Collecter les informations} : Interroger le chef traditionnel (après prise de rendez-vous protocolaire), les notables, les anciens, les jeunes. Consulter les archives locales (registres, arrêtés de nomination). Observer les lieux (palais, place des fêtes, cases des notables, objets symboliques : tambours, masques, trônes). \emph{Respecter l'éthique : demander l'autorisation pour photos/enregistrement, respecter les secrets d'initiation.}
\item \cle{Organiser les données} : Classer les réponses selon le plan du rapport. Croiser les sources orales (tradition) et écrites (administration). Identifier les continuités et ruptures avec le modèle précolonial étudié.
\item \cle{Rédiger un rapport structuré} :
  \begin{itemize}
  \item \textbf{Introduction} : Localisation, peuple, raison du choix, méthode.
  \item \textbf{Développement} : 
    \begin{enumerate}
    \item Le Chef et son mode de désignation (titre, lignée, rites, légitimité).
    \item L'équipe dirigeante (notables, sociétés, quartiers) et leurs fonctions précises.
    \item Les fonctions actuelles de la chefferie (justice, foncier, cohésion sociale, développement, interface administration/population).
    \end{enumerate}
  \item \textbf{Conclusion} : Bilan : la chefferie, institution vivante, héritière du précolonial, adaptée au contexte moderne. Propositions (musée local, documentation, valorisation).
  \end{itemize}
\item \cle{Présenter oralement} les résultats devant la classe (exposé de 10 min + 5 min questions). Évaluation : rigueur de l'enquête, clarté de l'exposé, utilisation du vocabulaire du cours (lamido, fon, notable, société secrète, tribut, etc.).
\end{enumerate}
\end{methode}

\begin{exoresolu}[Rapport d'enquête type : Chefferie de Bafia (exemple)]
\textbf{Énoncé.} Rédige le rapport d'enquête sur la chefferie de ta localité en suivant le plan imposé. (Note sur 10 : Introduction 1, Chef 2, Dignitaires 3, Rôle 2, Conclusion 1, Forme 1).
\tcblower
\textbf{Solution détaillée (modèle de rédaction attendue).}

\emph{Introduction.} Notre enquête, menée du 10 au 15 octobre 202X, porte sur la chefferie supérieure de Bafia (région du Centre, peuple Bafia/Rikpa). Elle vise à analyser la permanence des institutions traditionnelles. Nous avons interviewé Sa Majesté \emph{Tatadjou} (chef supérieur), trois notables (\emph{Mban}) et consulté l'arrêté de nomination n° 2015/... du MINATD.

\emph{1. Le Chef et son mode de désignation.} Le chef porte le titre de \textbf{Tatadjou} (« Père du peuple »). La désignation est \textbf{héréditaire matrilinéaire} : le successeur est choisi parmi les neveux du défunt (fils de ses sœurs) par le collège des \emph{Mban} (notables) et les chefs de clan, selon des rites secrets (retraite au \emph{Laak}, scarifications, remise des insignes : collier de léopard, lance, siège). L'intronisation est validée par un arrêté préfectoral. Son pouvoir est politique (chef de l'administration coutumière), judiciaire (tribunal coutumier) et sacré (garant de la fécondité, lien avec les ancêtres \emph{Ngok}).

\emph{2. Les dignitaires et l'administration.} Le Tatadjou est assisté de :
\begin{itemize}
\item Les \textbf{Mban} (9 notables principaux) : ministres de fait (Mban Nkwa : justice ; Mban Ntou : finances/terres ; Mban Ngwa : protocole/porte-parole). Ils forment le \emph{Conseil Royal}.
\item Les \textbf{Chefs de groupements/villages} (\emph{Bekom}) : relais territoriaux, percepteurs de la taxe coutumière (\emph{Nkwa}).
\item La \textbf{Société \emph{Mabuh}} (hommes) et \textbf{\emph{Takumbeng}} (femmes) : régulation sociale, police des mœurs, organisation des funérailles, médiation.
\item Le \textbf{Secrétaire du Chef} : interface avec l'administration moderne (état civil, élections).
\end{itemize}

\emph{3. Le rôle dans la vie locale.} La chefferie règle 80\% des conflits fonciers et successoraux (désengorgement des tribunaux). Elle organise la fête annuelle \emph{Ngoun} (récolte, unité). Elle mobilise la population pour l'entretien des routes rurales et la construction de l'école/centre de santé (travaux d'intérêt collectif \emph{Njangi}). Elle est consultée pour toute implantation d'entreprise (carrières, plantations).

\emph{Conclusion.} La chefferie de Bafia, bien qu'encadrée par la loi n° 77/245 sur l'organisation traditionnelle, conserve une autorité morale et administrative réelle. Elle illustre la résilience du modèle « chefferie des Grassfields » (chef sacré + notables + sociétés de régulation) adapté à la décentralisation actuelle.
\end{exoresolu}

\courssub{Fiche méthode 1 : Appliquer les notions de l'unité à une situation concrète}

\begin{methode}[Identifier et décrire une entité politique précoloniale]
Pour décrire correctement un lamidat, une chefferie des Grassfields ou le sultanat bamoun, suivre les étapes suivantes :
\begin{enumerate}
\item \cle{Nommer} précisément l'entité et la localiser (région, peuple, espace géographique).
\item Identifier le \cle{chef} (titre exact : lamido, fon/fom, sultan/mfon) et la nature de son pouvoir (politique, religieux, militaire, judiciaire).
\item Décrire l'\cle{organisation sociale et administrative} : hiérarchie, dignitaires, notables, quartiers ou provinces.
\item Préciser les \cle{fondements économiques} : agriculture, élevage, commerce, tributs, esclavage.
\item Conclure par les \cle{rapports avec l'extérieur} (guerres, alliances, commerce) et la place de l'entité dans l'espace camerounais précolonial.
\end{enumerate}
Réflexe : toujours donner au moins un exemple concret (un lamidat précis, une chefferie précise) pour illustrer la notion.
\end{methode}

\begin{exoresolu}[Décrire un lamidat peul du Nord-Cameroun]
\textbf{Énoncé (type Probatoire).} Après avoir défini le terme \emph{lamidat}, décris l'organisation politique et sociale d'un lamidat peul du Nord-Cameroun de ton choix. (4 points)
\tcblower
\textbf{Solution détaillée.}

\emph{Étape 1 — Définition.} Un \cle{lamidat} est une entité politique dirigée par un \emph{lamido} (du peul \emph{lamdo}, « celui qui commande »), chef à la fois politique et religieux. Les lamidats du Nord-Cameroun sont nés de la conquête peule (djihad) menée au début du XIX\textsuperscript{e} siècle par les disciples d'Ousman dan Fodio, notamment le Modibo Adama.

\emph{Étape 2 — Exemple et localisation.} Prenons le lamidat de \textbf{Garoua} (ou celui de Maroua, Ngaoundéré, Banyo, Tibati). Il est situé dans la plaine de la Bénoué, au Nord-Cameroun.

\emph{Étape 3 — Organisation politique.} À la tête du lamidat se trouve le \textbf{lamido}, qui détient le pouvoir politique, militaire, judiciaire et religieux (l'islam est la religion d'État). Il est assisté de dignitaires : le \emph{galadima} (conseiller principal), le \emph{modibo} (savant religieux), des chefs de guerre (\emph{sarkin yaki}). Le lamidat est subdivisé en quartiers et en circonscriptions dirigés par des chefs subalternes (\emph{lawanes} pour les sédentaires, \emph{djaouros} pour les pasteurs), qui rendent compte au lamido.

\emph{Étape 4 — Organisation sociale et économique.} La société est hiérarchisée : au sommet, les Peuls musulmans conquérants ; ensuite les populations converties ; à la base, les populations non islamisées (\emph{kirdi}), soumises au tribut et parfois réduites en esclavage. L'économie repose sur l'élevage bovin, l'agriculture (mil, sorgho, coton), le commerce caravanier et les tributs versés par les sujets.

\emph{Étape 5 — Conclusion.} Le lamidat est donc un État centralisé, de type théocratique et militaire, qui a structuré durablement le Nord-Cameroun avant l'arrivée des colonisateurs allemands.

\emph{Phrase de conclusion.} Ainsi, le lamidat de Garoua illustre l'existence, à la veille de la colonisation, d'États organisés et hiérarchisés au Cameroun.
\end{exoresolu}

\courssub{Fiche méthode 2 : Rédiger un commentaire de document historique}

\begin{methode}[Commenter un document (texte, carte, image) en Histoire]
\begin{enumerate}
\item \cle{Présenter le document} : nature (texte, récit de voyage, carte, photographie), auteur, date, source, contexte.
\item \cle{Dégager le thème} : de quoi parle le document ? Le relier au chapitre étudié (ici : les entités politiques précoloniales).
\item \cle{Analyser} : relever les informations essentielles (faits, acteurs, lieux, dates) et les expliquer à l'aide de ses connaissances du cours.
\item \cle{Apprécier la portée et les limites} : que nous apprend le document ? Est-il partial (regard d'un explorateur européen, par exemple) ?
\item \cle{Conclure} : répondre clairement à la question posée en une phrase synthétique.
\end{enumerate}
Réflexe : citer le document (« le texte indique que... », « l'auteur précise... ») pour justifier chaque affirmation.
\end{methode}

\begin{exoresolu}[Commentaire d'un extrait sur le royaume bamoun]
\textbf{Énoncé (type Probatoire).} Document : « Arrivé à Foumban, nous fûmes reçus par le sultan Njoya, souverain puissant, entouré de ses conseillers. Sa cour compte des scribes, des artistes, des guerriers. Le roi a fait inventer une écriture pour son peuple et fait construire un palais remarquable. » (Récit d'un explorateur allemand, vers 1902).

Présente le document, puis montre ce qu'il révèle de l'organisation du sultanat bamoun. (5 points)
\tcblower
\textbf{Solution détaillée.}

\emph{Étape 1 — Présentation.} Il s'agit d'un extrait de \textbf{récit de voyage} rédigé par un explorateur allemand vers 1902, époque de la pénétration coloniale allemande au Cameroun (le Kamerun allemand est proclamé en 1884). Le document décrit l'accueil réservé à l'explorateur à Foumban, capitale du royaume bamoun.

\emph{Étape 2 — Thème.} Le thème est l'organisation politique et culturelle du \cle{sultanat bamoun} à la veille de la colonisation, sous le règne du sultan \textbf{Njoya} (vers 1886-1931).

\emph{Étape 3 — Analyse.} Le document révèle plusieurs traits essentiels :
\begin{itemize}
\item « souverain puissant, entouré de ses conseillers » : le pouvoir est \textbf{centralisé} entre les mains du sultan (mfon), secondé par des dignitaires et des notables (le conseil des notables, les \emph{nji} et \emph{komguem} qui encadrent la société) ;
\item « des scribes, des artistes, des guerriers » : la cour est un centre administratif, militaire et artistique, signe d'un État structuré ;
\item « une écriture pour son peuple » : le sultan Njoya a fait créer l'écriture \textbf{shümom} (a-ka-u-ku), preuve d'un remarquable dynamisme culturel ;
\item « un palais remarquable » : le palais de Foumban (achevé vers 1917) symbolise la puissance royale et l'architecture bamoun.
\end{itemize}

\emph{Étape 4 — Portée et limites.} Le document montre que le royaume bamoun était un État organisé, lettré et ouvert aux échanges (Njoya collabora avec les Allemands). Mais il s'agit du regard d'un Européen : il faut le croiser avec d'autres sources, notamment la tradition orale bamoun.

\emph{Conclusion.} Ce récit confirme que le sultanat bamoun était, à la veille de la colonisation, un royaume centralisé et brillant, dirigé par un souverain réformateur, le sultan Njoya.
\end{exoresolu}

\courssub{Fiche méthode 3 : Construire un tableau comparatif}

\begin{methode}[Comparer les trois modèles d'organisation politique]
\begin{enumerate}
\item Choisir des \cle{critères de comparaison} identiques pour les trois entités : chef, nature du pouvoir, organisation sociale, économie, religion, exemples.
\item Remplir le tableau ligne par ligne, avec des informations \textbf{brèves et exactes}.
\item Dégager ensuite les \cle{points communs} (pouvoir centralisé, hiérarchie sociale, économie agraire) et les \cle{différences} (religion, degré de centralisation, mode de désignation du chef).
\item Rédiger une courte conclusion comparative.
\end{enumerate}
\end{methode}

\begin{exoresolu}[Tableau comparatif des trois entités politiques]
\textbf{Énoncé (type Probatoire).} À l'aide d'un tableau, compare le lamidat peul, la chefferie des Grassfields et le sultanat bamoun selon quatre critères : le chef, la nature du pouvoir, l'organisation sociale, l'économie. Dégage ensuite un point commun et une différence. (6 points)
\tcblower
\textbf{Solution détaillée.}

\begin{center}
\begin{tabularx}{\linewidth}{|l|X|X|X|}
\hline
\rowcolor{popblueL} \textbf{Critère} & \textbf{Lamidat peul} & \textbf{Chefferie des Grassfields} & \textbf{Sultanat bamoun} \\
\hline
Chef & Lamido & Fon (ou fom) & Sultan (mfon) \\
\hline
Nature du pouvoir & Politique et religieux (islam), héréditaire & Politique et sacré, entouré de sociétés secrètes & Politique, militaire et culturel, héréditaire \\
\hline
Organisation sociale & Hiérarchie stricte : Peuls, convertis, populations tributaires & Notables, sociétés régulatrices (takumbeng, mabusi), quartiers & Sultan, conseil des notables, scribes, guerriers, artisans \\
\hline
Économie & Élevage, agriculture, tributs, commerce caravanier & Agriculture (maïs, ignames), élevage, artisanat, marchés & Agriculture, artisanat d'art, commerce, tributs \\
\hline
\end{tabularx}
\end{center}

\emph{Point commun.} Les trois entités sont des formations politiques \textbf{centralisées et hiérarchisées}, dirigées par un chef dont le pouvoir est à la fois politique et sacré ou religieux, et reposant sur une économie agraire et le prélèvement de tributs.

\emph{Différence.} Le lamidat est un État \textbf{théocratique musulman} issu d'une conquête (djihad), alors que la chefferie des Grassfields repose sur des \textbf{croyances traditionnelles} et des sociétés secrètes qui encadrent et limitent le pouvoir du fon.

\emph{Conclusion.} Cette comparaison montre la diversité des modèles politiques qui structuraient le Cameroun avant la colonisation.
\end{exoresolu}

\courssub{Fiche méthode 4 : Construire et exploiter une frise chronologique}

\begin{methode}[Réaliser une frise chronologique en Histoire]
\begin{enumerate}
\item Relever les \cle{dates clés} du chapitre (fondations, règnes, conquêtes, arrivée des Européens).
\item Choisir une \cle{échelle} régulière (par exemple 1 cm = 5 ans sur 13 cm pour 1800-1890) et tracer une ligne du temps horizontale.
\item Placer chaque événement au bon endroit, avec un libellé \textbf{court} (date + fait).
\item Vérifier l'\cle{ordre chronologique} : aucune date ne doit être inversée.
\end{enumerate}
Réflexe : au Probatoire, une frise demandée doit comporter au minimum 4 à 5 dates exactes et correctement espacées.
\end{methode}

\begin{exoresolu}[Frise : le Cameroun à la veille de la colonisation]
\textbf{Énoncé (type Probatoire).} Construis une frise chronologique situant au moins cinq événements marquants de l'histoire politique du Cameroun précolonial, du début du XIX\textsuperscript{e} siècle à 1884. (4 points)
\tcblower
\textbf{Solution détaillée.}

Dates retenues : début du djihad peul au Nord-Cameroun (vers 1804-1809) ; arrivée du Modibo Adama et fondation des lamidats (vers 1809-1840) ; essor des chefferies des Grassfields (XIX\textsuperscript{e} s.) ; traités germano-douala (12 juillet 1884) ; début du règne de Njoya (1886).

\begin{popfigure}
\begin{tikzpicture}[scale=1]
% Axe principal : 13cm pour 90 ans (1800-1890). 1 an = 0.1444 cm. Origine 1800 à x=0.5
\draw[thick,popink,->] (0.5,0) -- (13.5,0);
% Graduations tous les 20 ans
\foreach \x/\d in {0.5/1800, 3.39/1820, 6.28/1840, 9.16/1860, 12.05/1880}
  \draw[popink] (\x,0.1) -- (\x,-0.1) node[below,font=\small]{\d};
% Événements (calculs : x = 0.5 + (année-1800)*0.1444)
% 1806 (milieu 1804-1809) -> 1.37
\fill[poporange] (1.37,0) circle (0.09); 
\node[above,font=\small,align=center] at (1.37,0.15) {1804-1809 : djihad\\peul (O. dan Fodio)};
% 1824 (milieu 1809-1840) -> 3.96
\fill[popblue] (3.96,0) circle (0.09); 
\node[below=0.55cm,font=\small,align=center] at (3.96,-0.15) {1809-1840 : Modibo Adama,\\fondation des lamidats};
% 1850 (essor Grassfields) -> 7.72
\fill[popgreen] (7.72,0) circle (0.09); 
\node[above,font=\small,align=center] at (7.72,0.15) {XIX\textsuperscript{e} s. : essor des\\chefferies des Grassfields};
% 1884 (traites) -> 12.63
\fill[poppurple] (12.63,0) circle (0.09); 
\node[below=0.55cm,font=\small,align=center] at (12.63,-0.15) {12 juil. 1884 : traités\\germano-douala};
% 1886 (Njoya) -> 12.92
\fill[popgold] (12.92,0) circle (0.09); 
\node[above,font=\small,align=center] at (12.92,0.15) {1886 : début du\\règne de Njoya};
\end{tikzpicture}
\legende{Frise chronologique : les entités politiques du Cameroun à la veille de la colonisation (1800-1890). Échelle : 1 cm $\approx$ 7 ans.}
\end{popfigure}

\emph{Vérification.} L'ordre chronologique est respecté : 1804-1809 $\rightarrow$ 1809-1840 $\rightarrow$ 1850 (essor) $\rightarrow$ 1884 $\rightarrow$ 1886. Chaque événement est daté et libellé brièvement.

\emph{Conclusion.} La frise montre que les grandes entités politiques camerounaises se sont constituées ou consolidées au XIX\textsuperscript{e} siècle, juste avant la colonisation allemande amorcée en 1884.
\end{exoresolu}

\courssub{Fiche méthode 5 : Mener une enquête sur une chefferie locale (TD)}

\begin{methode}[Étudier l'organisation d'une chefferie de la localité de l'établissement]
\begin{enumerate}
\item \cle{Préparer l'enquête} : choisir la chefferie, rédiger un guide d'entretien (questions sur le chef, les notables, les rites, l'histoire).
\item \cle{Collecter les informations} : interroger le chef traditionnel, les notables, les anciens ; consulter les documents disponibles ; observer les lieux (palais, place des rites).
\item \cle{Organiser les données} : identifier le chef et son mode de désignation, les dignitaires, les groupes ou sociétés, les fonctions économiques et culturelles.
\item \cle{Rédiger un rapport} structuré : introduction (présentation de la chefferie), développement (organisation politique, sociale, économique), conclusion (rôle actuel de la chefferie).
\item \cle{Présenter oralement} les résultats devant la classe et répondre aux questions.
\end{enumerate}
\end{methode}

\begin{exoresolu}[Rapport d'enquête sur une chefferie locale]
\textbf{Énoncé (type Probatoire).} Dans le cadre du TD, tu as enquêté sur la chefferie de ta localité. Rédige un rapport structuré présentant : le chef et son mode de désignation, les principaux dignitaires, le rôle de la chefferie dans la vie locale. (5 points)
\tcblower
\textbf{Solution détaillée (exemple de rédaction).}

\emph{Introduction.} Notre enquête porte sur la chefferie de notre localité. Elle vise à décrire son organisation politique et sociale et à montrer son rôle dans la vie quotidienne de la population.

\emph{1. Le chef et son mode de désignation.} La chefferie est dirigée par un \cle{chef traditionnel} (fon, fom, ou chef supérieur selon la région). Sa désignation est \textbf{héréditaire} : il est choisi parmi les descendants de la lignée royale, selon des rites précis, puis intronisé lors de cérémonies traditionnelles. Son pouvoir est à la fois politique, judiciaire et sacré.

\emph{2. Les dignitaires.} Le chef est assisté de \textbf{notables} : chefs de quartiers, conseillers, porte-parole (linguiste), chefs de sociétés traditionnelles. Chacun exerce une fonction précise : justice de proximité, organisation des travaux collectifs, cérémonies, levée des impôts ou contributions.

\emph{3. Le rôle dans la vie locale.} La chefferie règle les conflits entre habitants, organise les fêtes traditionnelles et les travaux communautaires, veille sur les terres et sert d'intermédiaire entre la population et l'administration moderne.

\emph{Conclusion.} L'enquête montre que la chefferie, héritière des entités politiques précoloniales, demeure une institution vivante qui structure la vie sociale de notre localité.
\end{exoresolu}

\courssub{Fiche méthode 6 : Rédiger et justifier un paragraphe argumenté}

\begin{methode}[Rédiger un paragraphe argumenté en Histoire]
\begin{enumerate}
\item Formuler une \cle{phrase d'idée} (la thèse du paragraphe) en une phrase claire.
\item Développer avec \textbf{deux ou trois arguments}, chacun appuyé sur un \cle{fait précis} (date, acteur, lieu).
\item Utiliser des connecteurs logiques : d'abord, ensuite, en effet, par exemple, enfin.
\item Conclure le paragraphe par une phrase de \cle{bilan} qui répond à la question posée.
\item Relire : orthographe des noms propres, exactitude des dates, aucune affirmation sans justification.
\end{enumerate}
\end{methode}

\begin{exoresolu}[Paragraphe argumenté : un espace politiquement organisé]
\textbf{Énoncé (type Probatoire).} Rédige un paragraphe argumenté (10 à 15 lignes) montrant que « le Cameroun, à la veille de la colonisation, était un espace politiquement organisé ». (5 points)
\tcblower
\textbf{Solution détaillée.}

\emph{Phrase d'idée.} Loin d'être un espace désorganisé, le Cameroun précolonial était structuré par de véritables États et chefferies.

\emph{Arguments et faits.} D'abord, au Nord, les \textbf{lamidats peuls} (Garoua, Maroua, Ngaoundéré, Banyo), issus du djihad du Modibo Adama au début du XIX\textsuperscript{e} siècle, constituaient des États centralisés dirigés par des lamidos détenant les pouvoirs politique, militaire et religieux. Ensuite, dans les \textbf{Grassfields} de l'Ouest, des chefferies puissantes comme celle de Bamenda ou de Bafut étaient dirigées par des \emph{fons} entourés de notables et de sociétés régulatrices qui encadraient la vie sociale et économique. Enfin, le \textbf{sultanat bamoun}, avec sa capitale Foumban, illustre le degré d'organisation atteint : le sultan Njoya (règne à partir de 1886) y fit créer une écriture, le shümom, édifia un palais et organisa une administration de cour.

\emph{Phrase de bilan.} Ces exemples prouvent que le territoire camerounais, avant 1884, était couvert d'entités politiques hiérarchisées et centralisées : la colonisation n'a donc pas rencontré un vide politique, mais des États en pleine activité.
\end{exoresolu}

\begin{attention}[Les 5 erreurs qui coûtent des points au Probatoire]
\begin{enumerate}
\item \textbf{Confondre les titres des chefs} : lamido (lamidat peul), fon ou fom (Grassfields), sultan ou mfon (Bamoun). Intervertir ces titres est une erreur éliminatoire dans une définition.
\item \textbf{Affirmer que le Cameroun précolonial n'avait « pas d'État »} : c'est faux ; il faut au contraire citer des entités précises (lamidats, chefferies, sultanat bamoun) avec des exemples localisés.
\item \textbf{Inverser les dates} sur une frise : le djihad peul (début XIX\textsuperscript{e} s.) précède les traités germano-douala (1884) et le règne de Njoya (à partir de 1886). Vérifier toujours l'ordre chronologique.
\item \textbf{Commenter un document sans le présenter} : oublier la nature, l'auteur, la date ou le contexte du document fait perdre les points de la première étape du commentaire.
\item \textbf{Rédiger des affirmations sans preuves} : chaque argument doit être justifié par un fait précis (date, acteur, lieu). Une copie qui généralise sans exemple (« les chefs étaient puissants ») ne peut pas obtenir la totalité des points.
\end{enumerate}
\end{attention}

\newpage

\coursec{Exercices}

\coursec{Le Cameroun à la veille de la colonisation}

\courssub{Je vérifie mes connaissances}

\begin{exercice}[QCM : les entités politiques précoloniales]
Pour chaque question, une seule réponse est correcte. Entoure la lettre correspondante.

\textbf{1.} Le chef d'un lamidat peul du Nord-Cameroun porte le titre de :
\begin{itemize}
\item[a)] fon
\item[b)] lamido
\item[c)] sultan
\end{itemize}

\textbf{2.} La capitale du royaume bamoun à la veille de la colonisation était :
\begin{itemize}
\item[a)] Bafoussam
\item[b)] Maroua
\item[c)] Foumban
\end{itemize}

\textbf{3.} Dans les Grassfields, le chef (fon) est assisté dans l'exercice du pouvoir par :
\begin{itemize}
\item[a)] des sociétés secrètes et des notables
\item[b)] une armée permanente de cavaliers
\item[c)] des gouverneurs de provinces nommés à vie
\end{itemize}

\textbf{4.} Le roi bamoun célèbre pour avoir créé une écriture et encouragé les arts est :
\begin{itemize}
\item[a)] Nsangu
\item[b)] Njoya
\item[c)] Bouba Ndjida
\end{itemize}
\end{exercice}

\begin{exercice}[Vrai ou faux, en justifiant]
Indique si chaque affirmation est vraie ou fausse, puis justifie ta réponse en une ou deux phrases.
\begin{enumerate}
\item À la veille de la colonisation, le territoire camerounais était dépourvu de toute organisation politique.
\item Les lamidats du Nord-Cameroun sont nés du jihad peul mené au début du XIX\textsuperscript{e} siècle.
\item La chefferie des Grassfields est dirigée par un chef élu pour un mandat de cinq ans.
\item Le sultanat bamoun disposait d'institutions centrales et d'une administration organisée.
\end{enumerate}
\end{exercice}

\begin{exercice}[Définitions]
Définis en deux ou trois phrases chacun des termes suivants : \cle{lamidat}, \cle{chefferie} (des Grassfields), \cle{sultanat}. Pour chaque terme, cite un exemple précis situé au Cameroun.
\end{exercice}

\begin{exercice}[Repères chronologiques]
Replace chacun des événements suivants sur une frise chronologique que tu traceras, en indiquant le siècle concerné :
\begin{enumerate}
\item le jihad d'Ousman dan Fodio et la formation des lamidats peuls du Nord-Cameroun ;
\item la fondation du royaume bamoun par Nchare Yen ;
\item l'arrivée des premiers Allemands à la côte camerounaise (1884) ;
\item le règne du roi Njoya à Foumban.
\end{enumerate}
Conclus en une phrase : que montre cette frise sur l'ancienneté des organisations politiques camerounaises ?
\end{exercice}

\begin{aretenir}[Synthèse : trois modèles politiques précoloniaux]
À la veille de la colonisation (1884), le Cameroun n'est pas un vide politique. On distingue trois grands types d'entités :
\begin{itemize}
\item \textbf{Les lamidats (Nord)} : nés du jihad d'Ousman dan Fodio (début XIXe), dirigés par un \cle{lamido} (chef politique et religieux), vassaux du Sultanat de Sokoto, organisation féodale, économie pastorale et commerciale.
\item \textbf{Les chefferies des Grassfields (Ouest/Nord-Ouest)} : très anciennes, dirigées par un \cle{fon} (chef sacré), pouvoir partagé avec des sociétés de régulation (Kwifoyn, Ngumba) et des notables, organisation en chefferies de premier et second rang.
\item \textbf{Le sultanat bamoun (Région de l'Ouest, capitale Foumban)} : royaume ancien (fondation ~XIVe s. par Nchare Yen), islamisé au XIXe, administration centrale (ministères), provinces, cour d'artisans, écriture (shümom) sous Njoya.
\end{itemize}
\end{aretenir}

\courssub{Je m'entraîne}

\begin{exercice}[Définition et distinction à partir d'un texte]
Lis le texte suivant :

\emph{« Au Nord, le lamido, chef politique et religieux, gouverne au nom de l'islam ; son autorité repose sur le jihad et sur la fidélité au sultan de Sokoto. À Foumban, le roi des Bamouns, à la tête d'un royaume ancien fondé par Nchare Yen, s'appuie sur une cour, des ministres et une administration qui contrôle les provinces et perçoit les tributs. »}

\begin{enumerate}
\item Définis les termes \cle{lamidat} et \cle{sultanat}.
\item Releve dans le texte deux éléments qui distinguent le lamidat du sultanat bamoun (fondement du pouvoir, origine de l'État).
\item Rédige trois phrases expliquant pourquoi ces deux États ne reposent pas sur les mêmes bases.
\end{enumerate}
\end{exercice}

\begin{exercice}[Localisation sur carte muette]
Sur une carte muette du Cameroun (fournie ou reproduite) :
\begin{enumerate}
\item Place trois lamidats du Nord : Rey-Bouba, Banyo, Maroua.
\item Place trois chefferies des Grassfields : Bafut, Bamenda (Mankon), Bali.
\item Place la capitale du royaume bamoun : Foumban.
\item Trace les limites approximatives des trois grandes zones culturelles (Nord soudanien, Grassfields, royaume bamoun) à l'aide de trois couleurs.
\item Rédige une légende complète et donne un titre à ta carte.
\end{enumerate}
\end{exercice}

\begin{popfigure}
\begin{center}
\begin{tikzpicture}[scale=0.7]
% Contour très simplifié du Cameroun
\draw[popdark, thick] (0,0) -- (2,0) -- (3,1) -- (4,2.5) -- (3.5,4) -- (2,5) -- (0.5,4.5) -- (-0.5,3) -- (-0.5,1) -- cycle;
% Fleuve Bénoué
\draw[popblue, thick] (1.5,0) .. controls (1.8,1.5) and (2,2.5) .. (2.2,3.5);
% Fleuve Sanaga
\draw[popblue, thick] (1,2) .. controls (1.5,3) .. (1.8,4);
% Lac Tchad (coin)
\draw[popblue, fill=popblue!20] (-0.5,0) rectangle (0,0.5);
% Villes / Points
\foreach \x/\y/\name in {
    0.5/0.8/Maroua, 1.5/1.5/Mora, 2.5/1.2/Kousseri,
    0.8/2.2/Garoua, 1.8/2.5/Ngaoundéré, 2.5/2.8/Meiganga,
    2.8/3.2/Banyo, 3.2/3.5/Rey-Bouba, 3.5/3.8/Tibati,
    1.2/3.8/Bafoussam, 0.8/4.2/Bamenda, 0.5/3.5/Bafut, 0.2/3.2/Bali,
    1.5/3.5/Foumban, 2.0/3.0/Bangangté, 2.5/2.0/Bertoua, 3.0/1.5/Batouri,
    1.0/1.0/Yagoua, 3.5/0.5/Kribi, 3.8/1.0/Ebolowa, 4.0/2.0/Yaoundé, 4.2/3.0/Douala
} {
    \fill[poporange] (\x,\y) circle (2pt);
    \node[font=\tiny, right=1pt, popdark] at (\x,\y) {\name};
}
% Zones culturales (hachures légères)
\pattern[pattern=north east lines, pattern color=poporange!30] (0,0) -- (2,0) -- (2.5,1.5) -- (1.5,2.5) -- (0.5,2) -- (-0.5,1) -- cycle; % Nord
\pattern[pattern=north west lines, pattern color=popgreen!30] (0.5,2) -- (1.5,2.5) -- (2.5,1.5) -- (3,2.5) -- (2.5,4) -- (0.5,4.5) -- (-0.5,3) -- cycle; % Grassfields
\pattern[pattern=dots, pattern color=poppurple!40] (1.2,3.2) circle (0.6); % Bamoun approx
% Légende interne
\node[font=\scriptsize, popdark, anchor=south west] at (-0.5,5) {Nord soudanien (hachures \textcolor{poporange}{orange})};
\node[font=\scriptsize, popdark, anchor=south west] at (-0.5,5.4) {Grassfields (hachures \textcolor{popgreen}{vert})};
\node[font=\scriptsize, popdark, anchor=south west] at (-0.5,5.8) {Royaume Bamoun (points \textcolor{poppurple}{violet})};
% Nord
\draw[->, thick] (-0.8,0.2) -- (-0.8,1.2) node[above, rotate=90, font=\tiny] {N};
\end{tikzpicture}
\end{center}
\legende{Carte de base du Cameroun pour l'exercice de localisation (contours, hydrographie, villes principales et zones culturelles indicatives).}
\end{popfigure}

\begin{exercice}[Tableau comparatif à compléter]
Recopie et complète le tableau suivant à l'aide de tes connaissances :

\begin{center}
\begin{tabularx}{\linewidth}{|l|X|X|X|}
\hline
\rowcolor{popblueL} \textbf{Critère} & \textbf{Lamidat} & \textbf{Chefferie des Grassfields} & \textbf{Sultanat bamoun} \\
\hline
Titre du chef & & & \\
\hline
Fondement du pouvoir & & & \\
\hline
Institutions principales & & & \\
\hline
Activités économiques dominantes & & & \\
\hline
\end{tabularx}
\end{center}

Conclus en trois phrases sur la diversité des organisations politiques du Cameroun à la veille de la colonisation.
\end{exercice}

\begin{exercice}[Étude de photographie : le palais de Foumban]
On étudie une photographie du palais royal de Foumban prise en 1917 : on y voit un vaste bâtiment de plan original, à façade ornée de piliers sculptés, entouré d'une place où se déroulent les grandes cérémonies ; le palais abrite la cour, les ateliers d'artisans (sculpteurs, forgerons, tisserands) et les services du royaume.
\begin{enumerate}
\item Décris en quelques lignes l'architecture du palais telle qu'elle apparaît sur la photographie.
\item Releve deux indices montrant que le palais est à la fois une résidence royale et un centre administratif et économique.
\item Dégage ce que ce document révèle du pouvoir du roi Njoya et du degré d'organisation du royaume bamoun.
\end{enumerate}
\end{exercice}

\courssub{Je me prépare au Bac}

\begin{exercice}[Question à réponse construite : le sultanat bamoun, un État organisé]
\textbf{Document 1.} \emph{« Le royaume bamoun, fondé par Nchare Yen, est dirigé depuis Foumban par un roi assisté de ministres et de conseillers. Le territoire est divisé en provinces confiées à des chefs qui lèvent les impôts et fournissent des hommes à l'armée. La cour compte de nombreux artisans : forgerons, sculpteurs, tisserands. »} (D'après un manuel d'histoire du Cameroun)

\textbf{Document 2.} \emph{« Sous le règne de Njoya (vers 1890-1933), le royaume connaît un essor remarquable : invention d'une écriture (le shümom), construction d'un nouveau palais, création d'écoles, rédaction d'une histoire et d'une coutume du royaume. »} (D'après les travaux sur le roi Njoya)

\textbf{Consignes :}
\begin{enumerate}
\item Présente chaque document en deux lignes (nature, auteur ou origine, sujet).
\item À partir du document 1, dégage deux éléments qui montrent l'existence d'institutions politiques et administratives dans le royaume bamoun.
\item À partir du document 2, montre que le règne de Njoya témoigne d'un État dynamique et rayonnant.
\item Rédige un paragraphe argumenté d'une quinzaine de lignes répondant à la question : \emph{montrez que le sultanat bamoun était un État organisé à la veille de la colonisation}.
\end{enumerate}
\end{exercice}

\begin{exercice}[Dissertation guidée : la diversité politique du Cameroun précolonial]
\textbf{Sujet :} \emph{« À la veille de la colonisation, le Cameroun était un espace politiquement organisé. » Discutez cette affirmation en vous appuyant sur l'étude du lamidat, de la chefferie des Grassfields et du sultanat bamoun.}

\textbf{Consignes :}
\begin{enumerate}
\item Définis les termes du sujet et délimite le cadre spatial et chronologique (annonce du plan en introduction).
\item Premier axe : présente le lamidat du Nord-Cameroun (chef, fondement religieux, organisation) avec un exemple précis.
\item Deuxième axe : présente la chefferie des Grassfields (fon, sociétés régulatrices, rôle des notables) avec un exemple précis.
\item Troisième axe : présente le sultanat bamoun (roi, administration, rayonnement culturel).
\item Rédige une conclusion qui répond à la question posée et souligne la richesse politique du Cameroun précolonial.
\end{enumerate}
\end{exercice}

\begin{exercice}[Travail dirigé : la chefferie de ma localité]
Dans le cadre du TD sur l'organisation d'une chefferie de la localité de ton établissement, tu réalises une enquête auprès des chefs traditionnels et des notables.

\begin{enumerate}
\item Élabore un guide d'entretien de cinq questions à poser au chef ou à un notable (identité du chef, mode de désignation, organes qui l'assistent, rôles de la chefferie, fêtes et rites).
\item À partir des réponses recueillies, réalise le croquis légendé de l'organisation de la chefferie : place le chef au sommet, puis les notables, les sociétés ou groupes d'âge, et la population ; indique par des flèches les relations entre ces éléments.
\item Rédige cinq lignes justifiant ta démarche d'enquête (choix des personnes interrogées, fiabilité des sources, recoupement des informations).
\item En t'appuyant sur ton croquis et sur le cours, indique deux ressemblances et une différence entre cette chefferie et le modèle général des chefferies des Grassfields.
\end{enumerate}
\end{exercice}

\newpage

\coursec{Corrigés des exercices}

\begin{corrige}[Exercice 1 — QCM : les entités politiques précoloniales]
\textbf{1. b) lamido} — Le chef d'un lamidat peul porte le titre de \cle{lamido} (terme dérivé de l'arabe \emph{amir} via le peul \emph{laamdo}), chef politique et religieux.

\textbf{2. c) Foumban} — La capitale du royaume bamoun est \cle{Foumban} (région de l'Ouest), siège du palais royal et de l'administration centrale depuis le XVIIIe siècle.

\textbf{3. a) des sociétés secrètes et des notables} — Dans les \cle{Grassfields}, le \cle{fon} exerce son pouvoir en collaboration avec des sociétés de régulation (ex. \emph{Kwifoyn}, \emph{Ngumba}) et un conseil de notables (\emph{tafoyn}, \emph{mafo}) qui limitent son autorité et assurent la continuité de l'État.

\textbf{4. b) Njoya} — Le roi \cle{Njoya} (règne ~1890–1933) est célèbre pour l'invention de l'écriture \cle{shümom}, la construction du palais de Foumban, la création d'écoles et la codification de l'histoire et des coutumes bamoun.
\end{corrige}

\begin{corrige}[Exercice 2 — Vrai ou faux, en justifiant]
\begin{enumerate}
\item \textbf{Faux.} À la veille de la colonisation (1884), le territoire camerounais était structuré en entités politiques organisées : lamidats au Nord, chefferies aux Grassfields, sultanat bamoun, chefferies sawa sur la côte, etc. Il n'y a pas de « vide politique ».
\item \textbf{Vrai.} Les lamidats du Nord-Cameroun (Rey-Bouba, Banyo, Tibati, Maroua, etc.) sont nés de l'expansion du \emph{jihad} mené par Ousman dan Fodio (début XIXe s.) et de ses lieutenants (Modibo Adama), qui ont implanté l'islam et l'administration sokotoise.
\item \textbf{Faux.} Le \cle{fon} (chef des Grassfields) n'est pas élu pour un mandat limité : il accède au trône par succession héréditaire (souvent matrilinéaire ou patrilinéaire selon les groupes) et règne à vie ; son pouvoir est sacré et encadré par les sociétés de régulation.
\item \textbf{Vrai.} Le sultanat bamoun disposait d'institutions centrales (roi, ministres/\emph{nji}, conseil, cour d'artisans), d'une administration provinciale (chefs de district levant l'impôt, recrutant l'armée) et d'une fiscalité organisée (tributs, corvées).
\end{enumerate}
\end{corrige}

\begin{corrige}[Exercice 3 — Définitions]
\begin{definition}[Lamidat]
Entité politique et religieuse du Nord-Cameroun, issue du jihad peul (début XIXe s.), dirigée par un \cle{lamido} (chef suprême, vassal du Sultan de Sokoto), organisant la société sur un mode féodal (aristocratie peule, populations soumises) et vivant d'élevage, d'agriculture et de commerce transsaharien. \textbf{Exemple :} le lamidat de \cle{Rey-Bouba} (fondé par Bouba Ndjida).
\end{definition}

\begin{definition}[Chefferie des Grassfields]
Entité politique ancienne des hautes terres de l'Ouest/Nord-Ouest, dirigée par un \cle{fon} (chef sacré, garant de la fécondité et de l'ordre), dont le pouvoir est partagé/contrôlé par des sociétés de régulation (Kwifoyn, Ngumba) et un conseil de notables ; organisation en chefferies de premier rang (ex. Bafut, Bamenda/Mankon) et de second rang. \textbf{Exemple :} la chefferie de \cle{Bafut} (Nord-Ouest).
\end{definition}

\begin{definition}[Sultanat bamoun]
Royaume ancien de la région de l'Ouest (fondation ~XIVe s. par Nchare Yen), islamisé au XIXe s., doté d'une administration centrale (ministères, cour, écriture shümom sous Njoya) et provinciale, d'une économie diversifiée (agriculture, artisanat, commerce) et d'un rayonnement culturel majeur. \textbf{Exemple :} le sultanat de \cle{Foumban} sous le roi Njoya.
\end{definition}
\end{corrige}

\begin{corrige}[Exercice 4 — Repères chronologiques]
\textbf{Frise chronologique (à tracer) :}

\begin{center}
\begin{tikzpicture}[scale=0.9]
% Axe principal
\draw[popdark, thick, ->] (0,0) -- (16,0) node[right, font=\small] {Années / Siècles};
% Graduations principales
\foreach \x/\label in {0/1400, 3/1600, 6/1800, 9/1884, 12/1900, 15/1933} {
    \draw (\x,0.1) -- (\x,-0.1) node[below, font=\scriptsize] {\label};
}
% Événements
% 1. Fondation royaume Bamoun ~1390 (Nchare Yen) -> approx 1400
\draw[poppurple, thick] (0.2,0.5) -- (0.2,1.5) node[above, font=\scriptsize, poppurple, align=center] {Fondation royaume\\Bamoun (Nchare Yen)\\~XIVe s.};
% 2. Jihad Ousman dan Fodio / lamidats -> 1804-1810 approx
\draw[poporange, thick] (5.8,0.5) -- (5.8,1.5) node[above, font=\scriptsize, poporange, align=center] {Jihad Ousman dan Fodio\\ \& lamidats Peuls\\début XIXe s.};
% 3. Arrivée Allemands 1884
\draw[popblue, thick] (8.8,0.5) -- (8.8,1.5) node[above, font=\scriptsize, popblue, align=center] {Arrivée Allemands\\côte camerounaise\\1884};
% 4. Règne Njoya ~1890-1933
\draw[popgreen, thick] (9.5,-0.5) -- (9.5,-1.5) node[below, font=\scriptsize, popgreen, align=center] {Règne du roi Njoya\\à Foumban\\~1890–1933};
% Siècles
\node[font=\small, popmuted] at (1.5,-0.6) {XVe s.};
\node[font=\small, popmuted] at (4.5,-0.6) {XVIIe–XVIIIe s.};
\node[font=\small, popmuted] at (7.5,-0.6) {XIXe s.};
\node[font=\small, popmuted] at (11,-0.6) {XXe s.};
\end{tikzpicture}
\end{center}

\textbf{Conclusion :} Cette frise montre que les organisations politiques camerounaises (royaume bamoun dès le XIVe s., chefferies Grassfields antérieures, lamidats au XIXe s.) sont \textbf{anciennes et bien enracinées} bien avant l'arrivée des Européens en 1884 ; le Cameroun précolonial n'est pas un espace vide mais un espace d'États structurés.
\end{corrige}

\begin{corrige}[Exercice 5 — Définition et distinction à partir d'un texte]
\begin{enumerate}
\item \textbf{Lamidat :} État du Nord-Cameroun né du jihad peul (début XIXe s.), dirigé par un \cle{lamido} (chef politique et religieux), vassal du Sultanat de Sokoto, à organisation féodale.\\
\textbf{Sultanat bamoun :} Royaume ancien de l'Ouest (fondé ~XIVe s. par Nchare Yen), islamisé au XIXe s., doté d'une administration centrale (ministres, cour) et provinciale, centré sur Foumban.

\item \textbf{Deux éléments de distinction relevés dans le texte :}
\begin{itemize}
\item \textit{Fondement du pouvoir :} le lamido gouverne « au nom de l'islam » et son autorité « repose sur le jihad » (légitimation religieuse/conquête) ; le roi bamoun s'appuie sur « un royaume ancien fondé par Nchare Yen » (légitimation dynastique/traditionnelle).
\item \textit{Origine de l'État :} le lamidat est une création issue de la conquête jihadiste (début XIXe s.) ; le sultanat bamoun est un « royaume ancien » (fondation ~XIVe s.), donc antérieur de plusieurs siècles.
\end{itemize}

\item \textbf{Trois phrases explicatives :}
Le lamidat tire sa légitimité de la conquête religieuse (jihad) et de la suzeraineté du Califat de Sokoto, ce qui en fait une entité d'implantation récente (début XIXe s.) à structure féodale islamisée. Le sultanat bamoun, au contraire, puise sa légitimité dans une dynastie fondatrice ancienne (Nchare Yen, ~XIVe s.) et dans la sacralité du roi, avec une administration laïque (ministres, provinces) qui s'est islamisée secondairement. Ainsi, le premier repose sur une autorité déléguée d'origine externe (Sokoto) et guerrière, le second sur une souveraineté interne, dynastique et administrative ancienne.
\end{enumerate}
\end{corrige}

\begin{corrige}[Exercice 6 — Localisation sur carte muette]
\textbf{Carte corrigée (description et code TikZ complet) :}

\begin{popfigure}
\begin{center}
\begin{tikzpicture}[scale=0.7]
% Contour très simplifié du Cameroun
\draw[popdark, thick] (0,0) -- (2,0) -- (3,1) -- (4,2.5) -- (3.5,4) -- (2,5) -- (0.5,4.5) -- (-0.5,3) -- (-0.5,1) -- cycle;
% Fleuve Bénoué
\draw[popblue, thick] (1.5,0) .. controls (1.8,1.5) and (2,2.5) .. (2.2,3.5);
\node[font=\tiny, popblue, rotate=30] at (2.5,1.5) {Bénoué};
% Fleuve Sanaga
\draw[popblue, thick] (1,2) .. controls (1.5,3) .. (1.8,4);
\node[font=\tiny, popblue, rotate=20] at (1.5,3.5) {Sanaga};
% Lac Tchad (coin)
\draw[popblue, fill=popblue!20] (-0.5,0) rectangle (0,0.5);
\node[font=\tiny, popblue] at (-0.2,0.2) {Lac Tchad};
% Villes / Points (toutes les villes du fond de carte)
\foreach \x/\y/\name in {
    0.5/0.8/Maroua, 1.5/1.5/Mora, 2.5/1.2/Kousseri,
    0.8/2.2/Garoua, 1.8/2.5/Ngaoundéré, 2.5/2.8/Meiganga,
    2.8/3.2/Banyo, 3.2/3.5/Rey-Bouba, 3.5/3.8/Tibati,
    1.2/3.8/Bafoussam, 0.8/4.2/Bamenda, 0.5/3.5/Bafut, 0.2/3.2/Bali,
    1.5/3.5/Foumban, 2.0/3.0/Bangangté, 2.5/2.0/Bertoua, 3.0/1.5/Batouri,
    1.0/1.0/Yagoua, 3.5/0.5/Kribi, 3.8/1.0/Ebolowa, 4.0/2.0/Yaoundé, 4.2/3.0/Douala} {
    \fill[poporange] (\x,\y) circle (2pt);
    \node[font=\tiny, right=1pt, popdark] at (\x,\y) {\name};
}
% --- ÉLÉMENTS DE L'EXERCICE ---
% 1. Trois lamidats du Nord : Rey-Bouba, Banyo, Maroua (marqués par un carré rouge)
\foreach \x/\y in {0.5/0.8, 2.8/3.2, 3.2/3.5} {
    \draw[poporange, thick] (\x-0.08,\y-0.08) rectangle (\x+0.08,\y+0.08);
}
% 2. Trois chefferies Grassfields : Bafut, Bamenda (Mankon), Bali (marqués par un triangle vert)
\foreach \x/\y in {0.5/3.5, 0.8/4.2, 0.2/3.2} {
    \draw[popgreen, thick] (\x,\y+0.1) -- (\x-0.08,\y-0.08) -- (\x+0.08,\y-0.08) -- cycle;
}
% 3. Capitale Bamoun : Foumban (étoile violette)
\node[star, star points=5, star point ratio=2.25, draw=poppurple, thick, fill=poppurple, minimum size=6pt] at (1.5,3.5) {};
\node[font=\tiny, poppurple, above=2pt] at (1.5,3.5) {Foumban*};
% 4. Limites approximatives des trois zones (polygones semi-transparents)
% Nord soudanien (hachures orange) - zone au nord de la ligne Garoua-Ngaoundéré-Meiganga
\fill[poporange!15, opacity=0.5] (0,0) -- (2,0) -- (3,1) -- (3.5,2) -- (2.5,2.5) -- (1.5,2.5) -- (0.5,2) -- (-0.5,1) -- cycle;
\pattern[pattern=north east lines, pattern color=poporange!50] (0,0) -- (2,0) -- (3,1) -- (3.5,2) -- (2.5,2.5) -- (1.5,2.5) -- (0.5,2) -- (-0.5,1) -- cycle;
% Grassfields (hachures vert) - zone montagneuse Ouest/Nord-Ouest
\fill[popgreen!15, opacity=0.5] (0.5,2) -- (1.5,2.5) -- (2.5,1.5) -- (3,2.5) -- (2.5,4) -- (0.5,4.5) -- (-0.5,3) -- cycle;
\pattern[pattern=north west lines, pattern color=popgreen!50] (0.5,2) -- (1.5,2.5) -- (2.5,1.5) -- (3,2.5) -- (2.5,4) -- (0.5,4.5) -- (-0.5,3) -- cycle;
% Royaume Bamoun (points violet) - autour de Foumban
\fill[poppurple!15, opacity=0.5] (1.2,3.2) circle (0.7);
\pattern[pattern=dots, pattern color=poppurple!60] (1.2,3.2) circle (0.7);
% Légende interne (dans la figure)
\node[font=\scriptsize, popdark, anchor=south west] at (-0.5,5.0) {\textbf{Légende :}};
\node[font=\scriptsize, popdark, anchor=south west] at (-0.5,5.4) {\textcolor{poporange}{\rule{3mm}{1.5mm}} Nord soudanien (lamidats)};
\node[font=\scriptsize, popdark, anchor=south west] at (-0.5,5.8) {\textcolor{popgreen}{\rule{3mm}{1.5mm}} Grassfields (chefferies)};
\node[font=\scriptsize, popdark, anchor=south west] at (-0.5,6.2) {\textcolor{poppurple}{\rule{3mm}{1.5mm}} Royaume Bamoun};
\node[font=\scriptsize, popdark, anchor=south west] at (-0.5,6.6) {\textcolor{poporange}{\framebox(4,4){}} Lamidats (Rey-Bouba, Banyo, Maroua)};
\node[font=\scriptsize, popdark, anchor=south west] at (-0.5,7.0) {\textcolor{popgreen}{\raisebox{1pt}{\tikz{\draw[popgreen,thick] (0,0) -- (0.1,0.15) -- (-0.1,0.15) -- cycle;}}} Chefferies (Bafut, Bamenda, Bali)};
\node[font=\scriptsize, popdark, anchor=south west] at (-0.5,7.4) {\textcolor{poppurple}{$\bigstar$} Foumban (capitale Bamoun)};
% Nord
\draw[->, thick] (-0.8,0.2) -- (-0.8,1.2) node[above, rotate=90, font=\tiny] {N};
% Titre de la carte
\node[font=\normalsize\bfseries, popdark, anchor=north] at (2,7.8) {Cameroun à la veille de la colonisation : entités politiques (vers 1884)};
\end{tikzpicture}
\end{center}
\legende{Carte corrigée : localisation des trois lamidats (carrés orange), trois chefferies Grassfields (triangles verts), capitale Foumban (étoile violette) et zones culturelles (hachures/points).}
\end{popfigure}

\textbf{Éléments attendus dans la légende et le titre :}
\begin{itemize}
\item Titre : « Cameroun à la veille de la colonisation : entités politiques (vers 1884) » ou équivalent.
\item Légende complète : symboles pour lamidats, chefferies Grassfields, capitale Bamoun ; zones culturelles (Nord soudanien, Grassfields, Royaume Bamoun) ; orientation (flèche Nord) ; échelle indicative (non tracée ici mais à mentionner : « Échelle approximative »).
\end{itemize}
\end{corrige}

\begin{corrige}[Exercice 7 — Tableau comparatif à compléter]
\begin{center}
\begin{tabularx}{\linewidth}{|l|X|X|X|}
\hline
\rowcolor{popblueL} \textbf{Critère} & \textbf{Lamidat} & \textbf{Chefferie des Grassfields} & \textbf{Sultanat bamoun} \\
\hline
Titre du chef & \cle{Lamido} (chef politique et religieux) & \cle{Fon} (chef sacré, garant de la fécondité) & \cle{Sultan} / Roi (Njoya), titre de \emph{Mfon} \\
\hline
Fondement du pouvoir & Jihad (conquête religieuse), légitimité islamique, vassalité à Sokoto & Tradition dynastique, sacralité du fon, contre-pouvoir des sociétés secrètes (Kwifoyn) et notables & Dynastie ancienne (Nchare Yen), sacralité royale, islamisation secondaire (XIXe), administration centrale \\
\hline
Institutions principales & Lamido, conseil de notables (titre \emph{galadima}, \emph{sarkin}), cadis (juges islamiques), chefs de districts (\emph{ardos}) & Fon, sociétés de régulation (Kwifoyn, Ngumba), conseil de notables (tafoyn, mafo), chefs de quartiers/vassaux & Sultan, ministres (\emph{nji} : \emph{Nji Mfor}, \emph{Nji Nchinda}, etc.), cour d'artisans, chefs de provinces (\emph{nduna}), armée, écriture (shümom) \\
\hline
Activités économiques dominantes & Élevage (zébu), agriculture (mil, sorgho, coton), commerce transsaharien (sel, bétail, esclaves), artisanat (cuir, tissage) & Agriculture intensive (maïs, igname, café, palmier à huile), élevage (petit bétail), artisanat (sculpture, poterie, tissage), commerce local et régional & Agriculture (maïs, igname, huile de palme), artisanat de cour très développé (fonte du bronze, sculpture, tissage, perles), commerce (kola, sel, artisanat), tributs provinciaux \\
\hline
\end{tabularx}
\end{center}

\textbf{Conclusion (trois phrases) :}
À la veille de la colonisation, le Cameroun présente une grande diversité politique : au Nord, des lamidats féodaux islamisés nés du jihad ; aux Grassfields, des chefferies anciennes à pouvoir partagé entre fon sacré et sociétés de régulation ; à l'Ouest, un sultanat bamoun centralisé, doté d'une administration bureaucratique et d'un rayonnement culturel original (écriture, arts). Ces trois modèles coexistent sans s'unifier, témoignant de l'adaptation des sociétés à leurs milieux (sahélien, montagnard, forestier) et à leurs histoires propres. Cette richesse politique sera un facteur majeur dans la résistance et la négociation face à la pénétration coloniale.
\end{corrige}

\begin{corrige}[Exercice 8 — Étude de photographie : le palais de Foumban]
\begin{enumerate}
\item \textbf{Description de l'architecture :} Le palais de Foumban (photographie de 1917) est un vaste bâtiment à plan rectangulaire ou carré, à un ou deux étages, dont la façade principale est rythmée par une série de piliers sculptés supportant un toit à larges débords. Les murs sont en terre crue (banco) ou en briques de terre, enduits. La façade présente des motifs géométriques et figuratifs (représentant le roi, des animaux, des symboles de pouvoir) sculptés en bas-relief sur les piliers et les panneaux. Une grande cour ouverte (place) devant le palais accueille les cérémonies publiques. L'ensemble dénote une architecture monumentale, soignée, alliant tradition locale (banco, toiture) et innovations (plan symétrique, décor sculpté systématique).

\item \textbf{Deux indices de double fonction (résidence + centre admin/éco) :}
\begin{itemize}
\item La présence d'\emph{ateliers d'artisans} (forgerons, sculpteurs, tisserands) \emph{à l'intérieur même du palais} : cela montre que la production artisanale de prestige (objets de cour, regalia, tissus ndop) est une activité économique contrôlée directement par le roi.
\item Le fait que le palais \emph{abrite la cour, les services du royaume (ministères, archives, écoles) et que la place sert aux grandes cérémonies} : le palais est le siège physique du gouvernement (administration, justice, diplomatie) et le lieu de manifestation du pouvoir politique (cérémonies, réception des tributs).
\end{itemize}

\item \textbf{Ce que le document révèle du pouvoir de Njoya et de l'organisation du royaume :}
Le palais monumental, centre névralgique réunissant résidence, administration, justice, production artisanale et vie cérémonielle, révèle que le roi Njoya exerce un \textbf{pouvoir centralisé, sacré et bureaucratique}. L'existence d'ateliers intégrés prouve une économie de cour planifiée et un contrôle royal sur les ressources symboliques et matérielles. L'architecture même (plan régulier, décor codifié) traduit une pensée étatique structurée, capable de mobiliser la main-d'œuvre et les savoir-faire pour un projet dynastique. Le royaume bamoun apparaît ainsi comme un \textbf{État organisé, rayonnant, à la fois traditionnel et innovant} (écriture, écoles, musées), capable de s'adapter (islamisation, contacts allemands) sans perdre sa cohérence interne.
\end{enumerate}
\end{corrige}

\begin{corrige}[Exercice 9 — Question à réponse construite : le sultanat bamoun, un État organisé]
\textbf{Barème indicatif : 10 points}
\begin{itemize}
\item Présentation des documents (2 pts : 1 pt par doc — nature, origine, sujet)
\item Exploitation Doc 1 (2 pts : 2 éléments distincts + justes)
\item Exploitation Doc 2 (2 pts : dynamisme + rayonnement)
\item Rédaction paragraphe argumenté (4 pts : thèse, arguments structurés, vocabulaire précis, conclusion)
\end{itemize}

\textbf{Points de vigilance :}
\begin{itemize}
\item Citer explicitement les documents (« Document 1 », « Document 2 »).
\item Ne pas confondre « nature » (extrait de manuel, extrait d'ouvrage historique) et « auteur ».
\item Pour le paragraphe : utiliser les deux documents, articuler institutions + dynamisme culturel, conclure sur le caractère « organisé » avant 1884 (même si Njoya règne après, les structures existaient avant).
\end{itemize}

\medskip
\textbf{Correction détaillée :}

\textbf{1. Présentation des documents :}
\begin{itemize}
\item \textbf{Document 1 :} Extrait d'un manuel d'histoire du Cameroun (nature : source secondaire, synthèse pédagogique), auteur non nommé (collectif d'historiens/enseignants), sujet : organisation politique et administrative du royaume bamoun (institutions centrales, provinciales, cour artisanale).
\item \textbf{Document 2 :} Extrait d'un ouvrage historique / travaux universitaires sur le roi Njoya (nature : source secondaire, étude spécialisée), origine : travaux sur Njoya (ex. travaux de Nji Oumarou Nchare, ou synthèses récentes), sujet : réalisations intellectuelles, architecturales et éducatives sous le règne de Njoya (1890–1933).
\end{itemize}

\textbf{2. Deux éléments d'institutions politiques/administratives (Doc 1) :}
\begin{itemize}
\item Existence d'un \textbf{gouvernement central} : le roi est « assisté de ministres et de conseillers » (ex. Nji Mfor, Nji Nchinda), ce qui marque une séparation des fonctions (exécutif, justice, finances, armée).
\item \textbf{Administration territoriale structurée} : le territoire est « divisé en provinces confiées à des chefs » qui perçoivent les impôts et lèvent des hommes pour l'armée, preuve d'une fiscalité et d'une conscription organisées à l'échelle de l'État.
\end{itemize}

\textbf{3. Dynamisme et rayonnement sous Njoya (Doc 2) :}
\begin{itemize}
\item \textbf{Innovation intellectuelle majeure :} invention de l'écriture \cle{shümom} (syllabaire original), permettant la rédaction de l'histoire, de la coutume, de pharmacopées, et l'enseignement dans des écoles créées à cet effet.
\item \textbf{Réalisations matérielles et symboliques :} construction d'un \textbf{nouveau palais} monumental (siège du pouvoir), création d'\textbf{écoles} (formation des cadres), codification de l'\textbf{histoire et de la coutume} (constitution écrite avant l'heure), développement des arts de cour (bronze, sculpture, tissage). Cela montre un État qui se dote d'outils de mémoire, de formation et de légitimation modernes tout en s'appuyant sur sa tradition.
\end{itemize}

\textbf{4. Paragraphe argumenté (environ 15 lignes) :}

Le sultanat bamoun était, à la veille de la colonisation, un État pleinement organisé, doté d'institutions politiques stables et d'une capacité d'innovation remarquable. Le document 1 montre que le royaume repose sur une administration centrale structurée autour du roi (Mfon) secondé par des ministres spécialisés (Nji), et sur un maillage territorial efficace : le pays est divisé en provinces gérées par des chefs délégués qui assurent la perception de l'impôt et le recrutement militaire, garantissant ainsi la ressource fiscale et la défense. La cour elle-même intègre des corps de métier (forgerons, sculpteurs, tisserands), ce qui place l'économie de prestige sous contrôle direct du pouvoir. Le document 2 complète ce tableau en révélant le dynamisme du règne de Njoya (1890–1933) qui, s'inscrivant dans la continuité dynastique, dote l'État d'outils modernes : une écriture originale (le shümom), un système éducatif, une historiographie officielle et une architecture monumentale. Cette double capacité — maintenir des institutions traditionnelles performantes (administration, armée, fiscalité) et impulser des innovations culturelles majeures (écriture, écoles, arts) — prouve que le sultanat bamoun n'était pas une chefferie rudimentaire mais un \textbf{État centralisé, bureaucratisé et rayonnant}, capable de s'adapter aux défis de son temps (islamisation, pression coloniale) tout en affirmant sa souveraineté interne.
\end{corrige}

\begin{corrige}[Exercice 10 — Dissertation guidée : la diversité politique du Cameroun précolonial]
\textbf{Barème indicatif : 20 points}
\begin{itemize}
\item Introduction (4 pts) : définition termes, délimitation spatio-temporelle, problématique, annonce plan.
\item Axe 1 Lamidat (4 pts) : chef, fondement, organisation, exemple précis.
\item Axe 2 Chefferie Grassfields (4 pts) : fon, sociétés, notables, exemple précis.
\item Axe 3 Sultanat Bamoun (4 pts) : roi, administration, rayonnement culturel.
\item Conclusion (4 pts) : réponse à la question, richesse politique, ouverture.
\end{itemize}

\textbf{Points de vigilance :}
\begin{itemize}
\item Respecter la structure dissertation (intro, développement en 3 parties équilibrées, conclusion).
\item Ne pas mélanger les modèles dans un même paragraphe.
\item Citer des exemples précis (Rey-Bouba, Bafut, Foumban/Njoya).
\item Utiliser le vocabulaire du cours : lamido, fon, sultan, jihad, Kwifoyn, Nji, shümom, vassalité, tribut, société de régulation.
\item Chronologie : « à la veille de la colonisation » = vers 1884 (mais structures antérieures).
\end{itemize}

\medskip
\textbf{Proposition de dissertation rédigée :}

\textbf{Introduction}
[Accroche] L'historiographie coloniale a longtemps présenté l'Afrique comme un continent « sans histoire » ni États avant l'arrivée des Européens. [Définitions] Le « Cameroun précolonial » désigne l'espace géographique de l'actuel Cameroun avant le traité de protectorat germano-douala de juillet 1884. « Politiquement organisé » signifie l'existence d'entités dotées d'institutions stables (chef, administration, justice, fiscalité, armée) assurant l'ordre interne et les relations extérieures. [Délimitation] Cadre spatial : Nord (bénoué), Grassfields (Ouest/Nord-Ouest), région de Foumban (Ouest). Cadre chronologique : des origines (XIVe s. pour Bamoun) à 1884. [Problématique] Comment la diversité des modèles politiques (lamidats, chefferies Grassfields, sultanat bamoun) prouve-t-elle que le Cameroun n'était pas un vide politique mais un espace d'États structurés et originaux ? [Annonce du plan] Nous analyserons tour à tour le lamidat du Nord (I), la chefferie des Grassfields (II) et le sultanat bamoun (III) pour montrer la richesse de l'organisation politique camerounaise.

\textbf{Premier axe : Le lamidat du Nord-Cameroun — une féodalité islamisée issue du jihad}
Né de l'expansion du jihad d'Ousman dan Fodio (1804) mené par Modibo Adama, le lamidat est une entité politique et religieuse. Le \cle{lamido} (ex. Bouba Ndjida à Rey-Bouba, ou le lamido de Banyo) en est le chef suprême : il détient l'autorité politique (nomination des \emph{ardos} chefs de districts, perception de la \emph{jangali} — impôt sur le bétail) et l'autorité religieuse (imamat, application de la charia via des cadis). Vassal du Sultan de Sokoto (califat), il lui verse un tribut annuel. La société est stratifiée : aristocratie peule (conquérants), populations autochtones intégrées (Mboum, Koutine) ou soumises (Kirdi, « païens » raflés pour l'esclavage). L'économie repose sur l'élevage bovin, l'agriculture (mil, sorgho, coton) et le commerce transsaharien. Exemple précis : le lamidat de \cle{Rey-Bouba}, fondé vers 1805 par Bouba Ndjida, couvre une vaste zone entre Bénoué et Mayo-Rey, structuré en districts (\emph{lamidats} de second rang) et prospère grâce au commerce du bétail et de l'ivoire.

\textbf{Deuxième axe : La chefferie des Grassfields — un pouvoir partagé entre sacré et contre-pouvoirs}
Dans les hautes terres de l'Ouest et du Nord-Ouest, les chefferies (fon) sont très antérieures au XIXe siècle (XVe–XVIIe s.). Le \cle{fon} (ex. le fon de Bafut, ou le fon de Mankon/Bamenda) est un chef sacré, garant de la fécondité de la terre et de la cohésion sociale ; il est le « père » de son peuple. Mais son pouvoir n'est pas absolu : il est encadré, voire contrôlé, par des \textbf{sociétés de régulation} (la \cle{Kwifoyn} — société masculine de nuit, police religieuse et judiciaire ; le \cle{Ngumba} — société d'élite) et par un \textbf{conseil de notables} (\emph{tafoyn}, \emph{mafo}) qui participent à la prise de décision, à la justice et à la succession. L'organisation territoriale distingue chefferies de premier rang (capitales régionales, ex. Bafut, Bamenda, Bali, Kumbo) et chefferies de second rang (vassales). L'économie intensive (cultures de terroirs, café, palmier à huile) et l'artisanat (sculpture, poterie, tissage ndop) soutiennent la puissance du fon. Exemple précis : la chefferie de \cle{Bafut} (Nord-Ouest), où le fon règne avec le Kwifoyn et un conseil de neuf notables, et où la fête annuelle de l'\emph{Abin} (nouvel an) renouvelle le pacte social.

\textbf{Troisième axe : Le sultanat bamoun — un État centralisé, bureaucratique et rayonnant}
Au cœur de la région de l'Ouest, le royaume bamoun, fondé au XIVe siècle par \cle{Nchare Yen} (dynastie Nchare), a sa capitale à \cle{Foumban}. Islamisé progressivement au XIXe siècle (sous le sultan Njoya vers 1890), il conserve une structure étatique originale. Le \cle{sultan} (Mfon) est chef de l'État, de l'armée et de la religion ; il s'appuie sur une \textbf{administration centrale} composée de ministres (\emph{Nji}) aux portefeuilles définis (Nji Mfor — premier ministre/intérieur, Nji Nchinda — justice, Nji Poutgnigni — finances, Nji Gbetnkom — armée, etc.). Le territoire est divisé en \textbf{provinces} dirigées par des chefs (\emph{nduna}) nommés, qui lèvent l'impôt (produits agricoles, artisanat) et fournissent des contingents. La \textbf{cour d'artisans} (forgerons, sculpteurs, tisserands, perliers) produit les regalia et objets de prestige, gérés par le palais. Sous Njoya, le royaume innove : invention de l'écriture \cle{shümom} (syllabaire), création d'écoles, rédaction de l'histoire et de la coutume (code juridique), construction du palais monumental. Le sultanat bamoun est ainsi un État « moderne avant l'heure », alliant tradition dynastique et rationalité administrative.

\textbf{Conclusion}
[Bilan] Les trois modèles étudiés — lamidat féodal islamisé, chefferie à pouvoir partagé des Grassfields, sultanat bureaucratique bamoun — démontrent sans ambiguïté que le Cameroun à la veille de la colonisation (1884) était un \textbf{espace politiquement organisé}, dense et diversifié. [Réponse] Chaque région a produit des institutions adaptées à son histoire (jihad, migration, dynastie), à son milieu (sahélien, montagnard, forestier) et à ses ressources. [Ouverture] Cette richesse politique explique la variété des réactions face à la colonisation (résistance armée au Nord, diplomatie et adaptation à Foumban, négociations aux Grassfields) et constitue le socle sur lequel s'est construit l'État camerounais contemporain, qui reconnaît aujourd'hui encore la légitimité des chefferies traditionnelles.
\end{corrige}

\begin{corrige}[Exercice 11 — Travail dirigé : la chefferie de ma localité]
\textit{Note : cet exercice étant un TD de terrain, la correction ci-dessous donne les \textbf{attendus types} et une \textbf{méthodologie} ; les réponses concrètes dépendent de la localité de l'établissement (ex. Bafoussam, Dschang, Bamenda, Maroua, etc.). L'élève doit adapter le guide et le croquis à sa réalité locale.}

\begin{enumerate}
\item \textbf{Guide d'entretien (5 questions types) :}
\begin{enumerate}
\item \textit{Identité et accession :} « Quel est votre titre officiel ? Comment êtes-vous devenu chef (héritage, élection, désignation par le conseil, investiture) ? Depuis quand occupez-vous cette fonction ? »
\item \textit{Organes d'assistance et de contrôle :} « Quels sont les groupes ou personnes qui vous aident à gouverner (notables, sociétés secrètes, conseils de quartier, groupes d'âge, femmes titulaires) ? Quel est leur rôle précis (conseil, exécution, contrôle, justice) ? »
\item \textit{Fonctions et pouvoirs :} « Quelles sont vos attributions principales (justice coutumière, gestion foncière, fiscalité locale, organisation des fêtes, représentation auprès de l'administration moderne) ? Comment s'articulent vos pouvoirs avec ceux de l'administration déléguée (sous-préfet, maire) ? »
\item \textit{Ressources et économie de la chefferie :} « De quels moyens dispose la chefferie (terres, redevances, contributions volontaires, subventions, revenus d'activités économiques) pour fonctionner et réaliser des projets collectifs ? »
\item \textit{Rites, fêtes et symboles de l'unité :} « Quelles sont les grandes cérémonies annuelles ou exceptionnelles (fête des récoltes, intronisation, culte des ancêtres, danse des masques) qui marquent la vie de la chefferie ? Quels en sont les symboles (trône, objets sacrés, masques, lieux sacrés) ? »
\end{enumerate}

\item \textbf{Croquis légendé de l'organisation (description type) :}
\begin{center}
\begin{tikzpicture}[scale=0.8, every node/.style={font=\small}]
% Chef au sommet
\node[draw, ellipse, fill=poporange!20, thick, minimum width=3cm, minimum height=1.2cm, align=center] (chef) at (0,4){\textbf{CHEF (Fon / Lamido / Sultan)}\\(Garant de l'unité, justice, rites)};
% Notables / Conseil
\node[draw, rectangle, fill=popblue!20, thick, minimum width=3.5cm, minimum height=1cm, align=center] (notables) at (-4,2){NOTABLES / CONSEIL\\(Conseillers, porte-parole,\\gestion quotidienne)};
\node[draw, rectangle, fill=popblue!20, thick, minimum width=3.5cm, minimum height=1cm, align=center] (notables2) at (4,2){SOCIÉTÉS DE RÉGULATION /\\GROUPES D'ÂGE\\(Kwifoyn, Ngumba, Kwifon...\\Contrôle, police, initiation)};
% Population / Quartiers
\node[draw, rectangle, rounded corners, fill=popgreen!20, thick, minimum width=4cm, minimum height=1cm, align=center] (pop1) at (-5,0){QUARTIERS / VILLAGES\\(Chefs de quartier, chefs de famille)};
\node[draw, rectangle, rounded corners, fill=popgreen!20, thick, minimum width=4cm, minimum height=1cm, align=center] (pop2) at (0,0){POPULATION\\(Contribuables, cultivateurs,\\artisans, jeunes)};
\node[draw, rectangle, rounded corners, fill=popgreen!20, thick, minimum width=4cm, minimum height=1cm, align=center] (pop3) at (5,0){ASSOCIATIONS\\Féminines, jeunes, tontines,\\Églises/Mosquées};
% Flèches
\draw[->, thick, popdark] (chef.south west) -- (notables.north east) node[midway, above left, font=\tiny] {Consulte, délègue};
\draw[->, thick, popdark] (chef.south east) -- (notables2.north west) node[midway, above right, font=\tiny] {Contrôle, collabore};
\draw[->, thick, popdark] (notables.south) -- (pop1.north) node[midway, left, font=\tiny] {Relaye ordres};
\draw[->, thick, popdark] (notables2.south) -- (pop2.north) node[midway, above, font=\tiny] {Sanctionne, encadre};
\draw[<->, thick, popmuted] (pop1.east) -- (pop2.west) node[midway, below, font=\tiny] {Solidarité, travail};
\draw[<->, thick, popmuted] (pop2.east) -- (pop3.west) node[midway, below, font=\tiny] {Participation};
\draw[->, dashed, popdark] (pop1.north west) .. controls (-6,3) and (-4,3.5) .. (chef.west) node[midway, left, font=\tiny] {Tributs, doléances, info};
\draw[->, dashed, popdark] (pop3.north east) .. controls (6,3) and (4,3.5) .. (chef.east) node[midway, right, font=\tiny] {Soutien, alertes};
\end{tikzpicture}
\end{center}
\textbf{Légende :} Flèches pleines = autorité descendante / délégation ; Flèches doubles = relations horizontales ; Flèches pointillées = remontée d'informations, tributs, légitimité. Adapter les noms des boîtes à la réalité locale (ex. « Lamido » au Nord, « Fon » à l'Ouest, « Chef de 3e degré » au Sud).

\item \textbf{Justification de la démarche d'enquête (5 lignes type) :}
« J'ai interrogé le chef traditionnel (ou son représentant désigné) et deux notables reconnus (un homme, une femme si possible) pour croiser les points de vue sur l'organisation. Le choix de ces personnes repose sur leur légitimité coutumière et leur connaissance directe des institutions. La fiabilité est assurée par le recoupement de leurs réponses entre elles et avec les données du cours (modèle Grassfields / lamidat). J'ai privilégié l'entretien semi-directif pour laisser s'exprimer les conceptions locales tout en guidant sur les thèmes du programme. Les informations ont été notées sur place et validées par une relecture collective en classe. »

\item \textbf{Comparaison avec le modèle général des Grassfields (deux ressemblances, une différence) :}
\begin{itemize}
\item \textbf{Ressemblance 1 :} Présence d'un chef sacré (Fon) au sommet, garant de la tradition et de la fécondité, assisté d'un conseil de notables (tafoyn/mafo) qui participe à la gouvernance.
\item \textbf{Ressemblance 2 :} Existence de sociétés de régulation (Kwifoyn, Ngumba ou équivalents locaux) qui jouent un rôle de police coutumière, de justice et de contre-pouvoir vis-à-vis du Fon.
\item \textbf{Différence :} (Variable selon la localité) \textit{Exemple :} Dans ma chefferie (ex. Bafoussam, chefferie de 1er rang), le Fon dispose d'une administration plus bureaucratique (secrétaires, services techniques) et d'une emprise foncière plus forte qu'une chefferie de 2e rang du modèle classique ; \textit{ou} : ma chefferie (ex. zone Sawa ou Nord) n'a pas de sociétés secrètes de type Kwifoyn mais des groupes d'âge ou des conseils de famille comme organes de régulation.
\end{itemize}
\end{enumerate}
\end{corrige}

\newpage

\coursec{Fiche bilan}

\begin{popfigure}
\begin{center}\tcbox[colback=popblue,colframe=popblue,coltext=white,arc=9pt,boxsep=4pt,left=6pt,right=6pt]{\bfseries\small Le Cameroun à la veille de la colonisation}\end{center}
\vspace{-2pt}
\begin{tcbraster}[raster columns=3,raster equal height=rows,raster column skip=6pt,raster row skip=6pt,enhanced,blankest]
\begin{tcolorbox}[enhanced,colback=poporange,colframe=poporange,coltext=white,boxrule=0.9pt,arc=3pt,left=4pt,right=4pt,top=3pt,bottom=3pt,boxsep=1pt,halign=left]\bfseries\scriptsize Le lamidat peul (Nord) : État islamique, Lamido, cavaleries\end{tcolorbox}
\begin{tcolorbox}[enhanced,colback=popgreen,colframe=popgreen,coltext=white,boxrule=0.9pt,arc=3pt,left=4pt,right=4pt,top=3pt,bottom=3pt,boxsep=1pt,halign=left]\bfseries\scriptsize La chefferie des Grassfields : fon, sociétés secrètes, cases\end{tcolorbox}
\begin{tcolorbox}[enhanced,colback=poppurple,colframe=poppurple,coltext=white,boxrule=0.9pt,arc=3pt,left=4pt,right=4pt,top=3pt,bottom=3pt,boxsep=1pt,halign=left]\bfseries\scriptsize Le sultanat bamoun : sultan, cour, royaume centralisé\end{tcolorbox}
\begin{tcolorbox}[enhanced,colback=popgold,colframe=popgold,coltext=white,boxrule=0.9pt,arc=3pt,left=4pt,right=4pt,top=3pt,bottom=3pt,boxsep=1pt,halign=left]\bfseries\scriptsize Traits communs : hiérarchie, justice, armée, impôt\end{tcolorbox}
\begin{tcolorbox}[enhanced,colback=poppink,colframe=poppink,coltext=white,boxrule=0.9pt,arc=3pt,left=4pt,right=4pt,top=3pt,bottom=3pt,boxsep=1pt,halign=left]\bfseries\scriptsize Diversité des modèles politiques\end{tcolorbox}
\begin{tcolorbox}[enhanced,colback=popblueL,colframe=popblueL,coltext=white,boxrule=0.9pt,arc=3pt,left=4pt,right=4pt,top=3pt,bottom=3pt,boxsep=1pt,halign=left]\bfseries\scriptsize Sources : traditions orales, récits, vestiges\end{tcolorbox}
\begin{tcolorbox}[enhanced,colback=poporangeL,colframe=poporangeL,coltext=white,boxrule=0.9pt,arc=3pt,left=4pt,right=4pt,top=3pt,bottom=3pt,boxsep=1pt,halign=left]\bfseries\scriptsize TD : chefferie de ma localité\end{tcolorbox}
\end{tcbraster}

\legende{Carte mentale de la leçon : les entités politiques du Cameroun précolonial.}
\end{popfigure}

\begin{bilan}
\textbf{Définitions clés}
\begin{itemize}
  \item \cle{Lamidat} : État politique musulman dirigé par un \cle{Lamido}, issu des conquêtes peules du XIX\up{e} siècle au Nord-Cameroun (ex. : Rey-Bouba, Banyo, Ngaoundéré).
  \item \cle{Chefferie des Grassfields} : entité politique de l'Ouest et du Nord-Ouest dirigée par un chef sacré (\emph{fon} ou \emph{fo}), entouré de notables et de sociétés régulatrices.
  \item \cle{Sultanat bamoun} : royaume centralisé de Foumban, dirigé par un sultan, fondé par Nchare Yen, avec une cour organisée et une armée.
  \item \cle{Société secrète (ou régulatrice)} : association (ex. : \emph{Tukah}, \emph{Mfam}) qui aide le chef à maintenir l'ordre et à exécuter les décisions.
  \item \cle{Tradition orale} : source historique transmise de bouche à oreille (généalogies, récits, éloges) ; elle doit être croisée avec d'autres sources.
\end{itemize}

\textbf{Repères à connaître par cœur}
\begin{center}
\begin{tabularx}{\linewidth}{|l|X|X|}\hline
\rowcolor{popblueL} \textbf{Entité} & \textbf{Chef} & \textbf{Caractéristique majeure} \\ \hline
Lamidat (Nord) & Lamido & État islamique, impôt (\emph{zakat}), cavalerie \\ \hline
Chefferie (Grassfields) & Fon / Fo & Chef sacré, sociétés secrètes, organisation en quartiers \\ \hline
Sultanat bamoun & Sultan (Mfon) & Royaume centralisé, cour, écriture \emph{shümom} sous Njoya \\ \hline
\end{tabularx}
\end{center}

\textbf{Méthodes express}
\begin{itemize}
  \item \textbf{Commentaire de document} : présenter (nature, auteur, date) $\rightarrow$ dégager les idées $\rightarrow$ montrer ce que le document révèle de l'organisation politique.
  \item \textbf{Comparer deux entités} : critères fixes = chef, institutions, justice, armée, impôt, religion.
  \item \textbf{Enquête de TD} : interroger les anciens et notables, observer le palais et les cases, noter les titres des dignitaires, croiser les informations.
\end{itemize}
\end{bilan}

\begin{aretenir}[Checklist avant le Bac]
\begin{itemize}
  \item $\square$ Je sais définir : lamidat, chefferie, sultanat, société secrète.
  \item $\square$ Je sais situer géographiquement les trois entités (Nord, Grassfields, Foumban).
  \item $\square$ Je connais le titre du chef dans chaque entité (Lamido, fon, sultan).
  \item $\square$ Je peux citer au moins deux lamidats du Nord-Cameroun.
  \item $\square$ Je sais que Nchare Yen est le fondateur du royaume bamoun et Njoya un sultan célèbre.
  \item $\square$ Je connais le rôle des sociétés secrètes dans les Grassfields.
  \item $\square$ Je sais comparer les trois modèles selon des critères précis.
  \item $\square$ Je sais expliquer la valeur et les limites de la tradition orale.
  \item $\square$ Je peux décrire l'organisation d'une chefferie de ma localité (chef, notables, quartiers).
  \item $\square$ Je sais rédiger une introduction et une conclusion en commentaire de document.
\end{itemize}
\end{aretenir}

\findecours

\end{document}
